\documentclass{article}

\usepackage{neurips_2026}

\usepackage[utf8]{inputenc} % allow utf-8 input
\usepackage[T1]{fontenc}    % use 8-bit T1 fonts
\usepackage{hyperref}       % hyperlinks
\usepackage{url}            % simple URL typesetting
\usepackage{booktabs}       % professional-quality tables
\usepackage{amsfonts}       % blackboard math symbols
\usepackage{nicefrac}       % compact symbols for 1/2, etc.
\usepackage{microtype}      % microtypography
\usepackage{xcolor}         % colors

\usepackage{multirow}
\usepackage{amssymb}
\usepackage{graphicx}
\usepackage{natbib}
\usepackage{doi}
\usepackage{amsmath}
\usepackage{physics}
\usepackage{wrapfig}
\usepackage{bm}
\usepackage{algorithm,algpseudocode}
\usepackage{forest}
\usepackage{tcolorbox}
\usepackage{tikz-cd}
\usetikzlibrary{calc}
\usepackage{enumitem}

\usepackage{amsthm}

\newtheorem{theorem}{Theorem}
\newtheorem{lemma}{Lemma}
\newtheorem{proposition}{Proposition}
\newtheorem{corollary}{Corollary}

\theoremstyle{definition}
\newtheorem{definition}{Definition}
\newtheorem{assumption}{Assumption}

\theoremstyle{remark}
\newtheorem{remark}{Remark}
\newtheorem{example}{Example}

% Blackboard bold sets
\newcommand{\R}{\mathbb{R}}
\newcommand{\N}{\mathbb{N}}
\newcommand{\Z}{\mathbb{Z}}
\newcommand{\Q}{\mathbb{Q}}
\newcommand{\C}{\mathbb{C}}

% Labels
\newcommand{\alg}{\mathrm{alg}}
\newcommand{\meta}{\mathrm{meta}}
\newcommand{\fix}{\mathrm{fix}}

% Probability
\newcommand{\Prb}{\mathbb{P}}
\newcommand{\E}{\mathbb{E}}
\DeclareMathOperator{\Var}{Var}
\DeclareMathOperator{\Cov}{Cov}

% Operators
\DeclareMathOperator{\diag}{diag}
\DeclareMathOperator{\supp}{supp}

\newcommand{\argmax}{\operatorname*{arg\,max}}
\newcommand{\argmin}{\operatorname*{arg\,min}}
\newcommand{\ind}{\mathbf{1}}

% Inner product
\newcommand{\inner}[2]{\left\langle #1, #2 \right\rangle}

% Calligraphic sets
\newcommand{\cA}{\mathcal{A}}
\newcommand{\cS}{\mathcal{S}}
\newcommand{\cT}{\mathcal{T}}
\newcommand{\cF}{\mathcal{F}}
\newcommand{\cE}{\mathcal{E}}

% RL notation
\newcommand{\SSS}{\mathcal{S}}
\newcommand{\AAA}{\mathcal{A}}
\newcommand{\PP}{\mathcal{P}}
\newcommand{\RRR}{\mathcal{R}}
\newcommand{\gammaRL}{\gamma}

% Confidence bounds
\newcommand{\olQ}{\overline{Q}}
\newcommand{\ulQ}{\underline{Q}}
\newcommand{\olV}{\overline{V}}
\newcommand{\ulV}{\underline{V}}

% Misc
\newcommand{\eps}{\varepsilon}
\newcommand{\dt}{\mathrm{d}t}
\newcommand{\given}{\,\bigl\vert\,}

% Fractions
\newcommand{\half}{\tfrac12}
\newcommand{\thalf}{\tfrac32}

% Notes
\newcommand{\note}[1]{\paragraph{Note.} #1}

\theoremstyle{remark}
\newtheorem{conjecture}[theorem]{Conjecture}


\newcommand{\todo}[1]{\textcolor{red!70!black}{[\textbf{TODO:} #1]}}
\newcommand{\Cold}{\mathsf{C}}
\newcommand{\Exp}{\mathsf{E}}
\newcommand{\Hyb}{\mathsf{H}}
\newcommand{\wid}{w}

\title{Information Value Is Not Policy Value:\\ When LLM Warm-Starts Help Linear Contextual Bandits}


\author{%
  David S.~Hippocampus\thanks{Use footnote for providing further information
    about author (webpage, alternative address)---\emph{not} for acknowledging
    funding agencies.} \\
  Department of Computer Science\\
  Cranberry-Lemon University\\
  Pittsburgh, PA 15213 \\
  \texttt{hippo@cs.cranberry-lemon.edu} \\
}


\begin{document}
\maketitle
\begin{abstract}
A recent line of work reports ``reward multiplication'': routing the early
rounds of a linear contextual bandit to a frozen LLM policy and then handing
control to LinUCB can beat both the LLM and cold-started LinUCB, sometimes by a
large factor. We give an exact accounting of when this can happen and of how
large it can be. The regret difference between cold LinUCB and the handoff
hybrid splits into the expert's direct advantage during the routed rounds plus
the value of the data those rounds leave in LinUCB's design matrix; the
difference between the expert and the hybrid is the adaptive learner's
advantage on the tail. Data has value only through the confidence widths it
lowers, widths fall only through coverage, and coverage has a price: the least
regret any available action pays per unit of coverage of a direction. Our
main theorem bounds the online multiplier over cold LinUCB, on any instance
with a finite arm set, by the ratio of the price cold exploration pays for the
directions the warm start certifies to the price the expert pays for the same
coverage, plus cold LinUCB's regret on the arms the warm start does not
certify. A diverging online multiplier therefore requires a diverging price
ratio: an action of vanishing gap that covers a direction cold exploration can
buy only dearly. The bound is attained, through either channel, on instances
where such probes exist; and on the canonical separating instance of the
literature, which multiplies offline by a factor of the horizon, no online
handoff realizes anything, because every useful observation there costs a
constant amount of regret. A catalog of eleven instances, each removing one
ingredient, shows that recurrence, cheapness, correct targeting, and the
expert's suboptimality on the tail are each necessary.
We then show that this price accounting is \emph{complete} under the
model it assumes: when rewards are linear, the response sum of a warm start
is determined by its Gram matrix, so two experts whose routed rounds cover the
same directions leave warm starts with the same law. Any difference between
equally covering experts is therefore a departure from the linear model, and
we exhibit the departure in closed form: a \emph{transfer} term, the ridge
projection of the model's residuals onto the actions the expert chose, which
is the only expert-dependent part of the estimate at handoff, and which we
bound. A warm start therefore acts through two channels, coverage and
transfer, and on real data the transfer is the larger. On AmazonCat-13K three
LLM experts of equal accuracy raise the router's accuracy over cold LinUCB by
factors of $1.71$, $2.72$ and $3.29$, which is $6.5\%$, $16\%$ and $21\%$ less
regret. Coverage alone saves about $5\%$, the same for all three experts, as
the theory requires; the rest, and every significant difference between the
experts, comes from the transfer, and removing it by any of three independent
mechanisms collapses those differences. Which expert makes the better warm
start is decided by the residuals its choices carry, not by its accuracy.
\end{abstract}


\section{Introduction}
\label{sec:intro}

Sequential decision making studies how an agent adapts to an unknown
environment through repeated interaction \cite{zhu2022largeaction,zhu2022contextual},
with applications in recommendation \cite{li2010contextual,agarwal2016making},
healthcare \cite{tewari2017ads,svensson2023sequential}, and dialogue
\cite{li2016deep}. Large pretrained language models supply prior knowledge that
can guide early decisions
\cite{brown2020language,ouyang2022training,achiam2023gpt,
alamdari2024jumpstarting,abbasiyadkori2024believe}.

A striking observation is that combining a frozen LLM with an adaptive linear
bandit can strictly dominate either component, sometimes by orders of magnitude
\cite{chen2024llmdecision,huangcontextual}. Yet bandit warm-starting
\cite{bastani2021mostly,zhang2019warm} and confidence-gated model selection
\cite{auer2002nonstochastic,agarwal2017corralling,pacchiano2020model} are older
objects, and that literature does not report comparable multipliers. The
tension dissolves once one separates two things that are easy to conflate:
the \emph{policy value} of an expert (how much reward its actions earn) and
the \emph{information value} of an expert (how much its actions teach an
adaptive learner that keeps their covariance). The two can be wildly
different. A separating instance we construct to model the phenomenon
documented by \cite{chen2024llmdecision}, stated as
Proposition~\ref{prop:three-way}, is the cleanest illustration:
a constant policy that is wrong on every round, and whose regret is therefore
linear, produces a design matrix under which LinUCB's regret collapses from
linear to constant. The expert there has no policy value at all; its data has
enormous information value.

This paper makes that separation exact and uses it to say when a handoff from
a frozen expert to LinUCB can multiply reward. Writing $\Cold$ for cold LinUCB,
$\Exp$ for the frozen expert, and $\Hyb_\tau$ for the hybrid that follows
$\Exp$ for $\tau$ rounds while feeding every observation to LinUCB and then
hands over, we have the identities
\[
R_T(\Cold)-R_T(\Hyb_\tau)=A_\tau+V_\tau,
\qquad
R_T(\Exp)-R_T(\Hyb_\tau)=B_\tau,
\]
where $A_\tau$ is the expert's reward advantage over cold LinUCB on the routed
rounds, $V_\tau$ is the regret LinUCB saves after the handoff because of the
data the expert left behind, and $B_\tau$ is the adaptive learner's advantage
over the frozen expert on the tail (Proposition~\ref{prop:decomp}). Reward
multiplication over both base policies means $A_\tau+V_\tau$ and $B_\tau$ are
simultaneously large. The first is where the trade-off lives: the expert's
actions must be cheap (small regret while routed) and informative (large
$V_\tau$), and on many instances those two requirements pull in opposite
directions.

What governs $V_\tau$ is a round-level quantity that is standard in the
analysis of optimistic algorithms but rarely isolated. On the confidence event,
the instantaneous regret of LinUCB at round $t$ with design matrix $V$ is at
most
\[
\Phi_t(V)\;:=\;\max_{a\in\mathcal A_t}\min\bigl\{\Delta_{t,a},\;2\beta_t\|x_{t,a}\|_{V^{-1}}\bigr\},
\]
the largest gap among competitors whose confidence width still exceeds their
gap (Lemma~\ref{lem:width-limited}). Everything that the prior literature
attributes to ``dimension,'' ``coherence,'' or ``gap structure'' enters through
$\Phi_t$: the gap through the first argument of the minimum; the dimension
through how many distinct directions can keep a width large; coherence
through how much one observation shrinks the width of another direction, which
is an exact Sherman--Morrison identity (Lemma~\ref{lem:transfer}); and a warm
start through the Loewner ordering $(\lambda I+M)^{-1}\preceq(\lambda I)^{-1}$
(Lemma~\ref{lem:loewner}). We therefore do not posit separate structural axes.
We derive them.

Two consequences follow. First, the determinant-ratio (elliptical potential)
bound that underlies every $\widetilde O(\sqrt{rT})$-type guarantee cannot by
itself certify a multiplicative gain: with $\tau=o(T)$ the ratio of the warm
and cold determinant bounds tends to one, so any order-of-magnitude effect must
come from competitors whose widths the warm start pushes \emph{below their
gaps}, exactly as in Proposition~\ref{prop:three-way}
(Proposition~\ref{thm:handoff} and the discussion after it). Second, the offline
separation of Proposition~\ref{prop:three-way} has no online counterpart on
the same instance: for \emph{every} expert policy and every $\tau$, the
handoff hybrid has $V_\tau=0$, because each round introduces a fresh direction
that no earlier observation can cover (Proposition~\ref{prop:one-pass}).
Third, a warm start $(M,c)$ enters LinUCB through two objects, the design
matrix and the response sum, and under linear rewards the second is a
deterministic function of the first plus noise whose law depends only on the
first (Proposition~\ref{prop:response-sum}). Coverage is therefore the whole
of information value in the model our theory assumes, and any difference
between equally covering experts is a misspecification effect with an explicit
form (Section~\ref{sec:two-channels}). On real data this second channel is not
a correction: it carries most of the gain.

Widths fall only through coverage, and coverage has a price: for a direction
$u$, the least regret any action in the current set pays per unit of
$\langle x,u\rangle^2$ (Equation~\eqref{eq:price}). No policy, expert or
learner, obtains coverage for less (Lemma~\ref{lem:no-free-coverage}), and
cold LinUCB pays for each direction at the price of the arm optimism forces it
to pull, stopping once that arm's width falls below its gap. Our main theorem
(Theorem~\ref{thm:general-price}) puts the two facts together: on any
instance with a finite arm set, the multiplier of a handoff over cold LinUCB
is at most the ratio of cold LinUCB's price to the expert's price on the
directions the warm start certifies, plus cold LinUCB's regret on the arms
it does not certify, divided by the acquisition regret. The confidence radius
cancels. The theorem does not care whether the saving is collected during
the routed rounds ($A_\tau$) or afterwards ($V_\tau$), nor whether the
routing is fixed or adaptive, since it holds for any comparison policy at
any cutoff, and it makes the shape of the answer precise: an online handoff multiplies by a
\emph{constant} that is a ratio of prices, and the constant diverges only
if the expert has access to actions of vanishing gap that cover directions
cold exploration can buy only dearly (Corollary~\ref{cor:diverge}). Offline
warm starts, whose acquisition is not charged, are the only route to a
multiplier of order the horizon.

We then show that the bound is attained. On an instance in which the
directions that trap cold LinUCB recur and a cheap ``probe'' covering them is
offered early, a frozen expert that plays the probes and hands off has regret
smaller than cold LinUCB's by a factor of order $c^2/(\eps L^2)$, the ratio of
the bad arm's price $1/L^2$ to the probe's price $\eps/c^2$, and smaller than
the expert's own by a factor that grows with the horizon
(Theorem~\ref{thm:price}); the same ratio is realized on the tail when the
routed window keeps cold LinUCB busy elsewhere
(Proposition~\ref{prop:cat-tail-saving}). Making the probe dear collapses the
multiplier to at most two. A catalog of eleven instances
(Section~\ref{sec:catalog}, Table~\ref{tab:catalog}), all built from the
canonical one by changing a single ingredient, isolates what each clause of
this mechanism contributes, and each experimental configuration of
Section~\ref{sec:experiments} has a row in it.

\paragraph{Contributions.}
\begin{itemize}[itemsep=1pt, topsep=1pt, leftmargin=*]
\item \textbf{A price bound for the online multiplier}
(Theorem~\ref{thm:general-price}): on any finite-arm instance the advantage
of a handoff over cold LinUCB is at most a ratio of coverage prices on the
certified directions plus the uncertified part of cold LinUCB's regret, so
diverging online multipliers require vanishing prices
(Corollary~\ref{cor:diverge}).
\item \textbf{An exact decomposition of the hybrid's advantage}
(Proposition~\ref{prop:decomp}) into the expert's direct advantage $A_\tau$,
the value of the transferred data $V_\tau$, and the adaptive tail advantage
$B_\tau$, with the round-level quantity that governs $V_\tau$
(Lemma~\ref{lem:width-limited}) and exact identities for how warm starts and
ordinary exploration reduce widths (Lemmas~\ref{lem:loewner}
and~\ref{lem:transfer}).
\item \textbf{A separation of information value from policy value.}
Proposition~\ref{prop:three-way} (offline transfer by a factor of the
horizon) is paired with Proposition~\ref{prop:one-pass} (no online transfer
on the same instance, for every expert, randomized or adaptive).
\item \textbf{Instances attaining the price bound} through each channel
(Theorem~\ref{thm:price}, Proposition~\ref{prop:cat-tail-saving}), and a
catalog of eleven instances with complete proofs (Section~\ref{sec:catalog},
Appendix~\ref{app:catalog}) that shows which ingredient each negative case is
missing and places each experimental configuration.
\item \textbf{Two channels, and the transfer in closed form}
(Section~\ref{sec:two-channels}): under linear rewards the response sum of a
warm start is determined by its Gram matrix, so the price ratio is a complete
accounting of information value; under misspecification the estimate acquires
a \emph{transfer} term (Proposition~\ref{prop:transfer}), the ridge
projection of the model's residuals onto the expert's chosen actions, which is
the only expert-dependent part of the warm start, together with a regret bound
that holds against any reference parameter
(Proposition~\ref{prop:transfer-regret}).
\item \textbf{Experiments showing that the transfer is the larger channel in
practice} (Section~\ref{sec:experiments}), on two datasets and twelve
dataset/model configurations (ten evaluation seeds per row, with every
confidence parameter chosen on a disjoint block of ten tuning seeds). We measure the decomposition directly and report both the accuracy
multiplier that the literature calls reward multiplication and the regret
ratio that the theorems bound. Coverage delivers a small gain that is the same
for every expert; the transfer carries most of the gain and every significant
difference between experts of equal accuracy.
\end{itemize}
A detailed related-work discussion is in Appendix~\ref{sec:related}.


\section{Preliminaries}
\label{sec:preliminaries}

We work in a standard linear contextual bandit
\cite{abbasi2011improved,chu2011contextual,lattimore2020bandit}. At each round
$t=1,\dots,T$, the learner observes a finite action set $\mathcal A_t$ where
each $a\in\mathcal A_t$ has a feature vector $x_{t,a}\in\mathbb R^d$ with
$\|x_{t,a}\|_2\le L$, selects $A_t\in\mathcal A_t$, and observes
$r_t=\langle\theta^\star,x_{t,A_t}\rangle+\eta_t$, where $\theta^\star$ is
unknown with $\|\theta^\star\|_2\le S$ and $\eta_t$ is conditionally mean-zero
and $\sigma$-sub-Gaussian. Write $\mu_t(a)=\langle\theta^\star,x_{t,a}\rangle$,
$a_t^\star\in\arg\max_a\mu_t(a)$, and $\Delta_{t,a}=\mu_t(a_t^\star)-\mu_t(a)$
for the suboptimality gap of $a$ at round $t$. For a policy $\pi$ and an
interval $I\subseteq[T]$ of rounds, $R_I(\pi)=\sum_{t\in I}\Delta_{t,\pi_t}$ is
its pseudo-regret on $I$, and $R_T(\pi)=R_{[T]}(\pi)$. The language aspect of
the applications we target enters only through the feature map: context and
action embeddings from pretrained models are combined into $x_{t,a}$, and the
frozen expert may be an LLM. Nothing in the theory uses language structure.

\paragraph{LinUCB and its warm start.}
LinUCB maintains $V_t=V_1+\sum_{s<t}x_{s,A_s}x_{s,A_s}^\top$,
$b_t=b_1+\sum_{s<t}r_sx_{s,A_s}$, $\widehat\theta_t=V_t^{-1}b_t$, and plays
$A_t\in\arg\max_a\langle\widehat\theta_t,x_{t,a}\rangle+\beta_t\|x_{t,a}\|_{V_t^{-1}}$,
where $\|x\|_{V^{-1}}=\sqrt{x^\top V^{-1}x}$ is the confidence width of $x$ and
$\beta_t$ is the confidence radius. \emph{Cold-started} LinUCB, written
$\Cold$, uses $V_1=\lambda I_d$, $b_1=0$. A \emph{warm start} $(M,c)$ uses
$V_1=\lambda I_d+M$ and $b_1=c$ with $M=\sum_sx_sx_s^\top\succeq0$ and
$c=\sum_sr_sx_s$ the Gram matrix and response sum of some logged data. We say
the warm start is \emph{offline} if that data is not counted in the evaluated
horizon and \emph{online} if it is.

\paragraph{Frozen expert and handoff hybrid.}
A frozen expert $\Exp$ proposes $a_t^{\Exp}\in\mathcal A_t$ at each round by a
rule that is never updated (for an LLM, the action whose embedding is closest
to the LLM's output; see Section~\ref{sec:exp-setup}). For $\tau\in\{0,\dots,T\}$
the \emph{handoff hybrid} $\Hyb_\tau$ plays $a_t^{\Exp}$ for $t\le\tau$, updates
$(V_t,b_t)$ with every played action and observed reward including those
rounds, and plays LinUCB for $t>\tau$. Thus $\Hyb_0=\Cold$ and $\Hyb_T=\Exp$,
and the LinUCB phase of $\Hyb_\tau$ is LinUCB with the online warm start
$(M_\tau,c_\tau)=\bigl(\sum_{t\le\tau}x_{t,A_t}x_{t,A_t}^\top,\sum_{t\le\tau}r_tx_{t,A_t}\bigr)$.
The deterministic router used in our experiments
(Appendix~\ref{app:meta-algorithm}) chooses $\tau$ adaptively and may
interleave the two; the theory below is for a fixed $\tau$, which is the
right object for understanding why a multiplier appears, and we comment on
adaptive $\tau$ after Proposition~\ref{thm:handoff}.


\section{Main Results}
\label{sec:main-results}

\subsection{The Decomposition}
\label{sec:catalog}

Everything in this paper is organized by one identity.

\paragraph{Setup for the decomposition.}
Recall the three policies: $\Cold$ is LinUCB initialized from $\lambda I_d$
(no prior data); $\Exp$ is a frozen expert that proposes actions by a fixed
rule (e.g.\ an LLM) and is never updated; and $\Hyb_\tau$ is the handoff
hybrid that follows $\Exp$ for the first $\tau$ rounds, feeds every
observation into LinUCB's state $(V_t,b_t)$, and then hands control to LinUCB
for rounds $\tau{+}1,\dots,T$. For any policy $\pi$ and interval $I$, we
write $R_I(\pi)=\sum_{t\in I}\Delta_{t,\pi_t}$ for its pseudo-regret on $I$,
and $R_T(\pi)=R_{[T]}(\pi)$.

\begin{proposition}[Decomposition of the hybrid's advantage]
\label{prop:decomp}
For any instance, any expert $\Exp$, any $\tau$, and any realization of the
noise,
\begin{align*}
R_T(\Cold)-R_T(\Hyb_\tau)
&=\underbrace{R_{[1,\tau]}(\Cold)-R_{[1,\tau]}(\Exp)}_{A_\tau}
 +\underbrace{R_{[\tau+1,T]}(\Cold)-R_{[\tau+1,T]}(\Hyb_\tau)}_{V_\tau},\\
R_T(\Exp)-R_T(\Hyb_\tau)
&=\underbrace{R_{[\tau+1,T]}(\Exp)-R_{[\tau+1,T]}(\Hyb_\tau)}_{B_\tau}.
\end{align*}
\end{proposition}

$A_\tau$ is the expert's \emph{direct} advantage: positive if the expert earns
more than cold LinUCB would have on the routed rounds, negative if the expert's
actions cost more regret than cold exploration. $V_\tau$ is the
\emph{information} value of the routed rounds: it is nonzero only through the
warm start $(M_\tau,c_\tau)$ that they leave in LinUCB's state, since apart
from that state $\Hyb_\tau$ and $\Cold$ run the same algorithm on the same
future. $B_\tau$ is the adaptive learner's advantage on the tail and is large
only when the expert remains substantially suboptimal late while the
warm-started learner does not. We write $R_{[1,\tau]}(\Exp)=\sum_{t\le\tau}\Delta_{t,a_t^{\Exp}}$
for the \emph{acquisition regret} of the warm start: the price, counted in the
online horizon, of the data in $M_\tau$.

\begin{remark}
The identities of Proposition~\ref{prop:decomp} hold pathwise in the
reward-table model of Section~\ref{sec:template}, in which two policies can
be run on the same noise, and therefore also in expectation. They also hold
when the routed rounds are not an initial segment: for any set
$\mathcal T\subseteq[T]$ of rounds on which the hybrid follows the expert,
fixed in advance or chosen adaptively (as by the router of
Appendix~\ref{app:meta-algorithm}), replacing $[1,\tau]$ by $\mathcal T$ and
$[\tau+1,T]$ by its complement gives the same two identities, because each
side is a difference of sums over a partition of $[T]$. In particular,
$R_T(\Hyb_\tau)\le\min\{R_T(\Cold),R_T(\Exp)\}$ if and only if
$A_\tau+V_\tau\ge0$ and $B_\tau\ge0$, and
$\min\{R_T(\Cold),R_T(\Exp)\}/R_T(\Hyb_\tau)\to\infty$ along a sequence of
instances---formally, a sequence
$(d_T,\theta^\star_T,\{\mathcal A^{(T)}_t\}_{t\le T},\Exp_T,\tau_T)_{T\to\infty}$,
which is how every diverging ratio in this paper is to be read---if and only
if both $A_\tau+V_\tau$ and $B_\tau$ exhaust the corresponding base regret up
to lower-order terms. If instead the warm start
$(M,c)$ is built from data outside the horizon (an \emph{offline} warm start
$\Hyb^{\mathrm{off}}$, LinUCB started from $(\lambda I+M,c)$), the routed
rounds are not charged and the identities read
$R_T(\Cold)-R_T(\Hyb^{\mathrm{off}})=:V^{\mathrm{off}}$ and
$R_T(\Exp)-R_T(\Hyb^{\mathrm{off}})=:B^{\mathrm{off}}$, with no $A$ term.
\end{remark}

The proposition is bookkeeping and says nothing by itself. Its content is in
what the three terms can be. Reward multiplication over both base policies
means $A_\tau+V_\tau$ is almost all of $R_T(\Cold)$ and $B_\tau$ is almost all
of $R_T(\Exp)$, so every question about when a handoff multiplies is a
question about which of $A_\tau$, $V_\tau$, $B_\tau$ can be large at once. We
answer it with eleven instances, each obtained from the canonical one
(Proposition~\ref{prop:three-way}) by changing a single ingredient
(Table~\ref{tab:catalog}); exactly two of them multiply over both base
policies, they realize the same price ratio through the two channels
$A_\tau$ and $V_\tau$ (Theorem~\ref{thm:price},
Proposition~\ref{prop:cat-tail-saving}), and Theorem~\ref{thm:general-price}
says that no instance can do better than a price ratio. The pattern across
the rows is the following.


\begin{center}
\fbox{\parbox{0.95\linewidth}{\textbf{Thesis.}
A handoff multiplies regret savings over cold LinUCB if and only if the
expert's routed actions buy coverage of the directions cold LinUCB would
otherwise pay for, at a lower price than cold exploration pays, early enough
for those directions to recur; and it multiplies over the expert if and only
if, in addition, the expert is worse than the warm-started learner on the tail.
}}
\end{center}

The catalog has eleven rows rather than the $2^3=8$ of a binary table
because several terms admit qualitatively distinct failure modes---for
instance, $V_\tau$ can vanish because the expert buys the same directions cold
exploration buys at the same price (row~4), because it buys the wrong
directions (row~7), or because it buys the right directions too dearly
(row~8)---because the offline/online distinction (rows~1--2) and the two
ways of killing the tail term (rows~3 and~6) each warrant separate entries,
and because the positive case appears twice, with the saving collected in
the routed window (row~5) and on the tail (row~$5'$).

\begin{table}[t]
\centering
\scriptsize
\setlength{\tabcolsep}{2.5pt}
\begin{tabular}{@{}clccccl@{}}
\toprule
\# & Instance & $A_\tau$ & $V_\tau$ & $B_\tau$ & Ratio to $\Cold$ & Ingredient isolated\\
\midrule
1 & canonical, offline (Prop.~\ref{prop:three-way}) & not charged & $V^{\mathrm{off}}\asymp T$ & $B^{\mathrm{off}}\asymp T$ & $\asymp T$ & data is free\\
2 & canonical, online (Prop.~\ref{prop:one-pass}) & $\le\tau$ & $0$ & $\le0$ & none & directions never recur\\
3 & canonical, oracle (Cor.~\ref{cor:cat-oracle}) & $\tau$ & $0$ & $-(T-\tau)$ & additive & expert too good\\
4 & recurring, bad-arm expert (Prop.~\ref{prop:cat-mimic}) & $0$ & $0$ & $\ge0$ & $\le961$ & same price as cold\\
5 & recurring, cheap probes (Thm.~\ref{thm:price}) & $\ge Km_0-\eps\tau_0$ & $\ge0$ & $T-\tau_0$ & $\asymp c^2/(\eps L^2)$ & price ratio, in the window\\
$5'$ & as 5, cold kept busy in the window (Prop.~\ref{prop:cat-tail-saving}) & $-\eps\tau_0$ & $\ge Km_0$ & $T-\tau_0$ & $\asymp c^2/(\eps L^2)$ & price ratio, on the tail\\
6 & as 5, expert optimal on tail (Prop.~\ref{prop:cat-tail}) & as 5 & as 5 & $0$ & as 5 & expert optimal late\\
7 & nuisance probes (Prop.~\ref{prop:cat-wrong}) & & $<R_{[\tau_0+1,T]}(\Cold)$ & & $\le65$ & wrong directions\\
8 & as 5 with $\eps\ge c^2/L^2$ (Thm.~\ref{thm:price}) & & & & $\le2$ & probe too dear\\
9 & cold already cheap (Prop.~\ref{prop:cat-cheap}) & & & & $o(1)$ avg.\ reward & nothing to buy\\
10 & bad arms with gap $\eta$ (Prop.~\ref{prop:cat-gap}) & & & & $\le1+c^2/(\eps L^2)$ & gap rescales prices\\
\bottomrule
\end{tabular}
\caption{The catalog: what the three terms of Proposition~\ref{prop:decomp} can be.
Row 1 is the offline identity, in which $A$ is not charged; rows 2--10 are
online handoffs $\Hyb_\tau$. Rows 3 and 6 kill $B_\tau$ from the two sides.
Rows 4, 7, 8 kill the multiplier over cold LinUCB in three different ways:
the expert buys what cold buys at the same price, buys directions nobody will
need, or buys the right directions too dearly. Row 9 is the case in which cold
LinUCB's own bill is already negligible in average reward. Row 10 shows that
the gap is a rescaling of both prices. Rows 5 and $5'$ are the positive
cases; they realize the same price ratio, collected as $A_\tau$ in row 5 and
as $V_\tau$ in row $5'$, and Theorem~\ref{thm:general-price} shows no
instance exceeds a price ratio. Blank cells are not claimed. Constants are
those proved in Appendix~\ref{app:catalog}, not optimized.}
\label{tab:catalog}
\end{table}


\subsection{Rows 1--3: The Canonical Instance}
\label{sec:canonical}

All proofs are in Appendices~\ref{app:proofs} and~\ref{app:catalog}.


\paragraph{Rows 1--2: the canonical instance, offline and online.}
The canonical instance, which we construct to model the separation reported
by \cite{chen2024llmdecision}, uses
horizon $n$, a single reward direction $g=e_1$, and $n$ orthogonal ``bad''
arms $b_t=Le_{t+1}$ with $L>1$.

\begin{proposition}[Offline three-way separation]
\label{prop:three-way}
Fix $n\ge1$, $L>1$, $\lambda>0$. There exist $d\ge n+1$, $\theta^\star$, and
two-action sets $\{\mathcal A_t\}_{t=1}^n$ such that, for any sub-Gaussian
noise: (1) a constant policy $\pi_{\mathrm{const}}$ has pseudo-regret exactly
$n$; (2) cold LinUCB also incurs regret $n$, selecting the same action as
$\pi_{\mathrm{const}}$ each round; (3) for $m$ large enough, LinUCB
warm-started with the design matrix that $\pi_{\mathrm{const}}$ generates from
$mn$ \emph{offline} rollouts incurs $O(1)$ regret with probability $\ge1-\delta$.
\end{proposition}

The proof (Appendix~\ref{app:proof-three-way}) is short and worth reading
through the lens of Section~\ref{sec:preliminaries}. The true parameter is as
simple as possible: one reward-bearing direction. Cold LinUCB fails anyway,
because at each round the fresh arm $b_t$ has width $\|b_t\|_{V_t^{-1}}=L/\sqrt\lambda$
regardless of what was played before, while $g$ has width $1/\sqrt\lambda$ and
estimated value $0$, so optimism picks $b_t$ every time. The offline design
matrix $M=mL^2\sum_{j\ge2}e_je_j^\top$ does not teach LinUCB that $g$ is good;
its response sum is zero in expectation. It teaches LinUCB that the
\emph{alternatives} are not uncertain: $\|b_t\|^2_{(\lambda I+M)^{-1}}=L^2/(\lambda+mL^2)$,
which falls below the gap once $m$ is moderate, and LinUCB then plays $g$.
In the language of Proposition~\ref{prop:decomp} this is the offline identity with
$V^{\mathrm{off}}\asymp T$ and $B^{\mathrm{off}}\asymp T$; the acquisition of
the data is not charged.

The same instance is hostile to online acquisition. Every round introduces a
direction that no earlier observation touches, so the data any expert can
collect during the first $\tau$ rounds says nothing about the arms that appear
after round $\tau$.

\begin{proposition}[No online transfer on the canonical instance]
\label{prop:one-pass}
Consider the instance of Proposition~\ref{prop:three-way} with horizon $T$
and confidence radius $\beta$, and let $\kappa:=\sigma\sqrt{2\log(2T/\delta)}/\sqrt\lambda$.
If $\beta(L-1)/\sqrt\lambda>1+\kappa$, then with probability at least
$1-\delta$, for \emph{every} expert policy $\Exp$ and every $\tau\in\{0,\dots,T\}$,
the handoff hybrid $\Hyb_\tau$ plays $b_t$ at every round $t>\tau$. Consequently
$R_{[\tau+1,T]}(\Hyb_\tau)=R_{[\tau+1,T]}(\Cold)=T-\tau$ and $V_\tau=0$.
\end{proposition}

The condition on $\beta$ is the same optimism condition that makes cold LinUCB
fail in the first place (it only strengthens the inequality
$\beta L/\sqrt\lambda>\beta/\sqrt\lambda$ by the largest value the estimate of
$g$ can take). The proposition says that the offline separation of
Proposition~\ref{prop:three-way} cannot be realized by any online router on
this instance: not by a bad expert, whose data has $V_\tau=0$, and not by a
good one, which has $B_\tau\le0$ because it is already optimal. Online
multiplication needs a third kind of expert, whose routed actions are cheap
\emph{and} cover directions that will matter later; row 5 constructs one.


\paragraph{Row 3: an oracle expert on the canonical instance.}
\begin{corollary}[Oracle expert on the canonical instance]
\label{cor:cat-oracle}
On the instance of Proposition~\ref{prop:one-pass}, under its hypothesis, let
$\Exp$ play $g$ at every round. With probability at least $1-\delta$, for every
$\tau$: $R_T(\Exp)=0$, $R_T(\Hyb_\tau)=T-\tau$, $A_\tau=\tau$, $V_\tau=0$,
$B_\tau=-(T-\tau)$. The hybrid improves on cold LinUCB by exactly the rounds
it routes and on the expert by nothing.
\end{corollary}


\subsection{Prices and the Main Theorem}
\label{sec:price}

Rows 4--10 share a vocabulary. The word ``price'' has a precise meaning. For a direction $u\in\R^d$ and an
action set $\mathcal A_t$, write
\begin{equation}
\label{eq:price}
\pi_t(u)\;:=\;\min\Bigl\{\frac{\Delta_{t,a}}{\langle x_{t,a},u\rangle^2}:\;a\in\mathcal A_t,\ \langle x_{t,a},u\rangle\neq0\Bigr\},
\qquad \pi(u):=\inf_t\pi_t(u),
\end{equation}
the least pseudo-regret any policy must pay at round $t$ to add one unit to
the coverage $\sum_s\langle x_{s,A_s},u\rangle^2$ of $u$. Coverage is what
reduces widths (Lemma~\ref{lem:loewner}), so $\pi(u)$ is the exchange rate
between regret and width reduction in direction $u$. The following lemma
says that no policy, expert or otherwise, gets coverage for less.

\begin{lemma}[No free coverage]
\label{lem:no-free-coverage}
For any policy $\pi$, any interval $I$ of rounds, and any direction $u$,
$\sum_{t\in I}\langle x_{t,\pi_t},u\rangle^2\le R_I(\pi)/\pi(u)$.
Consequently, if $\Hyb_\tau$ hands off with design matrix $V_{\tau+1}=\lambda I+M_\tau$, then
for every unit vector $u$,
$\|u\|^2_{V_{\tau+1}^{-1}}\ \ge\ 1\big/\bigl(\lambda+R_{[1,\tau]}(\Exp)/\pi(u)\bigr)$.
\end{lemma}

The lemma is one half of the main theorem. The other half is that cold
LinUCB buys coverage of each direction at the price of the arm optimism
forces it to pull, and stops once that arm's width falls below its gap.
Together they bound the online multiplier by a ratio of prices on the
directions the warm start certifies, plus cold LinUCB's regret on the
directions it does not.

\paragraph{Setting for the theorem.}
Let $\mathcal X\subset\R^d$ be a finite set of arm vectors with
$\|x\|_2\le L$, let $\mathcal A_t\subseteq\mathcal X$ for every $t$, and
suppose an arm $x^\star\in\mathcal X$ with
$\mu(x^\star)=\max_{x\in\mathcal X}\mu(x)$ belongs to every $\mathcal A_t$.
Then the gap of an arm is the same at every round,
$\Delta(x):=\mu(x^\star)-\mu(x)$, and we write
$S:=\{x\in\mathcal X:\Delta(x)>0\}$. For $x\in S$ put $u_x:=x/\|x\|_2$ and
\[
p(x):=\frac{\Delta(x)}{\|x\|_2^2},\qquad
c_x:=\frac{4\beta^2\|x\|_2^2}{\Delta(x)^2},\qquad
N_x:=\#\{t\in[T]:x\in\mathcal A_t\},
\]
where $\beta:=\beta_T$ is the largest confidence radius. Here $p(x)$ is the
price of covering $u_x$ by pulling $x$ itself, and $c_x$ is the coverage of
$u_x$ at which the width of $x$ equals half its gap, so that the confidence
bound alone excludes $x$ once the coverage $\sum_s\langle x_s,u_x\rangle^2$
exceeds $c_x-\lambda$. For a set $U$ of unit vectors write
$\mathrm{cov}_U(x):=\sum_{u\in U}\langle x,u\rangle^2$. The theorem compares
cold LinUCB with an arbitrary policy $\pi$ at an arbitrary cutoff
$t_0\in\{0,\dots,T\}$; the handoff hybrid $\Hyb_\tau$ with $t_0=\tau$ is the
case of interest, and then rounds $1,\dots,t_0$ are the expert's, but nothing
in the statement or proof uses this. Write $x^\pi_t$ for the arm $\pi$ plays
at round $t$ and define
\[
\pi_U(t_0):=\min\Bigl\{\frac{\Delta(x^\pi_t)}{\mathrm{cov}_U(x^\pi_t)}:
\ t\le t_0,\ \Delta(x^\pi_t)>0,\ \mathrm{cov}_U(x^\pi_t)>0\Bigr\},
\qquad
C^{\mathrm{free}}_U(t_0):=\sum_{t\le t_0:\,\Delta(x^\pi_t)=0}\mathrm{cov}_U(x^\pi_t),
\]
the price $\pi$ pays for coverage of $U$ up to $t_0$ (the minimum over an
empty set is $+\infty$, and then $1/\pi_U(t_0):=0$), and the coverage of $U$
that $\pi$ obtains for free by playing optimal arms. Finally let
$U_{\mathrm{cert}}(t_0)$ be the set of directions $u_x$, $x\in S$, with
$2\beta\|x\|_{(V^\pi_{t_0+1})^{-1}}<\Delta(x)$, where
$V^\pi_{t_0+1}=\lambda I+\sum_{t\le t_0}x^\pi_t (x^\pi_t)^\top$ is $\pi$'s design
matrix after round $t_0$: the directions $\pi$'s first $t_0$ rounds certify.
For $\pi=\Hyb_\tau$ and $t_0=\tau$ this is the hybrid's matrix at handoff,
$\pi_U(\tau)$ is the expert's price, and $R_{[1,t_0]}(\pi)$ is the acquisition
regret.

\begin{theorem}[Price bound for the online multiplier]
\label{thm:general-price}
In the setting above, on the confidence event of cold LinUCB, the following
hold for every policy $\pi$, every cutoff $t_0\in\{0,\dots,T\}$, and every
$U\subseteq U_{\mathrm{cert}}(t_0)$.
\begin{enumerate}[itemsep=1pt,topsep=1pt,leftmargin=*]
\item[(a)] \emph{(Cold LinUCB pays its own price, once per direction.)}
$R_T(\Cold)\le\sum_{x\in S}\Delta(x)\min\{m_x,N_x\}$, where
$m_x:=(c_x-\lambda)/\|x\|_2^2+1$ if $c_x\ge\lambda$ and $m_x:=0$ otherwise.
\item[(b)] \emph{(Certification costs at least the price.)}
\[
\sum_{u_x\in U}(c_x-\lambda)_+\ \le\ \sum_{u_x\in U}\sum_{t\le t_0}\langle x^\pi_t,u_x\rangle^2
\ \le\ \frac{R_{[1,t_0]}(\pi)}{\pi_U(t_0)}+C^{\mathrm{free}}_U(t_0).
\]
\item[(c)] \emph{(Price bound.)} With $p_U:=\max_{u_x\in U}p(x)$ and the
price ratio $P_U:=p_U/\pi_U(t_0)$,
\[
R_T(\Cold)\ \le\ P_U\,R_{[1,t_0]}(\pi)\ +\ p_U\,C^{\mathrm{free}}_U(t_0)\ +\ \sum_{u_x\in U}\Delta(x)
\ +\sum_{x\in S:\,u_x\notin U}\Delta(x)\min\{m_x,N_x\}.
\]
Since $R_T(\pi)\ge R_{[1,t_0]}(\pi)$, whenever $R_{[1,t_0]}(\pi)>0$,
\begin{equation}
\label{eq:general-price}
\frac{R_T(\Cold)}{R_T(\pi)}\ \le\ P_U\ +\
\frac{p_U\,C^{\mathrm{free}}_U(t_0)+\sum_{u_x\in U}\Delta(x)+\sum_{x\in S:\,u_x\notin U}\Delta(x)\min\{m_x,N_x\}}{R_{[1,t_0]}(\pi)} .
\end{equation}
For $\pi=\Hyb_\tau$, $t_0=\tau$, this bounds the multiplier of the handoff
hybrid by the price ratio between cold exploration and the expert on the
directions the warm start certifies.
\end{enumerate}
\end{theorem}

The proof is in Appendix~\ref{app:proof-general-price}; part (a) is the
width-limited lemma of Section~\ref{sec:width} combined with the
Sherman--Morrison identity, part (b) is Lemma~\ref{lem:no-free-coverage}
summed over the certified directions, and (c) is arithmetic. Only the
confidence event of cold LinUCB is used: certification is a property of the
matrix $V^\pi_{t_0+1}$ and (c) does not require $\pi$ to exploit it. Because
$\pi$ is arbitrary, the bound covers the fixed-$\tau$ hybrid, the adaptive
router of Appendix~\ref{app:meta-algorithm} (take $t_0$ to be its last routed
round), a CORRAL-style master, and cold LinUCB itself; every one of them
beats cold LinUCB by at most the price ratio of the coverage it has certified
by $t_0$, plus the uncertified residual. We stress that
\eqref{eq:general-price} bounds the ratio $R_T(\Cold)/R_T(\pi)$ of the two
policies' regrets on the same instance, rather than comparing an upper bound
for one against an upper bound for the other: two policies sharing an
asymptotic upper bound may still differ by a diverging factor, so a
statement about the multiplier cannot be obtained from matching rates and
must bound one policy's regret in terms of the other's, as (c) does.

\begin{remark}[Two relaxations]
\label{rem:relaxations}
(i) The factor $2$ in the definition of $U_{\mathrm{cert}}$ is the crude
one. On the confidence event with the true self-normalized radius
$\beta^{\mathrm{true}}_t\le\beta$, the proof of Lemma~\ref{lem:width-limited}
gives $\Delta_{t,A_t}\le(\beta_t+\beta^{\mathrm{true}}_t)\|x_{t,A_t}\|_{V_t^{-1}}$
for the played arm, so the theorem holds with $c_x$ replaced by
$(\beta+\beta^{\mathrm{true}})^2\|x\|_2^2/\Delta(x)^2$ throughout and
$U_{\mathrm{cert}}$ enlarged to the directions with
$(\beta+\beta^{\mathrm{true}})\|x\|<\Delta(x)$; $P_U$ is unchanged. (ii) If the
optimal arm is not offered at every round, so that the gap of $x$ varies
with $t$, write $\Delta^-(x)$ and $\Delta^+(x)$ for its smallest and largest
values over rounds at which $x$ is offered. Part (a) holds with $c_x$ computed
from $\Delta^-(x)$ and each pull charged $\Delta^+(x)$, and (b) with
$\Delta^-(x)$; the price ratio then carries a factor $\Delta^+(x)/\Delta^-(x)$,
which is $1$ under the standing assumption.
\end{remark}

The residual in \eqref{eq:general-price} has three parts: the coverage the
expert obtains for free by playing optimal arms (zero when the optimal arm is
orthogonal to the certified directions, as on every instance of this
section), one pull of each certified arm, and cold LinUCB's regret on the arms
whose widths the warm start does not push below their gaps. Only the last can
be large, and it is exactly the part of the hybrid's advantage that is not
bought at a price: on those arms the hybrid faces the same confidence bound
as cold LinUCB, and no coverage argument lower-bounds its tail regret on them.
Theorem~\ref{thm:general-price} is therefore a statement about the certified
saving, and it holds whether that saving is collected in the routed window
($A_\tau$) or on the tail ($V_\tau$): rows~5 and~$5'$ below realize the same
price ratio through the two channels. The confidence radius cancels in $P_U$
because it scales $c_x$, the coverage both policies must reach, identically.

\begin{corollary}[Diverging online multipliers require vanishing prices]
\label{cor:diverge}
Consider a sequence of instances, policies $\pi$ and cutoffs $t_0$ along
which $R_T(\Cold)/R_T(\pi)\to\infty$ and, for some choice of
$U\subseteq U_{\mathrm{cert}}(t_0)$ on each instance, the residual in
\eqref{eq:general-price} is at most a fixed multiple of $R_{[1,t_0]}(\pi)$.
Then $P_U\to\infty$: the price $\pi_U(t_0)$ that $\pi$ pays for the certified
coverage vanishes relative to the price $p_U$ cold LinUCB pays for it, so
$\pi$ plays actions whose gap is negligible against the coverage they supply.
For the handoff hybrid this means the expert has access to such actions.
On the canonical instance of Proposition~\ref{prop:three-way} no online expert
certifies any direction that is offered after the handoff, so
$U_{\mathrm{cert}}(\tau)=\emptyset$ for every $\tau<T$ and all of cold LinUCB's
regret is in the residual; that is the price-theoretic content of
Proposition~\ref{prop:one-pass}. Offline, where the acquisition is not
charged, the same instance multiplies by a factor of the horizon.
\end{corollary}

\begin{proof}
From \eqref{eq:general-price}, $P_U\ge R_T(\Cold)/R_T(\pi)-C$ for the
fixed multiple $C$, so $P_U\to\infty$; since $p_U\le\max_{x\in S}p(x)$ is a
property of the instance, this is $\pi_U(t_0)\to0$ relative to $p_U$.
On the canonical instance the arms offered after round $\tau$ are
$b_t=Le_{t+1}$, $t>\tau$, and no arm played at or before round $\tau$ has a
component along $e_{t+1}$, so $\|b_t\|_{V_{\tau+1}^{-1}}=L/\sqrt\lambda$ and
$2\beta L/\sqrt\lambda>1=\Delta(b_t)$ under the hypothesis of
Proposition~\ref{prop:one-pass}; hence no such $b_t$ is certified.
\end{proof}

\paragraph{When do cheap probes exist?}
The positive rows need an action of small gap whose feature vector carries a
non-negligible component along a direction that cold LinUCB can otherwise
cover only by pulling an arm of large gap. Two observations make this less
special than it looks. First, the probe need not be a single designated arm:
any near-optimal action whose embedding overlaps the bad direction serves,
and in the embedding spaces used in practice (Section~\ref{sec:experiments})
near-optimal actions overlap many others. Second, the expert does not need to
know that the probe is informative; it needs only to prefer it, which any
policy with prior knowledge of the reward does when the probe is nearly as
good as the best arm. What the expert supplies, and cold optimism cannot, is
the \emph{order} in which coverage is bought: optimism ranks a fresh bad arm
above a cheap probe because the bad arm's width is larger, and pays the bad
arm's price first (Lemma~\ref{lem:cat-trap} in the appendix); a policy that
already believes the bad arm is bad takes the probe at once. The price ratio
$c^2/(\eps L^2)$ of the rows below is the exchange rate between those two
orders, and Theorem~\ref{thm:general-price} says it is the only exchange rate
available: an expert whose actions are not cheaper per unit of certified
coverage than cold exploration's cannot multiply, however accurate it is.

\subsection{The Two Channels of a Warm Start}
\label{sec:two-channels}

A warm start $(M,c)$ enters LinUCB through two objects: the design matrix
$V_1=\lambda I+M$, which fixes every confidence width, and the response sum
$b_1=c$, which fixes the estimate $\widehat\theta_1=V_1^{-1}c$. Everything so
far has bounded the value of a warm start through the first. We now separate
the two. The conclusion is that in the model of
Section~\ref{sec:preliminaries} the first is the whole story, so
Theorem~\ref{thm:general-price} is a complete accounting of information value
there; and that the second exists only when the reward is not linear in the
features, in which case it has an explicit form.

\begin{proposition}[Width-limited regret, two channels]
\label{prop:two-channel}
At a round $t$ at which LinUCB acts, let
$\rho_t:=\|\widehat\theta_t-\theta^\star\|_{V_t}$ be the realized estimation
error, so $\rho_t\le\beta_t$ on $\mathcal E_\delta$. Then on $\mathcal E_\delta$
\begin{equation}
\label{eq:phi-two}
\Delta_{t,A_t}\ \le\ \Phi_t(V_t,\rho_t)\ :=\
\max_{a\in\mathcal A_t}\min\bigl\{\Delta_{t,a},\ (\beta_t+\rho_t)\|x_{t,a}\|_{V_t^{-1}}\bigr\},
\end{equation}
and $\Phi_t(V,\rho)$ is non-increasing in $V$ and in $\rho$. A warm start acts
on the first argument through $M$ and on the second through $c$. Since
$0\le\rho_t\le\beta_t$ we have $\Phi_t(V_t,\rho_t)\ge\tfrac12\Phi_t(V_t)$: the
response sum can tighten the width-limited bound by at most a factor of two,
whereas the design matrix can reduce it by the price ratio, which is unbounded.
\end{proposition}

\begin{proposition}[The response sum is determined by the design matrix]
\label{prop:response-sum}
Let the warm start be built from pairs $(x_s,r_s)_{s\le\tau}$ with
$r_s=\langle\theta^\star,x_s\rangle+\eta_s$ and $\eta_s$ i.i.d.\
$\mathcal N(0,\sigma^2)$. Then
\[
c=M\theta^\star+\xi,\qquad \xi:=\sum_{s\le\tau}\eta_sx_s\sim\mathcal N(0,\sigma^2M),
\qquad
\widehat\theta_1-\theta^\star=-\lambda V_1^{-1}\theta^\star+V_1^{-1}\xi .
\]
Hence the law of $(V_1,b_1)$ given $M$ does not depend on which policy
generated the logged actions: two experts, or two routed sets, with the same
Gram matrix leave warm starts with the same law, and the tail trajectories of
the corresponding hybrids---hence $V_\tau$ and $B_\tau$---have the same law.
\end{proposition}

For $\sigma$-sub-Gaussian noise the distributional identity is replaced by the
observation that the self-normalized bound of \cite{abbasi2011improved}, the
only property of $\xi$ used in this paper, depends on the features only
through $M$; every bound in Sections~\ref{sec:width}--\ref{sec:handoff} is
therefore unchanged when one expert is exchanged for another of equal coverage.

\begin{corollary}[Coverage is complete under linear rewards]
\label{cor:coverage-complete}
Under the model of Section~\ref{sec:preliminaries}, two experts whose routed
rounds leave the same Gram matrix have the same $V_\tau$ in law. Any
difference in $V_\tau$ between experts of equal coverage is a departure from
the linear model.
\end{corollary}

\begin{proof}
Immediate from Proposition~\ref{prop:response-sum}: equal Gram matrices give
warm starts $(V_1,b_1)$ of the same law, and the LinUCB phase is a
deterministic function of $(V_1,b_1)$ and of the future contexts and noise,
which are independent of the warm start.
\end{proof}

\paragraph{Misspecification opens the second channel.}
Suppose $r(x)=\langle\theta^\star,x\rangle+\varepsilon(x)+\eta$ for a fixed
reference $\theta^\star$---for instance the best linear predictor---and a
residual $\varepsilon$; the binary Precision@1 reward of
Section~\ref{sec:experiments} is of this form for every $\theta^\star$, and
badly so, since no linear function of $\phi(x,a)$ reproduces an indicator.
Regret is still measured against the true best action, with
$\bar r:=\langle\theta^\star,\cdot\rangle+\varepsilon(\cdot)$ and
$\Delta_{t,a}=\bar r(x_{t,a^\star_t})-\bar r(x_{t,a})$.

\begin{proposition}[The transfer]
\label{prop:transfer}
Under the misspecified model a warm start built from $(x_s,r_s)_{s\le\tau}$
satisfies
\[
\widehat\theta_1=\theta^\star-\lambda V_1^{-1}\theta^\star+m+V_1^{-1}\xi,
\qquad
m:=V_1^{-1}\sum_{s\le\tau}\varepsilon(x_s)\,x_s .
\]
Given $M$, the \emph{transfer} $m$ is the only term that depends on which
actions were logged: it is the ridge projection of the residuals the expert's
own picks incur. Under linear rewards $\varepsilon\equiv0$, $m=0$, and
Proposition~\ref{prop:response-sum} is recovered.
\end{proposition}

\begin{proposition}[What the transfer can buy]
\label{prop:transfer-regret}
Under the misspecified model, at a round $t$ at which LinUCB acts and for any
reference $\widetilde\theta\in\R^d$, write
$\widetilde\rho_t:=\|\widehat\theta_t-\widetilde\theta\|_{V_t}$ and
$\widetilde\varepsilon(x):=\bar r(x)-\langle\widetilde\theta,x\rangle$. Then
\[
\Delta_{t,A_t}\ \le\ (\beta_t+\widetilde\rho_t)\|x_{t,A_t}\|_{V_t^{-1}}
\ +\ (\widetilde\rho_t-\beta_t)_+\|x_{t,a^\star_t}\|_{V_t^{-1}}
\ +\ \bigl[\widetilde\varepsilon(x_{t,a^\star_t})-\widetilde\varepsilon(x_{t,A_t})\bigr],
\]
simultaneously for every reference. The design matrix fixes the widths; the
response sum fixes $\widetilde\rho_t$ for every reference at once. The
transfer helps exactly when it moves $\widehat\theta$ toward a reference whose
residual is smaller on the optimal arm than on the arms LinUCB would otherwise
play, and the bound is width-limited whenever some reference has
$\widetilde\rho_t\le\beta_t$ and a nonpositive residual difference on the
played arm.
\end{proposition}

Proposition~\ref{prop:transfer-regret} is a statement about what the second
channel \emph{can} deliver, not a prediction of what a given expert will
deliver: $m$ depends on the residuals at the actions the expert chooses, which
are not known before it is run. Section~\ref{sec:amazon} measures $m$ on real
data, finds that it accounts for essentially the whole estimate at handoff,
and shows that removing it collapses the differences between experts of equal
accuracy; Appendix~\ref{app:diagnostics} reports what this costs us in
predictive power.

\subsection{The Template and Rows 4--10}
\label{sec:template}


\paragraph{The template.}
All instances below are built from the canonical one by three modifications:
directions may recur, a cheap informative action may be available, and that
action may disappear. Fix $K\ge1$ and set $d=K+1$ (or $d=2K+1$ when nuisance
directions are present). The parameter is $\theta^\star=e_1$. Directions
cycle: $j(t):=2+((t-1)\bmod K)$. Three actions may be offered at round $t$:
\[
g:=e_1\ (\mu=1),\qquad
b_t:=Le_{j(t)}\ (\mu=0),\qquad
p_t:=(1-\eps)e_1+c\,e_{j(t)}\ (\mu=1-\eps),
\]
the \emph{good} arm, the \emph{bad} arm, and the \emph{probe}. The instance
$\mathcal I(K,L,c,\eps,\tau_0)$ has $\mathcal A_t=\{g,b_t,p_t\}$ for $t\le\tau_0$
and $\mathcal A_t=\{g,b_t\}$ for $t>\tau_0$. The price of $e_j$, $j\ge2$, is
$\min\{1/L^2,\eps/c^2\}$ while probes are offered and $1/L^2$ afterwards. The
expert is the frozen linear policy $\Exp_w$ that plays
$\arg\max_{a\in\mathcal A_t}\langle\widehat\theta_{\Exp},x_{t,a}\rangle$ with
$\widehat\theta_{\Exp}=e_1+w\sum_{j\ge2}e_j$; ties do not occur for the values
of $w$ we use. Throughout, $\tau_0=Kn$ for an integer $n$, so each direction is
probed exactly $n$ times by a policy that plays every probe.

\paragraph{Standing conventions.}
Noise is i.i.d.\ $\sigma$-sub-Gaussian and independent of the actions; we use
the reward-table model in which the $m$-th pull of a given action vector
receives the $m$-th entry of an i.i.d.\ sequence attached to that vector
\cite[\S4.6]{lattimore2020bandit}, which has the same law as the sequential
model and lets us compare two policies on the same noise. We fix
\begin{equation}
\label{eq:conventions}
L\ge8,\quad c\in(0,1],\quad 0<\eps<c/L,\quad
\lambda\ \ge\ \max\bigl\{\sigma^2d\,\ell_0,\ 8\sigma^2\ell_1,\ 60\bigr\},\quad
\beta\equiv2\sqrt\lambda,
\end{equation}
with $\ell_0:=\log(1+TL^2)+2\log(1/\delta)$ and $\ell_1:=\log(2dT/\delta)$.
The confidence radius is constant. The lower bound on $\lambda$ makes
$2\sqrt\lambda$ at least the radius of \cite{abbasi2011improved} at every
round, so the confidence event of Section~\ref{sec:width} holds for each
policy with probability $\ge1-\delta$, and it makes the noise on any $\lambda/30$
pulls of one arm smaller than a quarter of that arm's confidence bonus. The
constant $m_0:=\lfloor\lambda/30\rfloor\ge2$ is the number of pulls of a
fresh bad arm that optimism forces before the good arm or the probe can be
preferred to it. Conditions on $T$ appear in each statement; they only say
that the horizon is long enough for the claimed rounds to take place.


\paragraph{Row 4: an expert that mimics cold exploration.}
\begin{proposition}[Expert mimics cold exploration: same price, no gain]
\label{prop:cat-mimic}
Let $\mathcal I(K,L,\cdot,\cdot,0)$ be the instance with no probes
($\mathcal A_t=\{g,b_t\}$ for all $t$, $d=K+1$), and let $\Exp$ be any policy
that plays $b_t$ at every round (for instance $\Exp_w$ with $w>1/L$). With
probability at least $1-2\delta$: (i) for every $\tau\le Km_0$,
$\Hyb_\tau$ and $\Cold$ play the same action at every round, so
$A_\tau=V_\tau=0$ and $R_T(\Hyb_\tau)=R_T(\Cold)$; (ii) for every
$\tau>Km_0$, $R_T(\Hyb_\tau)\ge\tau>Km_0$ while $R_T(\Cold)\le K(4\beta^2+1)$,
so $R_T(\Cold)/R_T(\Hyb_\tau)\le(16\lambda+1)/m_0\le961$.
\end{proposition}

\paragraph{Row 5: cheap recurring probes, the positive case.}
On $\mathcal I(K,L,c,\eps,\tau_0)$ the expert $\Exp_w$ with
$w\in(1/L,(1-\eps)/(L-c))$ plays every probe while probes are offered and
every bad arm afterwards (Lemma~\ref{lem:expert-routes}). Its routed actions
cost $\eps$ each, cover $e_{j(t)}$ at rate $c^2$ per round, and recur; after
$\tau_0$ it is the worst possible policy. The hybrid that hands off exactly
when the probes disappear therefore tests every clause of the thesis at once.

\begin{theorem}[Online multiplication, and the price ratio as the multiplier]
\label{thm:price}
Let $\mathcal I=\mathcal I(K,L,c,\eps,Kn)$ under \eqref{eq:conventions}, with
$\Exp=\Exp_w$ for some $w\in(1/L,(1-\eps)/(L-c))$, and $\tau=\tau_0=Kn$. With
probability at least $1-3\delta$ the following hold simultaneously.
\begin{enumerate}[itemsep=1pt,topsep=1pt,leftmargin=*]
\item[(i)] \emph{(Expert.)} $R_T(\Exp)=\eps Kn+(T-Kn)$.
\item[(ii)] \emph{(Cold, lower bound.)} If $n\ge m_0$ then
$R_{[1,Kn]}(\Cold)\ge Km_0$; in particular $A_{\tau_0}\ge Km_0-\eps Kn$.
\item[(iii)] \emph{(Cold, upper bound.)} For every $n\ge1$, $R_T(\Cold)\le K(8\beta^2+1)+\eps Kn$.
\item[(iv)] \emph{(Hybrid.)} If $n\ge36\lambda L^2/c^2$ and $K\ge36L^2$, then
$\Hyb_{\tau_0}$ plays $g$ at every round $t>\tau_0$, so
$R_T(\Hyb_{\tau_0})=\eps Kn$, $V_{\tau_0}=R_{[Kn+1,T]}(\Cold)\ge0$, and
$B_{\tau_0}=T-Kn$.
\end{enumerate}
Consequently, for $n\ge36\lambda L^2/c^2$ and $K\ge36L^2$,
\begin{equation}
\label{eq:price-ratio}
\frac{m_0}{\eps n}\ \le\ \frac{R_T(\Cold)}{R_T(\Hyb_{\tau_0})}\ \le\ 1+\frac{8\beta^2+1}{\eps n},
\qquad
\frac{R_T(\Exp)}{R_T(\Hyb_{\tau_0})}=1+\frac{T-Kn}{\eps Kn},
\end{equation}
and at the smallest admissible $n=\lceil36\lambda L^2/c^2\rceil$ the first
ratio is $\Theta(c^2/(\eps L^2))$, with explicit constants
$c^2/(2220\,\eps L^2)\le R_T(\Cold)/R_T(\Hyb_{\tau_0})\le1+c^2/(\eps L^2)$
(the constants of this section are not optimized; see
Appendix~\ref{app:catalog}, ``Remarks on the constants'').
\end{theorem}

The bracket \eqref{eq:price-ratio} is the price ratio. Cold exploration buys
coverage $L^2$ of a bad direction per unit of regret, so its price is $1/L^2$;
the probe buys $c^2$ per $\eps$, price $\eps/c^2$; the ratio of the two prices
is $c^2/(\eps L^2)$, and the multiplier over cold LinUCB is that ratio up to
constants, and is $O(1)$ when the ratio is. The confidence radius does not
appear: both policies must push the same widths below the same level, and
$\beta$ scales the required coverage identically for both. The multiplier
over the expert is $(T-\tau_0)/(\eps\tau_0)$ and diverges as the tail
lengthens, because the expert keeps paying for bad arms that the warm-started
learner has learned to avoid. The lower bound in (ii) is attained inside the
routed window, so on this instance $A_{\tau_0}>0$ whenever $\eps n<m_0$: the
expert's probes are cheaper than the exploration cold LinUCB is forced into
at the same rounds, and the data they leave is what prevents that exploration
from resuming after $\tau_0$. On this instance the saving is carried by
$A_{\tau_0}$, not by $V_{\tau_0}$, and for a structural reason: the
certificate in (iv) needs coverage of order $\beta^2L^2$ of each bad
direction, which the probe supplies at $c^2$ per round while cold LinUCB's
forced pulls supply it at $L^2$ per round. The window long enough for the
expert to certify ($n\ge36\lambda L^2/c^2$ rounds per direction) is therefore
longer by a factor of order $L^2/c^2$ than the exploration cold LinUCB is
forced into ($m_0\le$ pulls $\le 8\beta^2+1$ per direction), and cold LinUCB
completes that exploration inside the window whenever the bad arm is offered
there (exact-arithmetic simulations of the template give
$R_{[Kn+1,T]}(\Cold)=0$ across the admissible range). Theorem~\ref{thm:general-price}
does not distinguish the two channels: the certified saving is bounded by the
same price ratio wherever it is collected. Row~$5'$ shows that the tail
channel is attainable with the same ratio when the routed window keeps cold
LinUCB away from the directions it will be trapped on later.

\paragraph{Row $5'$: the same price ratio, collected on the tail.}
For the saving to appear as $V_\tau$, the routed window must contain rounds
at which cold LinUCB is \emph{not} offered the bad arms it will later be
trapped on, and is busy with something else; otherwise, as just explained, it
completes its forced exploration inside the window. The following instance
keeps cold LinUCB busy with \emph{distractors}: arms of gap zero whose
feature vectors carry a large component in a fresh direction, so that
optimism prefers them to everything else while the expert probes. Fix
$K,K',n\in\N$, put $d=1+K+K'$, $\theta^\star=e_1$, $L':=12$, and
$\tau_0:=Kn$. Directions $J:=\{2,\dots,K+1\}$ are the bad directions and
$J':=\{K+2,\dots,K+K'+1\}$ the distractor directions, cycled by
$j(t):=2+((t-1)\bmod K)$ and $j'(t):=K+2+((t-1)\bmod K')$. The arms are
\[
g=e_1\ (\mu=1),\quad p_j=(1-\eps)e_1+c\,e_j\ (\mu=1-\eps),\quad
b_j=Le_j\ (\mu=0),\quad b'_{j'}=e_1+L'e_{j'}\ (\mu=1),
\]
and the instance $\mathcal J(K,K',L,c,\eps,n)$ offers
$\mathcal A_t=\{g,p_{j(t)},b'_{j'(t)}\}$ for $t\le\tau_0$ and
$\mathcal A_t=\{g,b_{j(t)}\}$ for $t>\tau_0$. The expert plays the probe while
probes are offered and the bad arm afterwards; the linear policy
$\widehat\theta_{\Exp}=e_1+w\sum_{j\in J}e_j-\sum_{j'\in J'}e_{j'}$ with
$w\in(1/L,(1-\eps)/(L-c))$ does this, by the computation of
Lemma~\ref{lem:expert-routes} together with $\langle\widehat\theta_{\Exp},b'_{j'}\rangle=1-L'<1$.
The prices of $e_j$, $j\in J$, are $\eps/c^2$ in the window and $1/L^2$ on
the tail, as in Theorem~\ref{thm:price}; the distractor directions have
price $0$ and no bad arm uses them.

\begin{proposition}[The price ratio collected as $V_\tau$]
\label{prop:cat-tail-saving}
Let $\mathcal J=\mathcal J(K,K',L,c,\eps,n)$ under \eqref{eq:conventions}
strengthened to $L\ge13$ and $\lambda\ge16\sigma^2d\,\ell_0$, with
$n\ge36\lambda L^2/c^2$, $K\ge36L^2$, $K'\ge L'^2Kn/\lambda$, and
$T-\tau_0\ge Km_0$. With probability at least $1-3\delta$ the following hold
simultaneously.
\begin{enumerate}[itemsep=1pt,topsep=1pt,leftmargin=*]
\item[(i)] \emph{(Expert.)} $R_T(\Exp)=\eps Kn+(T-Kn)$.
\item[(ii)] \emph{(Cold.)} Cold LinUCB plays the distractor at every round
$t\le\tau_0$ and the bad arm at each of the first $m_0$ rounds at which every
tail direction is offered, so $R_{[1,\tau_0]}(\Cold)=0$ and
$Km_0\le R_{[\tau_0+1,T]}(\Cold)\le K(4\beta^2+1)$.
\item[(iii)] \emph{(Hybrid.)} $\Hyb_{\tau_0}$ plays $g$ at every round
$t>\tau_0$, so $R_T(\Hyb_{\tau_0})=\eps Kn$.
\end{enumerate}
Consequently $A_{\tau_0}=-\eps Kn<0$, $Km_0\le V_{\tau_0}\le K(16\lambda+1)$,
$B_{\tau_0}=T-Kn$, and
\[
\frac{m_0}{\eps n}\ \le\ \frac{R_T(\Cold)}{R_T(\Hyb_{\tau_0})}\ \le\ \frac{16\lambda+1}{\eps n},
\]
which at the smallest admissible $n$ lies between $c^2/(2220\,\eps L^2)$ and
$c^2/(2\eps L^2)$: the bracket of Theorem~\ref{thm:price} up to its constant,
now carried entirely by $V_{\tau_0}$.
\end{proposition}

The proof (Appendix~\ref{app:proof-tail-saving}) has one new ingredient: on
the confidence event with the true radius, which the strengthened condition
on $\lambda$ makes at most $\tfrac58\beta$, every index satisfies
$U_t(x)\ge\mu(x)+\tfrac38\beta\|x\|_{V_t^{-1}}$, and the distractor's width is
at least $L'/\sqrt{2\lambda}$ as long as it has been pulled at most
$\lambda/L'^2$ times, which the condition on $K'$ guarantees. Its index then
exceeds $7$ while the indices of $g$ and of the probe are below $4.3$ and
$5.6$. Cold LinUCB never probes, reaches the tail with every bad direction
untouched, and is trapped there exactly as in Theorem~\ref{thm:price}(ii);
the hybrid's window is the same as in Theorem~\ref{thm:price}(iv). In the
language of Theorem~\ref{thm:general-price}, rows~5 and~$5'$ have the same
certified set $U=\{e_j:j\in J\}$, the same prices, and the same $P_U=c^2/(\eps L^2)$;
they differ only in which rounds cold LinUCB's counterfactual exploration
falls on. Data streams with many more unresolved directions than any
affordable window can certify (Section~\ref{sec:amazon}) are of the second
kind.


\paragraph{Row 6: the same expert, optimal on the tail.}
\begin{proposition}[Expert optimal on the tail: large $V_\tau$, no gain over the expert]
\label{prop:cat-tail}
Let $\mathcal I=\mathcal I(K,L,c,\eps,Kn)$ with $n\ge36\lambda L^2/c^2$ and
$K\ge36L^2$, and let $\Exp=\Exp_w$ with $w\in(\eps/c,1/L)$, which plays every
probe and then $g$ at every round. With probability at least $1-3\delta$,
$R_T(\Exp)=R_T(\Hyb_{\tau_0})=\eps Kn$ and $B_{\tau_0}=0$, while $A_{\tau_0}$
and $V_{\tau_0}$ are exactly as in Theorem~\ref{thm:price}.
\end{proposition}

\paragraph{Row 7: probes that cover the wrong directions.}
\begin{proposition}[Probes cover the wrong directions]
\label{prop:cat-wrong}
Let $d=2K+1$ and replace the probe by
$p'_t:=(1-\eps)e_1+c\,e_{K+j(t)}$, whose second coordinate is a nuisance
direction that no bad arm uses, with $\mathcal A_t=\{g,b_t,p'_t\}$ for
$t\le\tau_0=Kn$ and $\{g,b_t\}$ afterwards. Let $\Exp$ play $p'_t$ for
$t\le\tau_0$ and $b_t$ afterwards (for instance the linear policy with
$\widehat\theta_{\Exp}=e_1+w\sum_{j\ge2}e_j$, $w\in(1/L,(1-\eps)/(L-c))$).
Assume $n\ge36\lambda L^2/c^2$, $K\ge36L^2$, and $T-\tau_0\ge K\lfloor\lambda/2\rfloor$.
With probability at least $1-3\delta$,
\[
R_T(\Hyb_{\tau_0})\ \ge\ \eps Kn+K\lfloor\lambda/2\rfloor,
\qquad
R_T(\Cold)\ \le\ K(4\beta^2+1)+\eps Kn,
\qquad
\frac{R_T(\Cold)}{R_T(\Hyb_{\tau_0})}\le65 .
\]
The principal angle between the expert's Gram matrix and the span of the bad
directions is $90^\circ$ and $V_{\tau_0}\le R_{[\tau_0+1,T]}(\Cold)-K\lfloor\lambda/2\rfloor$.
\end{proposition}

\paragraph{Row 8: probes that are too dear.}
\begin{corollary}[Row 8: probes no cheaper than cold exploration]
\label{cor:cat-dear}
In the setting of Theorem~\ref{thm:price}, if $\eps\ge c^2/L^2$ (the probe's
price $\eps/c^2$ is at least cold exploration's price $1/L^2$) then
$R_T(\Cold)/R_T(\Hyb_{\tau_0})\le2$ for every $n\ge36\lambda L^2/c^2$.
\end{corollary}

\paragraph{Row 9: cold LinUCB already cheap.}
\begin{proposition}[Cold LinUCB already cheap: no multiplier in average reward]
\label{prop:cat-cheap}
On the no-probe instance $\mathcal I(K,L,\cdot,\cdot,0)$, with probability
at least $1-\delta$, $R_T(\Cold)\le K(4\beta^2+1)=K(16\lambda+1)$. Hence for
\emph{every} policy $\pi$, $(R_T(\Cold)-R_T(\pi))/T\le K(16\lambda+1)/T$, which with
$\lambda=\sigma^2d\ell_0$ is $O(\sigma^2K^2\ell_0/T)$.
\end{proposition}

\begin{remark}[What certification costs on the template]
On the template, the certificate used in Theorem~\ref{thm:price}(iv), namely
that the confidence bound alone excludes the bad arm in direction $e_j$,
$2\beta L\|e_j\|_{V_{\tau+1}^{-1}}<1$, requires by
Lemma~\ref{lem:no-free-coverage} that
$R_{[1,\tau]}(\Exp)>\pi(e_j)\,(4\beta^2L^2-\lambda)$ for that direction, hence
a total acquisition regret exceeding
$K\min\{1/L^2,\eps/c^2\}(4\beta^2L^2-\lambda)$ for all $K$ directions. This is
Theorem~\ref{thm:general-price}(b) written out for the template.
\end{remark}

\paragraph{Row 10: the gap.}
\begin{proposition}[The gap enters only through the price]
\label{prop:cat-gap}
Fix $\eta\in(0,1]$ and replace $g$ by $g_\eta:=\eta e_1$ ($\mu=\eta$) and
$p_t$ by $p_{t,\eta}:=(1-\eps)\eta e_1+c\,e_{j(t)}$ ($\mu=(1-\eps)\eta$), so that
the gaps of $b_t$ and $p_t$ are $\eta$ and $\eps\eta$ and the prices of
$e_j$ are $\eta/L^2$ and $\eps\eta/c^2$. Let $\Exp=\Exp_w$ with
$w\in(\eta/L,(1-\eps)\eta/(L-c))$. Under \eqref{eq:conventions}, with
$n\ge36\lambda L^2/(\eta^2c^2)$, $K\ge36L^2/\eta^2$, and $\tau=\tau_0=Kn$, with
probability at least $1-3\delta$,
\[
\frac{\eta^2c^2}{2220\,\eps L^2}\ \le\ \frac{\eta Km_0}{\eps\eta Kn}\ \le\
\frac{R_T(\Cold)}{R_T(\Hyb_{\tau_0})}\ \le\ 1+\frac{8\beta^2/\eta^2+1}{\eps n}\ \le\ 1+\frac{c^2}{\eps L^2}.
\]
The upper bound is the price ratio $(\eta/L^2)/(\eps\eta/c^2)=c^2/(\eps L^2)$,
in which $\eta$ cancels; the lower bound loses a factor $\eta^2$ because our
lower bound on cold LinUCB counts only the $m_0$ pulls per direction that
optimism forces regardless of the gap.
\end{proposition}

\paragraph{Where the experiments sit.}
Each configuration of Section~\ref{sec:experiments} has a row. WikiLinks with
the Qwen expert is row~3: the expert is at the linear learner's ceiling, and
the hybrid's gain is additive. WikiLinks with Flan-T5 is row~9: the routed
data does reach useful directions ($V_\tau>0$ on every seed), but cold LinUCB
already resolves most of what matters and the multiplier is $1.01\times$; its
principal angle of $82^\circ$ alone would have suggested row~7, which the
decomposition rules out (Section~\ref{sec:wiki}). AmazonCat with Flan-T5 is
row~$5'$ for all three experts---multipliers $1.71\times$, $2.72\times$ and
$3.29\times$, routed rounds concentrated in the first few thousand, and $96$--$100\%$ of
the saving accruing after the last routed round
(Figure~\ref{fig:amazon-dynamics}), which is the tail channel. AmazonCat with
BGE features lies between the two: cold LinUCB already reaches about a third of
the correct answers, and once the gate is given a $c$ at which it engages it
adds a smaller tail gain (up to $1.23\times$ on held-out seeds). Section~\ref{sec:diagnostics} turns the row assignment into
quantities that can be estimated before any online round is played.



\subsection{What Governs the Regret of LinUCB at a Single Round}
\label{sec:width}

Fix the confidence event $\mathcal E_\delta$ on which
$\|\widehat\theta_t-\theta^\star\|_{V_t}\le\beta_t$ for all rounds $t$ at
which LinUCB acts, where $(V_t,b_t)$ are built from all data so far including
any warm start $(M,c)$ formed from the same kind of data; by the
self-normalized bound of \cite{abbasi2011improved} this event has probability
at least $1-\delta$ with
$\beta_t=\sigma\sqrt{2\log(\det(V_t)^{1/2}\det(\lambda I)^{-1/2}/\delta)}+\sqrt\lambda S$.
Define, for a positive definite $V$,
\begin{equation}
\label{eq:phi}
\Phi_t(V)\;:=\;\max_{a\in\mathcal A_t}\;\min\Bigl\{\Delta_{t,a},\;2\beta_t\|x_{t,a}\|_{V^{-1}}\Bigr\}.
\end{equation}

\begin{lemma}[Width-limited regret]
\label{lem:width-limited}
On $\mathcal E_\delta$, at any round $t$ at which LinUCB acts with design
matrix $V_t$, the instantaneous pseudo-regret satisfies
$\Delta_{t,A_t}\le\Phi_t(V_t)$. Moreover $\Phi_t$ is monotone: $V\preceq V'$
implies $\Phi_t(V')\le\Phi_t(V)$.
\end{lemma}

The lemma is the familiar fact that an optimistic algorithm can only play an
arm whose width exceeds half its gap, restated so that the gap and the width
sit in one expression. It makes three facts visible at once. A competitor with
a large gap but an even larger width is exactly as dangerous as its gap
(Proposition~\ref{prop:three-way}: $\Delta=1$, width $\beta L/\sqrt\lambda>1/2$).
A competitor with a tiny gap is cheap no matter how wide it is, so a heavy
mass of near-ties does not by itself mean large regret. And the only way to
reduce the regret LinUCB can incur on a given round is to reduce widths
\emph{on competitors whose widths exceed their gaps}; reducing the width of
$g$ in Proposition~\ref{prop:three-way} would have achieved nothing.

There are exactly two ways widths get reduced: through data that arrives as a
warm start, and through LinUCB's own observations. Both are exact linear
algebra.

\begin{lemma}[A warm start shrinks widths in the directions it covers]
\label{lem:loewner}
For any $M\succeq0$ and any $x$, $\|x\|_{(\lambda I+M)^{-1}}\le\|x\|_{(\lambda I)^{-1}}=\|x\|_2/\sqrt\lambda$.
If $P$ is an orthogonal projector, $M\succeq mP$, and $x\in\mathrm{Im}(P)$, then
$\|x\|^2_{(\lambda I+M)^{-1}}\le\|x\|_2^2/(\lambda+m)$. If instead $x\perp\mathrm{Im}(M)$,
then $\|x\|_{(\lambda I+M)^{-1}}=\|x\|_2/\sqrt\lambda$: a warm start does nothing
for directions it does not cover.
\end{lemma}

\begin{lemma}[One observation shrinks the width of another direction by an
exact factor]
\label{lem:transfer}
For $V\succ0$ and vectors $x,z$ with $z\ne0$,
\[
\|z\|^2_{(V+xx^\top)^{-1}}=\bigl(1-\rho_V(x,z)\bigr)\|z\|^2_{V^{-1}},
\qquad
\rho_V(x,z):=\frac{(z^\top V^{-1}x)^2}{(1+\|x\|^2_{V^{-1}})\,\|z\|^2_{V^{-1}}}\in[0,1).
\]
At a cold start, $\rho_{\lambda I}(x,z)=\cos^2(x,z)\cdot\|x\|_2^2/(\lambda+\|x\|_2^2)$.
\end{lemma}

Lemma~\ref{lem:transfer} is where ``feature coherence'' enters, and it shows
precisely how much of that notion survives. Observing $x$ reduces the width of
$z$ by the squared cosine \emph{in the $V^{-1}$ geometry}, scaled by how novel
$x$ itself is. The Euclidean cosine is the right quantity only at $V=\lambda I$;
as $V$ becomes anisotropic the relevant angle is measured in $V^{-1}$, and two
vectors with positive Euclidean inner product can have $\rho_V$ near zero (or
vice versa). In Proposition~\ref{prop:three-way} the Euclidean cosine between
any two bad arms is exactly zero, so $\rho=0$ and no observation ever helps a
future round; that is what makes the instance hard, not the ambient dimension
as such. A large ambient dimension is necessary for many mutually
$V^{-1}$-orthogonal directions to exist, but it is the lack of transfer, not
the dimension, that cold LinUCB pays for.

\paragraph{How the earlier structural conditions arise.}
Conditions sometimes presented as independent structural axes
(effective dimension, feature coherence, gap structure) are all readings of
Lemma~\ref{lem:width-limited} together with Lemmas~\ref{lem:loewner}
and~\ref{lem:transfer}: a uniform gap $\Delta_{\min}$ caps $\Phi_t$ at
$\Delta_{\max}$ but saves nothing while widths exceed $\Delta_{\min}/2$; a small
rank $r$ of the span of played directions caps the elliptical potential
$\sum_t\|x_{t,A_t}\|^2_{V_t^{-1}}$ at $2r\log(1+L^2T/(r\lambda))$ and therefore
the number of rounds at which any width can exceed a fixed level; and positive
coherence lower-bounds $\rho_{\lambda I}$ at the start of learning only. None
of the three is a condition under which a warm start cannot help by more than
a constant factor, and we do not claim one here; what they control is how
many rounds of cold exploration are needed before $\Phi_t$ falls below the
level the warm start could have delivered immediately.


\subsection{Regret of the Handoff Hybrid}
\label{sec:handoff}

\begin{proposition}[Handoff regret in exact form]
\label{thm:handoff}
Fix $\tau$ and an expert $\Exp$. On the event $\mathcal E_\delta$ (probability
at least $1-\delta$), the handoff hybrid satisfies
\begin{equation}
\label{eq:handoff-bound}
R_T(\Hyb_\tau)\;\le\;
\underbrace{\sum_{t\le\tau}\Delta_{t,a_t^{\Exp}}}_{\text{acquisition regret}}
\;+\;
\min\Biggl\{\sum_{t=\tau+1}^{T}\Phi_t(V_t),\;\;
2\beta_T\max\Bigl\{1,\tfrac{L}{\sqrt\lambda}\Bigr\}\sqrt{2(T-\tau)\,\log\frac{\det V_{T+1}}{\det V_{\tau+1}}}\Biggr\},
\end{equation}
where $V_{\tau+1}=\lambda I+M_\tau$ and the widths inside $\Phi_t$ are those
of the hybrid's own trajectory. The determinant ratio satisfies
$\log\det V_{T+1}/\det V_{\tau+1}\le\sum_{j=1}^{d}\log\bigl(1+L^2(T-\tau)/\sigma_j(V_{\tau+1})\bigr)$
with $\sigma_j$ the eigenvalues of $V_{\tau+1}$, and cold LinUCB satisfies the
same bound with $\tau=0$ and $V_1=\lambda I$.
\end{proposition}

We state the bound with the exact determinant ratio and make no asymptotic
simplification of it. Two remarks explain why.

\emph{The determinant term cannot certify a multiplier.} Suppose
$M_\tau\succeq\nu\tau P$ on an $r$-dimensional subspace containing all actions,
the situation analyzed in earlier work. Then the second term in
\eqref{eq:handoff-bound} is at most
$2\beta_T\sqrt{2r(T-\tau)\log(1+L^2(T-\tau)/(\lambda+\nu\tau))}$, and its ratio
to the cold bound $2\beta_T\sqrt{2rT\log(1+L^2T/\lambda)}$ is
\[
\Psi(\tau,\nu,T)=\sqrt{\frac{T-\tau}{T}}\cdot\sqrt{\frac{\log\bigl(1+L^2(T-\tau)/(\lambda+\nu\tau)\bigr)}{\log\bigl(1+L^2T/\lambda\bigr)}}.
\]
For $\tau=o(T)$ both logarithms are $\Theta(\log T)$ and $\Psi\to1$; for
$\tau=T^{\alpha}$ with $\nu\tau\gg\lambda$ it tends to $\sqrt{1-\alpha}$, bounded
away from zero. It is tempting to read $\Psi$ as
$\sqrt{\lambda/(\lambda+\nu\tau)}$ by comparing the arguments of the two
logarithms; that reading is wrong. The correct conclusion is the opposite
one: a warm start that improves only the $\sqrt{T}$-rate term buys at most a
constant factor, and any multiplier must come from the first term in the
minimum, that is, from competitors whose widths the warm start pushes below
their gaps so that $\Phi_t$ drops from the order of the gap to the order of
the (now small) width or to zero. That is what happens in
Proposition~\ref{prop:three-way}, where $\sum_t\Phi_t$ falls from $T$ to $O(1)$
while the determinant bound barely moves.

\emph{What the bound says about multiplication.} Combining
\eqref{eq:handoff-bound} with Proposition~\ref{prop:decomp}, the hybrid beats
cold LinUCB by more than a constant factor only if
\[
\sum_{t>\tau}\bigl(\Delta_{t,A^{\Cold}_t}-\Phi_t(V^{\Hyb}_t)\bigr)
\;\gg\;
\sum_{t\le\tau}\Delta_{t,a_t^{\Exp}}-R_{[1,\tau]}(\Cold),
\]
i.e., the width reduction that the routed data delivers on the competitors
cold LinUCB would have played must exceed the net regret paid to acquire it;
and it beats the expert only if the expert's tail regret exceeds the hybrid's.
Neither condition is implied by low expert accuracy being ``near-optimal
early'' or by an abstract coverage parameter. Both are properties of
\emph{which} actions the expert proposes, relative to the directions cold
exploration leaves unresolved, and of what those actions cost. This is the
quantity our diagnostics estimate.

\emph{Adaptive $\tau$ and interleaving.} The router of
Appendix~\ref{app:meta-algorithm} decides at each round whether to follow the
expert or LinUCB by comparing confidence widths, so the expert rounds need not
be contiguous. Proposition~\ref{thm:handoff} applies verbatim with
$\sum_{t\le\tau}$ replaced by the sum over expert-routed rounds and with the
LinUCB-phase terms summed over LinUCB-routed rounds, since
Lemma~\ref{lem:width-limited} is a round-level statement and the elliptical
potential argument only uses that $V_t$ is non-decreasing.


\subsection{Pre-Runtime Estimates}
\label{sec:diagnostics}

Proposition~\ref{thm:handoff} and Proposition~\ref{prop:decomp} point to three
quantities a practitioner would want before committing to LLM inference on a
long trajectory, all estimable from a labeled held-out sample
$\mathcal D=\{(\mathcal A_s,\{x_{s,a}\},\{r_{s,a}\})\}_{s=1}^{n}$ and the
expert's proposals $a_s^{\Exp}$ on it (formal definitions and consistency
statements in Appendix~\ref{app:diagnostics}):
\begin{enumerate}[itemsep=1pt,topsep=1pt,leftmargin=*]
\item \textbf{Width-limited regret under cold and warm designs.} Fit a ridge
estimate $\widehat\theta_{\mathrm{or}}$ on $\mathcal D$ to obtain gap estimates
$\widehat\Delta_{s,a}$, and form the expert's held-out Gram matrix
$\widehat M_{\Exp}=\sum_sx_{s,a_s^{\Exp}}x_{s,a_s^{\Exp}}^\top$. For a target
number $\tau$ of routed rounds, compute
$\widehat\Phi(\lambda I)$ and $\widehat\Phi(\lambda I+(\tau/n)\widehat M_{\Exp})$,
each the average over $s$ of
$\max_a\min\{\widehat\Delta_{s,a},2\beta\|x_{s,a}\|_{V^{-1}}\}$. Their difference
is the width reduction the expert's data would deliver on the competitors
that matter, before any online round is played.
\item \textbf{Acquisition regret.} $\widehat C_{\Exp}(\tau)=\tau\cdot\frac1n\sum_s\widehat\Delta_{s,a_s^{\Exp}}$,
which for binary reward is $\tau$ times the expert's held-out error rate.
\item \textbf{Tail advantage.} The expert's held-out accuracy compared with
the accuracy of $\widehat\theta_{\mathrm{or}}$, the latter being an upper
proxy for what the linear learner can reach; multiplication over the expert
is impossible if the expert is already at that ceiling.
\end{enumerate}
The prediction is that a handoff can multiply through the tail channel
($V_\tau$) only if the width reduction in item~1, scaled by the number of
future rounds on which it applies, exceeds the acquisition regret in item~2
by a wide margin, and item~3 leaves room; the concurrent channel ($A_\tau$)
is covered by the dry-run variant of Appendix~\ref{app:diagnostics}.
These are estimates, not certificates; the thresholds we use are stated with
the experiments and are provisional. Four cheaper statistics---effective
rank, median minimum cosine, gap-distribution tail, and the principal angle
between the expert's Gram matrix and the context covariance---are useful as
proxies and we report them alongside:
the rank bounds how many distinct unresolved directions can exist, the cosine
is $\rho_{\lambda I}$ up to the norm factor of Lemma~\ref{lem:transfer}, the
gap tail is the first argument of the minimum in \eqref{eq:phi}, and the
principal angle says whether $\widehat M_{\Exp}$ touches the directions that
dominate the online context covariance at all.


\section{Experiments}
\label{sec:experiments}

We measure the decomposition of Proposition~\ref{prop:decomp} and the two
channels of Section~\ref{sec:two-channels} on real data, running the router,
the LLM alone, and cold LinUCB on the same data stream, and ask how far the
quantities of Section~\ref{sec:diagnostics} anticipate the outcome. We sweep
the router's confidence parameter $c$ over several orders of magnitude so that
negative results are not hyperparameter artifacts. The setup mirrors
\cite{chen2024llmdecision}.

\textbf{Two scales for one effect.} The literature reports \emph{reward
multiplication}, the ratio of a hybrid's average reward to that of a base
policy \cite{chen2024llmdecision}; with binary Precision@1 this is a ratio of
accuracies, and we call it the \emph{accuracy multiplier}. The theorems bound
a different quantity, the ratio $R_T(\Cold)/R_T(\pi)$ of pseudo-regrets. Every
context has at least one correct action, so the best achievable reward is $1$
per round and $R_T(\pi)=T(1-\mathrm{acc}(\pi))$; alongside the accuracy
multiplier we report the \emph{regret ratio}
$(1-\max\{\mathrm{acc}_{\mathrm{LLM}},\mathrm{acc}_{\mathrm{Cold}}\})/(1-\mathrm{acc}_{\mathrm{Router}})$.
The two agree when accuracies are high and differ sharply when they are low:
on AmazonCat an accuracy multiplier of $3.29$ from a base of $8.4\%$ is a
regret ratio of $1.27$, that is $21\%$ less regret. Statements about the
theorems refer to the regret ratio. $A_\tau$, $V_\tau$ and $B_\tau$ are
differences of regret, counted in correct answers, and are the same on either
scale.

\subsection{Setup}
\label{sec:exp-setup}

\textbf{Datasets.} \textbf{OneShotWikiLinks-311}\footnote{\url{https://www.kaggle.com/generall/oneshotwikilinks}} ($T\!=\!622{,}000$,
$|\mathcal A|\!=\!311$) pairs two model families (Flan-T5, Qwen2.5) with the
same dataset, giving six configurations. \textbf{AmazonCat-13K}
\cite{bhatia2016extreme} ($T\!=\!50{,}000$, $|\mathcal A|\!=\!13{,}330$) pairs two embedding spaces
(BERT-768d, BGE-1024d) with the same categories. Reward is binary
Precision@1.

\textbf{Policies.} Each experiment tracks three policies on the same stream:
(1)~\textbf{Router} (Algorithm~\ref{alg:confidence-gated}), the deterministic
confidence-gated hybrid with expert design matrix $V^{(1)}=cI_d$, so the
expert's width is $u_t^{(1)}=\|\phi(x_t,a_t^{(1)})\|_2/\sqrt c$ and LinUCB's is
$u_t^{(0)}=\|\phi(x_t,a_t^{(0)})\|_{V_t^{-1}}$; the Router updates $(V_t,b_t)$ on
every round with the played action and observed reward, whichever policy
chose it, so its LinUCB component is exactly the online-warm-started learner
of Section~\ref{sec:preliminaries}; (2)~\textbf{LLM only}; (3)~\textbf{Cold
LinUCB} from $\lambda I_d$. Hyperparameters: $\lambda=0.1$; the exploration
scale is annealed per round as $\alpha_t=\max(0.1,0.9985^t)$ on both datasets,
reaching its floor at $t\approx1.5\times10^3$. A seed fixes the order of the
context stream and is shared by all policies in a run, so every comparison
between policies or experts is paired on seed. Reported results use
seeds $1$--$10$. The confidence parameter $c$ of every configuration is chosen
on the disjoint block of seeds $11$--$20$, by mean Router accuracy over the
three experts, and the chosen value is then evaluated on seeds $1$--$10$, so
that tuning and evaluation never share a seed. On AmazonCat the tuning block
selects $c=100$ for BERT features (mean Router accuracy $17.2$, $19.5$,
$22.2$, $21.4$ and $5.4\%$ at $c=1,10,10^2,10^3,10^4$), the value used
throughout, and $c=0.1$ for BGE features (Table~\ref{tab:bge-csweep}).
\todo{WikiLinks sweep} On WikiLinks the Flan-T5 grid is
$c\in\{1,10,10^2,10^3,10^4\}$ (Table~\ref{tab:wiki-csweep}) and the Qwen grid
$c\in\{0.01,0.1,1,10,100\}$ (Table~\ref{tab:bge-csweep}).
The LLM expert proposes the action whose embedding is closest to the LLM's
textual output, as in \cite{chen2024llmdecision}.

\textbf{Pre-runtime statistics} are computed on a held-out sample of
$5{,}000$ labeled contexts. Table~\ref{tab:diagnostics} reports the four
proxies of Appendix~\ref{app:proxies}, the held-out LLM accuracy (the
acquisition-cost term of Section~\ref{sec:diagnostics} up to the factor
$\tau$), and, for the configurations where we measured it, the price ratio
$P_U$ of Theorem~\ref{thm:general-price}. The proxies were designed for an
earlier version of this work and are retained for comparison only: on the
three AmazonCat Flan-T5 rows they agree to within $1.0^\circ$ of subspace
angle and are identical in the other three columns, while the observed
multipliers span $1.71\times$ to $3.29\times$. The price ratio is different
in kind, and is \emph{not} a pre-runtime quantity. We measure it after the run
from per-round logs, taking for $U$ the leading directions of the features
cold LinUCB played on rounds it lost after routing stopped, and the real
streams lie outside the fixed-arm setting of Theorem~\ref{thm:general-price};
we report it as a post-hoc explanatory statistic. It bounds the regret ratio,
and every observed regret ratio lies below it (Section~\ref{sec:amazon}); it is
equal across the three experts, as the coverage channel is; and it is by
construction blind to the transfer, which Section~\ref{sec:two-channels} shows
is what distinguishes these experts. Appendix~\ref{app:diagnostics} reports
the two estimators we built for the tail channel, $\widehat N$ and the dry run,
and how far each gets.


\begin{table}[h]
\centering
\caption{Pre-runtime statistics, with the observed multiplier for reference
(Flan-T5 rows over ten seeds; AmazonCat Qwen rows over seeds $1$--$10$, held
out from the choice of $c$; WikiLinks Qwen rows pending the sweep). The four proxies are inherited from an
earlier version of this work and are retained only for comparison;
Appendix~\ref{app:proxies} states what each does and does not bound. On the
three AmazonCat Flan-T5 rows they are identical in three columns and span
$1.0^\circ$ of subspace angle while the multipliers span $1.71\times$ to
$3.29\times$: they cannot separate the configurations. $P_U$ is the price
ratio of Theorem~\ref{thm:general-price}; unlike the other columns it is
measured \emph{after} the run, from per-round logs, and is a post-hoc
statistic. It bounds the regret ratio and is blind to the transfer, which is
what orders these three experts. The last column gives the accuracy multiplier
and, after the slash, the regret ratio. Boldface marks values in the direction
that leaves room for a warm start. LLM acc.\ is held-out top-1 accuracy; its
complement, times $\tau$, is the acquisition regret. }
\label{tab:diagnostics}
\scriptsize
\setlength{\tabcolsep}{3pt}
\begin{tabular}{@{}llccccccl@{}}
\toprule
Dataset & Config & $\widehat r\log T/\sqrt T$ & $\widehat{\bar\alpha}$ & Tail $\le\eta$ & Subspace $\angle$ & LLM acc.\ & $P_U$ & Observed (acc.\ / regret) \\
\midrule
WikiLinks & Flan-T5-small  & $0.05$ & $0.19$ & \textbf{heavy} & $\mathbf{82.2^\circ}$ & $7.2\%$ & --- & $1.01\times$ / $1.00$ \\
WikiLinks & Flan-T5-base   & $0.05$ & $0.19$ & \textbf{heavy} & $\mathbf{83.2^\circ}$ & $10.6\%$ & --- & $1.01\times$ / $1.00$ \\
WikiLinks & Flan-T5-large  & $0.05$ & $0.19$ & \textbf{heavy} & $\mathbf{83.0^\circ}$ & $12.4\%$ & --- & $0.88\times$ / $0.98$ \\
WikiLinks & Qwen-0.5B      & $0.04$ & $0.43$ & \textbf{heavy} & $\mathbf{81.9^\circ}$ & $18.7\%$ & --- & --- \\
WikiLinks & Qwen-3B        & $0.04$ & $0.43$ & \textbf{heavy} & $\mathbf{81.9^\circ}$ & $29.5\%$ & --- & --- \\
WikiLinks & Qwen-7B        & $0.04$ & $0.43$ & \textbf{heavy} & $\mathbf{81.9^\circ}$ & $38.7\%$ & --- & --- \\
\midrule
AmazonCat & Flan-T5-small (768d) & $\mathbf{0.16}$ & $\mathbf{0.04}$ & \textbf{heavy} & $\mathbf{85.1^\circ}$ & $5.2\%$ & $1.73$ & $1.71\times$ / $1.07$ \\
AmazonCat & Flan-T5-base (768d)  & $\mathbf{0.16}$ & $\mathbf{0.04}$ & \textbf{heavy} & $\mathbf{84.1^\circ}$ & $3.6\%$ & $1.59$ & $2.72\times$ / $1.19$ \\
AmazonCat & Flan-T5-large (768d) & $\mathbf{0.16}$ & $\mathbf{0.04}$ & \textbf{heavy} & $\mathbf{84.2^\circ}$ & $5.3\%$ & $1.58$ & $\mathbf{3.29\times}$ / $\mathbf{1.27}$ \\
%% TODO(check): the 83.1 angle below was reported once for the BGE/Qwen pair;
%% recompute it for each Qwen size as for the Flan-T5 rows.
AmazonCat & Qwen-0.5B (1024d) & $\mathbf{0.15}$ & $0.31$ & \textbf{heavy} & $\mathbf{83.1^\circ}$ & $21.4\%$ & --- & $1.00\times$ / $1.00$ \\
AmazonCat & Qwen-3B (1024d)   & $\mathbf{0.15}$ & $0.31$ & \textbf{heavy} & $\mathbf{83.1^\circ}$ & $22.4\%$ & --- & $\mathbf{1.23\times}$ / $\mathbf{1.13}$ \\
AmazonCat & Qwen-7B (1024d)   & $\mathbf{0.15}$ & $0.31$ & \textbf{heavy} & $\mathbf{83.1^\circ}$ & $27.2\%$ & --- & $1.19\times$ / $1.10$ \\
\bottomrule
\end{tabular}
\end{table}

\subsection{OneShotWikiLinks-311}
\label{sec:wiki}

\textbf{Setup.} $622{,}000$ context--entity pairs over $311$ entities.
\emph{Flan-T5 (BERT-768d features)}: embeddings from
\texttt{bert-base-nli-mean-tokens}, joint feature $\phi(x,a)=x\odot a\in\mathbb R^{768}$,
expert Flan-T5 \cite{chung2024scaling} at 60M/250M/780M.
\emph{Qwen2.5 (BGE-1024d features)}: embeddings from
\texttt{BAAI/bge-large-en-v1.5} \cite{xiao2024c}, expert Qwen2.5-Instruct
\cite{hui2024qwen2} at 0.5B/3B/7B.

\begin{table}[h]
\centering
\caption{Final cumulative accuracy on OneShotWikiLinks-311, mean$\pm$std over
seeds $1$--$10$. Flan-T5 rows are at $c=100$; the Qwen rows are filled at the
$c$ chosen on seeds $11$--$20$. LLM accuracy and cold LinUCB do not depend on
$c$. Mult.\ is the accuracy multiplier Router/$\max$(LLM, Cold); the regret
ratio is within $0.02$ of one on every Flan-T5 row. The lift column is the
Router's gain over cold LinUCB in percentage points. \todo{WikiLinks sweep}}
\label{tab:wiki-results}
\small
\begin{tabular}{@{}llccccc@{}}
\toprule
Family & Size & LLM & Cold LinUCB & Router & Lift (pp) & Mult.\ \\
\midrule
\multirow{3}{*}{Flan-T5} & small (60M) & $7.16\%$  & $10.71\pm 0.04\%$ & $10.84\pm 0.04\%$ & $+0.13$ & $1.01\times$ \\
& base (250M)   & $10.56\%$ & $10.71\pm 0.04\%$ & $10.87\pm 0.05\%$ & $+0.16$ & $1.01\times$ \\
& large (780M)  & $12.43\%$ & $10.71\pm 0.04\%$ & $10.88\pm 0.04\%$ & $+0.17$ & $0.88\times$ \\
\midrule
\multirow{3}{*}{Qwen2.5} & 0.5B  & $18.72\%$ & $18.88\pm 0.04\%$ & --- & --- & --- \\
& 3B  & $29.47\%$ & $18.88\pm 0.04\%$ & --- & --- & --- \\
& 7B  & $38.74\%$ & $18.88\pm 0.04\%$ & --- & --- & --- \\
\bottomrule
\end{tabular}
\end{table}

\textbf{Results.} Table~\ref{tab:wiki-results} reports final cumulative
accuracy at $c=100$. The Router routes to the LLM on $<0.3\%$ of rounds for
Flan-T5. Routing is driven by feature geometry, not by
reward: all three Flan-T5 sizes produce nearly identical routing curves
despite an order-of-magnitude spread in accuracy.

\textbf{Confidence sweep.} To rule out a suboptimal $c$, the Flan-T5
configurations are swept over $c\in\{1,10,10^2,10^3,10^4\}$ on the tuning seeds
$11$--$20$ (Table~\ref{tab:wiki-csweep}), and the Qwen configurations over
$c\in\{0.01,0.1,1,10,100\}$ (Table~\ref{tab:bge-csweep}).

\begin{table}[h]
\centering
\caption{Confidence sweep on the WikiLinks Flan-T5 (BERT-768d) configurations,
tuning seeds $11$--$20$: share of rounds routed to the LLM and final Router
accuracy at each $c$. Cold LinUCB does not depend on $c$. \todo{WikiLinks sweep}}
\label{tab:wiki-csweep}
\scriptsize
\setlength{\tabcolsep}{4pt}
\begin{tabular}{@{}lccccccc@{}}
\toprule
& \multicolumn{3}{c}{routed \%} & \multicolumn{3}{c}{Router accuracy} & \\
\cmidrule(lr){2-4}\cmidrule(lr){5-7}
$c$ & small & base & large & small & base & large & best mult.\ \\
\midrule
$1$ & --- & --- & --- & --- & --- & --- & --- \\
$10$ & --- & --- & --- & --- & --- & --- & --- \\
$10^2$ & --- & --- & --- & --- & --- & --- & --- \\
$10^3$ & --- & --- & --- & --- & --- & --- & --- \\
$10^4$ & --- & --- & --- & --- & --- & --- & --- \\
\bottomrule
\end{tabular}
\end{table}

\todo{WikiLinks sweep} The decomposition explains the shape. For Flan-T5 the routed rounds do leave
usable data---$V_\tau$ is $748$, $832$ and $886$ correct answers and is
positive on all ten seeds---but cold LinUCB reaches $10.71\%$ on its own, so
there is little left to buy: this is row~9 of Table~\ref{tab:catalog}, not the
wrong-directions failure of row~7. Nor does WikiLinks have a dominant entity
for a warm start to steer the learner toward, as AmazonCat does
(Section~\ref{sec:amazon}): the most frequent of the $311$ entities covers
$0.36\%$ of contexts. For Qwen the expert is far better than the
learner ($29$--$39\%$ against $19\%$), so $B_\tau<0$ for every $\tau$ and no
handoff can beat the expert. \todo{WikiLinks sweep}

\subsection{AmazonCat-13K}
\label{sec:amazon}

AmazonCat-13K has $13{,}330$ categories and $T=50{,}000$ rounds. We run two
feature spaces: \emph{Flan-T5 (BERT-768d)} and \emph{Qwen2.5 (BGE-1024d)};
hyperparameters as in Section~\ref{sec:exp-setup}.

\begin{table}[h]
\centering
\caption{AmazonCat-13K, $T=50{,}000$, mean $\pm$ standard deviation over ten
seeds. The accuracy multiplier (Acc.\ mult.) is the
ratio of the seed-averaged Router accuracy to the seed-averaged $\max\{$LLM,
cold$\}$, which is less sensitive to a single unlucky denominator than the
average of per-seed ratios; the regret ratio is the corresponding ratio of
regrets, $(1-\max\{\text{LLM},\text{cold}\})/(1-\text{Router})$, the quantity
Theorem~\ref{thm:general-price} bounds. Cold LinUCB does not depend on the expert, so its column is
shared within each feature space. BERT rows are at $c=100$, seeds $1$--$10$.
BGE rows are at $c=0.1$, selected on the disjoint seed block $11$--$20$ in
Table~\ref{tab:bge-csweep} and evaluated here on seeds $1$--$10$---at the
$c=100$ used for BERT the gate never engages and all six BGE rows are
degenerate. On BERT the expert's held-out accuracy orders the experts
differently from the multiplier; on BGE the two larger experts gain
significantly (paired $t=3.92$ and $2.59$ against cold) and the smallest does
not move ($t=0.01$).}
\label{tab:amazon-results}
\small
\begin{tabular}{@{}llcccccc@{}}
\toprule
Features & LLM & LLM acc.\ & Cold LinUCB & Router & LLM\% & Acc.\ mult.\ & Regret ratio \\
\midrule
\multirow{3}{*}{BERT-768d} & Flan-T5-small  & $5.20\%$ & $8.39\pm2.46\%$ & $14.36\pm6.12\%$ & $2.8\%$  & $1.71\times$ & $1.07$ \\
& Flan-T5-base  & $3.61\%$ & $8.39\pm2.46\%$ & $22.80\pm9.00\%$ & $3.1\%$  & $2.72\times$ & $1.19$ \\
& Flan-T5-large & $5.27\%$ & $8.39\pm2.46\%$ & $\mathbf{27.64\pm5.56\%}$ & $3.2\%$  & $\mathbf{3.29\times}$ & $\mathbf{1.27}$ \\
\midrule
\multirow{3}{*}{BGE-1024d} & Qwen-0.5B  & $21.38\%$ & $32.59\pm5.00\%$ & $32.62\pm6.33\%$ & $0.04\%$  & $1.00\times$ & $1.00$ \\
& Qwen-3B  & $22.39\%$ & $32.59\pm5.00\%$ & $\mathbf{40.20\pm4.45\%}$ & $0.05\%$  & $\mathbf{1.23\times}$ & $\mathbf{1.13}$ \\
& Qwen-7B  & $27.17\%$ & $32.59\pm5.00\%$ & $38.71\pm6.38\%$ & $0.05\%$  & $1.19\times$ & $1.10$ \\
\bottomrule
\end{tabular}
\end{table}

\begin{table}[h]
\centering
\caption{The decomposition of Proposition~\ref{prop:decomp} measured on the
BERT-768d runs, in correct answers, mean over ten seeds. $A_\tau$ is summed
over the routed rounds and $V_\tau,B_\tau$ over the rest, using the routed set
the Router actually selected (the identities hold for any routed set, by the
remark after Proposition~\ref{prop:decomp}). Every configuration collects
almost all of its advantage through $V_\tau$: the expert's own answers are
worth a few dozen correct rounds out of $1{,}578$ routed, while the data they
leave behind is worth thousands. $V_\tau$ is the only column that separates
the three experts, and the multiplying configuration is also the only one whose
$V_\tau$ is positive on every seed.}
\label{tab:amazon-decomp}
\small
\begin{tabular}{@{}lcccc@{}}
\toprule
LLM & $A_\tau$ & $V_\tau$ & $B_\tau$ & seeds with $V_\tau>0$ \\
\midrule
Flan-T5-small  & $47$ & $2{,}936\pm3{,}562$ & $4{,}582$ & $9/10$ \\
Flan-T5-base   & $27$ & $7{,}178\pm4{,}772$ & $9{,}599$ & $9/10$ \\
Flan-T5-large  & $53$ & $\mathbf{9{,}571\pm3{,}411}$ & $11{,}189$ & $\mathbf{10/10}$ \\
\bottomrule
\end{tabular}
\end{table}

\begin{table}[h]
\centering
\caption{Isolating the two channels on AmazonCat/BERT-768d, ten seeds. In
\emph{full} a routed round contributes both the covariance and the response
sum, as the Router does. In \emph{transfer removed} the round contributes the
covariance and the expert-independent part $M\theta_{\mathrm{or}}$ of the
response sum---the learner is fed
$\langle\theta_{\mathrm{or}},\phi\rangle$ instead of the observed reward on
routed rounds only, which deletes $m$ of Proposition~\ref{prop:transfer} and
changes nothing else. In \emph{no response sum} the round contributes the
covariance alone. Multipliers are accuracy multipliers, Router over cold
LinUCB, with the regret ratio in parentheses; $t$-statistics are paired over
seeds against small, on each scale, $t^\star=2.26$. With the transfer present the
experts separate; without it, by either mechanism, no pairwise difference
survives.}
\label{tab:two-channel}
\small
\begin{tabular}{@{}lccccc@{}}
\toprule
& \multicolumn{2}{c}{Router accuracy} & \multicolumn{3}{c}{Multiplier over cold} \\
\cmidrule(lr){2-3}\cmidrule(lr){4-6}
LLM & full & transfer removed & full & transfer removed & no response sum \\
\midrule
Flan-T5-small & $14.36\pm6.12\%$ & $9.71\pm2.08\%$  & $1.71\times$ ($1.07$) & $1.16\times$ ($1.02$) & $1.53\times$ ($1.05$) \\
Flan-T5-base  & $22.80\pm9.00\%$ & $15.21\pm9.13\%$ & $2.72\times$ ($1.19$) & $1.81\times$ ($1.08$) & $1.61\times$ ($1.06$) \\
Flan-T5-large & $27.64\pm5.56\%$ & $11.62\pm5.97\%$ & $3.29\times$ ($1.27$) & $1.38\times$ ($1.04$) & $1.53\times$ ($1.05$) \\
\midrule
\multicolumn{3}{@{}l}{spread of per-seed multipliers} & $1.74$ & $0.56$ & $0.11$ \\
\multicolumn{3}{@{}l}{large $-$ small, paired $t$} & $\mathbf{+5.19}$ & $+0.62$ & $+0.24$ \\
\multicolumn{3}{@{}l}{base $-$ small, paired $t$} & $\mathbf{+3.03}$ & $+1.25$ & $+0.23$ \\
\multicolumn{3}{@{}l}{large $-$ small, paired $t$, regret ratio} & $\mathbf{+5.66}$ & $+0.92$ & $+0.15$ \\
\multicolumn{3}{@{}l}{base $-$ small, paired $t$, regret ratio} & $\mathbf{+3.33}$ & $+1.74$ & $+0.33$ \\
\bottomrule
\end{tabular}
\end{table}

\begin{figure}[h]
\centering
\includegraphics[width=0.86\linewidth]{two_channel.png}
\caption{The two channels on AmazonCat/BERT-768d, ten seeds, with standard
errors over seeds. The three experts have held-out accuracies of $5.20\%$,
$3.61\%$ and $5.27\%$ and cover the same directions to within $5\%$, so
Corollary~\ref{cor:coverage-complete} leaves the transfer as the only place
they can differ. With it present (dark bars) they span $1.71\times$ to
$3.29\times$; with it removed by either mechanism they collapse onto a common
multiplier near $1.5\times$ and no pairwise difference survives. The two
transfer-free conditions differ in what else they delete---\emph{transfer
removed} keeps the expert-independent part $M\theta_{\mathrm{or}}$ of the
response sum, \emph{coverage only} discards the response sum
entirely---yet they agree with each other on every expert, as they must if the
transfer is the only thing both lack.}
\label{fig:two-channel}
\end{figure}

\begin{table}[h]
\centering
\caption{Confidence sweep on the BGE-1024d configurations, tuning seeds
$11$--$20$, which select $c$. \todo{WikiLinks sweep} At the
$c=100$ used for BERT the expert width $u^{(1)}=\|\phi\|/\sqrt c$ is
permanently below LinUCB's, the gate never engages, and the Router equals the
LLM exactly on every seed ($V_\tau=0$, zero variance). Lowering $c$ raises the
expert's width and restores the gate. On AmazonCat, where cold LinUCB is
strong, the Router at $c=0.1$ beats it while routing $0.05\%$ of rounds, and
the same $c$ reproduces the gain on the evaluation seeds
(Table~\ref{tab:amazon-results}); at $c\ge1$ it is dragged down toward a
weaker expert. On WikiLinks the expert is
genuinely better than the learner---row 3 of Table~\ref{tab:catalog}, where no
combination can beat the better component.}
\label{tab:bge-csweep}
\scriptsize
\setlength{\tabcolsep}{4pt}
\begin{tabular}{@{}llccccccc@{}}
\toprule
& & \multicolumn{3}{c}{routed \%} & \multicolumn{3}{c}{Router accuracy} & \\
\cmidrule(lr){3-5}\cmidrule(lr){6-8}
Dataset & $c$ & 0.5B & 3B & 7B & 0.5B & 3B & 7B & best mult.\ \\
\midrule
\multirow{4}{*}{AmazonCat} & $0.01$ & $0.0$ & $0.0$ & $0.0$ & $34.77$ & $34.77$ & $34.77$ & $1.00\times$ \\
 & $\mathbf{0.1}$ & $0.04$ & $0.05$ & $0.05$ & $33.84$ & $39.68$ & $\mathbf{40.63}$ & $\mathbf{1.17\times}$ \\
 & $1$  & $75.2$ & $71.9$ & $73.5$ & $28.09$ & $29.26$ & $36.50$ & $1.05\times$ \\
 & $10$ & $100$ & $100$ & $100$ & $21.38$ & $22.39$ & $27.17$ & $0.78\times$ \\
\cmidrule(lr){1-9}
\multirow{5}{*}{WikiLinks} & $0.01$ & --- & --- & --- & --- & --- & --- & --- \\
 & $0.1$ & --- & --- & --- & --- & --- & --- & --- \\
 & $1$  & --- & --- & --- & --- & --- & --- & --- \\
 & $10$ & --- & --- & --- & --- & --- & --- & --- \\
 & $100$ & --- & --- & --- & --- & --- & --- & --- \\
\bottomrule
\end{tabular}
\end{table}

\begin{figure}[h]
\centering
\includegraphics[width=0.85\linewidth]{multiplication.png}
\caption{Learning dynamics on AmazonCat-13K (BERT-768d, $c=100$), mean over
ten seeds with shaded standard error. The shaded window contains every routed
round: all calls to the LLM fall inside the first $5{,}000$ rounds, after which
the expert is switched off. The three Flan-T5 experts (held-out accuracies
$5.20\%$, $3.61\%$ and $5.27\%$, the band at the bottom) are consulted on
equally few rounds, yet their Routers separate by nearly a factor of two; and
$96$--$100\%$ of each Router's advantage over cold LinUCB, counted in correct
answers, accrues after its last routed round. That is the row-$5'$ signature:
the saving is collected on the tail.}
\label{fig:amazon-dynamics}
\end{figure}

\textbf{The multiplying configuration.} Flan-T5-large reaches
$27.64\pm5.56\%$ against cold LinUCB's $8.39\pm2.46\%$ with only $3.2\%$ of
rounds routed to the LLM overall; routing is concentrated in the first
$\sim5{,}000$ rounds. Table~\ref{tab:amazon-decomp} gives the shape of a
positive case directly. $A_{\tau}=53$ is negligible: it is positive, not
negative, because the routed rounds fall early, while cold LinUCB is still
near $2\%$, so even a $5\%$ expert is marginally the better player there; but
at $1{,}578$ routed rounds the whole term is worth $53$ correct answers
against a total advantage of $9{,}624$. $V_{\tau}=9{,}571$ carries $99.4\%$ of
it, and is positive on all ten seeds. $B_{\tau}=11{,}189$ is large because the
expert stays at $5\%$ while the warm-started learner passes $27\%$. This is
the row-$5'$ signature of Proposition~\ref{prop:cat-tail-saving}: the
acquisition is nearly free, and the saving is collected on the tail, after
routing has stopped. Figure~\ref{fig:amazon-dynamics} shows the same
accounting in time: every routed round falls inside the first $5{,}000$, and
$96\%$ of the $9{,}624$ correct answers the Router gains over cold LinUCB are
gained after its last routed round, while the expert itself stays frozen at
$5.3\%$. The cumulative-accuracy curves part earlier than that only because an
early lead weighs heavily in a short prefix.

\textbf{Equal accuracy, unequal information---and which channel carries it.}
The three Flan-T5 experts have held-out accuracies $5.20\%$, $3.61\%$ and
$5.27\%$ and deliver accuracy multipliers of $1.71\times$, $2.72\times$ and
$3.29\times$ (regret ratios $1.07$, $1.19$ and $1.27$): accuracy orders them
differently from the multiplier, and the least accurate expert
beats a more accurate one by $60\%$. By Proposition~\ref{prop:decomp} the only
place they can differ is $V_\tau$, and Table~\ref{tab:amazon-decomp} confirms
it: $A_\tau$ is $47$, $27$ and $53$, indistinguishable as the equal accuracies
require, while $V_\tau$ is $2{,}936$, $7{,}178$ and $9{,}571$.

Section~\ref{sec:two-channels} says where that difference can live. The three
experts cover the same directions to within $5\%$: summed over the routed
rounds, the coverage of the directions cold LinUCB is trapped on is
$1.40\times10^5$, $1.46\times10^5$ and $1.47\times10^5$, at prices
$9.5\times10^{-3}$, $1.03\times10^{-2}$ and $1.04\times10^{-2}$, so the price
ratio of Theorem~\ref{thm:general-price}, measured post hoc as described in
Section~\ref{sec:exp-setup}, is $1.73$, $1.59$ and $1.58$ for the
three---equal to within measurement error, and ordered opposite to the observed
multipliers. It bounds the regret ratio, and the observed regret ratios lie
below it for every expert; since the real stream is outside the theorem's
setting we read this as consistency, not confirmation. By Corollary~\ref{cor:coverage-complete} the observed
separation must then come from the transfer, and
Table~\ref{tab:two-channel} and Figure~\ref{fig:two-channel} show that it
does.

Removing $m$ while leaving the binary problem otherwise intact collapses the
spread of per-seed multipliers from $1.74$ to $0.56$ and turns every
significant pairwise difference into a null result: large over small falls
from $t=5.19$ to $t=0.62$. Suppressing the response sum entirely gives the
same collapse, to a spread of $0.11$; and the two transfer-free conditions
agree with each other on every expert (paired $t=-1.18$, $+0.26$, $-0.60$),
as they must if $m$ is the only thing both lack. Removing the transfer costs
each expert in proportion to what it gained, $-2.22$ for large ($t=5.12$)
against $-0.63$ for small ($t=2.47$). A third route to the same conclusion
replaces the reward by $\langle\theta_{\mathrm{or}},\phi\rangle+\mathcal
N(0,\varsigma^2)$, making the model well specified so that
Proposition~\ref{prop:response-sum} applies exactly: at
$\varsigma\in\{0.001,0.01,0.499\}$ the three experts' regret ratios, measured
against the oracle's best action, span $0.019$, $0.010$ and $0.030$, against
$0.197$ on the real reward; the ratios themselves lie between $0.97$ and
$1.09$. Under linear rewards the experts are interchangeable, exactly as
predicted.

\textbf{Where the regret goes.} On the regret scale the two channels have
very different sizes. With the response sum suppressed, coverage alone gives
regret ratios $1.051$, $1.060$ and $1.051$: about $5\%$ less regret than cold
LinUCB, indistinguishable across experts, and below the price ratio for each,
as it must be. The full Router gives $1.070$, $1.187$ and $1.266$. For
Flan-T5-base and Flan-T5-large, therefore, most of the regret saved is saved
by the transfer, roughly $10$ of $16$ and $16$ of $21$ percentage points. The
separation between experts survives the change of scale intact: on per-seed
regret ratios, large beats small with paired $t=+5.66$ and base beats small
with $t=+3.33$ when the transfer is present, and neither difference is
significant with the transfer removed ($t=+0.92$, $+1.74$) or the response sum
suppressed ($t=+0.15$, $+0.33$); Figure~\ref{fig:two-channel} shows the same
collapse on the accuracy scale. The price ratio correctly predicts that
coverage makes the three experts interchangeable; it does not predict the size
of the coverage gain, which it bounds loosely.

Finally, the transfer is not a small correction to the estimate at handoff; on
this data it is essentially the whole of it. The ridge oracle assigns a
correct action a score of about $2.5\times10^{-3}$ where the reward is $1$, so
the residuals dominate: $\|m\|/\|\widehat\theta_{\tau+1}\|=0.9998$, and the
identity of Proposition~\ref{prop:transfer} holds on the logged trajectories
to $1.2\times10^{-14}$. Stripping $m$ makes the three estimates agree to four
decimals in held-out ranking quality ($0.2058$, $0.2058$, $0.2057$), while the
transfers themselves separate the experts with the same resolution the
multipliers have: base and large each beat small ($t=3.12$, $t=2.56$) and are
indistinguishable from one another ($t=0.36$), which is the pair the
multipliers also cannot separate.

\textbf{What the transfer carries.} On AmazonCat one category is present on
$29.8\%$ of contexts, so a policy that knew this in advance and always
predicted it would score $29.8\%$. An online learner does not know it, and
finding it is not easy: context-free bandits that try categories without
using the context reach $0.04\%$ (UCB1) and $2.31\%$ (optimistic greedy) over
the same $50{,}000$ rounds and $13{,}330$ categories, seeds $1$--$10$. The
Routers that multiply are the ones that discover this category. LinUCB
settles on one frequent category and predicts it on a large share of rounds,
and its accuracy then tracks that category's frequency: runs that settle on
the dominant category finish near $28\%$, while cold LinUCB settles on a
category present on $11\%$ of contexts, or on none, in $18$ of $20$ seeds.
Whether a Router settles on the dominant category is decided by a handful of
the expert's routed answers. Across $177$ Router runs on BERT features (every
$c$ and CORRAL, excluding configurations that route most rounds), a run whose
expert gave no correct routed answer on this category reached it in $3$ of
$45$ runs, and a run that received five or more reached it in all $54$. In a
controlled test, adding $15$ correct answers on this category to
Flan-T5-small ($+0.15$ points of held-out accuracy) doubled its multiplier
($1.33\times$ to $2.62\times$, paired $p=0.002$), while $15$ correct answers on
random categories gave no significant gain ($1.65\times$, $p=0.15$). This is
the transfer of Proposition~\ref{prop:transfer} made concrete: correct answers
on a frequent category are the residuals that move the estimate, and they are
why accuracy does not order the experts. Appendix~\ref{app:dominant} gives the
details.

\textbf{BGE-1024d.} At $c=100$ the expert width $u^{(1)}=\|\phi\|_2/\sqrt c$
is permanently below LinUCB's $u^{(0)}$, so the Router defers on every round
and equals the LLM exactly: $V_\tau=0$ with zero variance across seeds. This
is a failure of the gate's hyperparameter, not evidence about $V_\tau$, and it
affected all six BGE configurations. Table~\ref{tab:bge-csweep} sweeps $c$
down until the gate engages. The sweep is run on seeds $11$--$20$ and selects
$c=0.1$; the same value is then evaluated on seeds $1$--$10$. There cold LinUCB
reaches $32.59\%$, while the Router reaches $40.20\%$ with Qwen-3B and
$38.71\%$ with Qwen-7B, accuracy multipliers of $1.23\times$ and $1.19\times$
(regret ratios $1.13$ and $1.10$), while routing $0.05\%$ of rounds. Both gains
are significant (paired $t=3.92$ and $2.59$ against cold, $t^\star=2.26$). With
Qwen-0.5B the gate engages just as often and the Router does not move at all
($1.00\times$, $t=0.01$), which is the expert-dependence the BERT rows also
show. The accounting has the BERT signature where there is a gain: $A_\tau$ is
$5$ and $6$ correct answers while $V_\tau$ is $3{,}801$ and $3{,}056$, positive
on $9$ and $8$ of the ten seeds, against $V_\tau=13$ for Qwen-0.5B. The
multipliers are smaller than on BERT, which is what a stronger learner implies:
when cold LinUCB already reaches a third of the correct answers, there is less
left to buy. In the terms of the previous paragraph, cold LinUCB on BGE
features settles on the dominant category by itself in $7$ of $20$ seeds, so
the expert has less to add (Appendix~\ref{app:dominant}). Pushing $c$ higher
drags the Router toward a weaker expert and below cold LinUCB. On WikiLinks
the Qwen experts are better than the learner ($29.47\%$ and $38.74\%$ against
$18.88\%$), which is row~3 of Table~\ref{tab:catalog}: an expert at or above the
learner's ceiling, where $B_\tau\le0$ and no combination can beat the better
component. \todo{WikiLinks sweep}

\paragraph{Summary across configurations.}
Twelve configurations span two datasets, two model families (60M to 7B), two
embedding spaces and four orders of magnitude in $c$, with ten evaluation
seeds per row and every $c$ chosen on a disjoint block of ten tuning seeds. One configuration family
multiplies clearly: AmazonCat with BERT-768d and Flan-T5 (accuracy multipliers
$1.71\times$, $2.72\times$, $3.29\times$; regret ratios $1.07$, $1.19$,
$1.27$). AmazonCat with BGE-1024d and Qwen, evaluated on seeds held out from
the choice of $c$, gives a smaller but significant gain (up to $1.23\times$ and
$1.13$, paired $t=3.92$). The six WikiLinks configurations do not \todo{WikiLinks sweep};
neither does BGE at the $c$ that routes everything. In every configuration
that multiplies, the accounting has the same shape: $A_\tau$ is a few dozen
correct answers out of more than a thousand routed rounds, $V_\tau$ is in the
thousands, and $B_\tau$ is large because the expert stays poor while the
warm-started learner does not. Where multiplication fails, the failing term of
Proposition~\ref{prop:decomp} can be named: on WikiLinks with Flan-T5 cold
LinUCB already resolves the directions that matter, so $V_\tau$ is positive
but small; with Qwen the expert is above the learner's ceiling and $B_\tau<0$;
at large $c$ on BGE the gate never hands back and $V_\tau=0$ identically.

Across the three Flan-T5 experts the coverage channel is
expert-independent, as Corollary~\ref{cor:coverage-complete} requires, and the
transfer accounts for the whole of the observed separation and for most of the
regret saved (Table~\ref{tab:two-channel}, Figure~\ref{fig:two-channel}). The deterministic Router differs from the
randomized CORRAL master of \cite{agarwal2017corralling} and the CORRAL-style
routing of \cite{chen2024llmdecision} only in how it chooses the routed
rounds; both feed every observation to the bandit learner, so the
decomposition applies to both, and Proposition~\ref{prop:response-sum} adds
that two routers with the same routed Gram matrix leave warm starts of the
same law---any difference between routers must therefore appear in $A_\tau$ or
in which residuals their routed sets happen to collect.
Table~\ref{tab:baselines} runs both alternatives.

\begin{table}[h]
\centering
\caption{Alternative routing mechanisms on AmazonCat-13K (BERT-768d,
$c=100$), ten seeds, paired on seed with the gate. CORRAL is the randomized
log-barrier master of \cite{agarwal2017corralling} with importance-weighted
updates; the offline arms spend $\tau$ expert pulls on contexts drawn from
\emph{outside} the online stream and then run plain LinUCB, so their
acquisition is not charged as regret and $A_\tau=0$ identically.
$\tau=1{,}600$ matches the gate's routed count. Cold LinUCB is bit-identical
across all four arms on every seed, so each comparison isolates the router.
The last column is the paired $t$ statistic for gate minus that arm over the
ten shared seeds ($t^\star=2.26$). Mult.\ is the accuracy multiplier, with
the regret ratio in parentheses.}
\label{tab:baselines}
\small
\setlength{\tabcolsep}{4pt}
\begin{tabular}{@{}llrcccrr@{}}
\toprule
LLM & Router & queries & Router acc.\ & Mult.\ & $A_\tau$ & $V_\tau$ & $t$ vs.\ gate \\
\midrule
\multirow{4}{*}{Flan-T5-small}
 & Gate (Alg.\ 1)    & $1{,}406$  & $14.36\pm6.12\%$ & $1.71\times$ ($1.07$) & $47$   & $2{,}936$ & --- \\
 & CORRAL            & $22{,}860$ & $8.87\pm1.37\%$  & $1.06\times$ ($1.01$) & $-508$ & $744$     & $\mathbf{2.64}$ \\
 & Offline $\tau{=}1{,}600$ & $1{,}600$ & $11.17\pm5.71\%$ & $1.33\times$ ($1.03$) & $0$ & $1{,}386$ & $1.16$ \\
 & Offline $\tau{=}5{,}000$ & $5{,}000$ & $15.95\pm7.11\%$ & $1.90\times$ ($1.09$) & $0$ & $3{,}778$ & $-0.59$ \\
\midrule
\multirow{4}{*}{Flan-T5-base}
 & Gate (Alg.\ 1)    & $1{,}570$ & $22.80\pm9.00\%$ & $2.72\times$ ($1.19$) & $27$  & $7{,}178$ & --- \\
 & CORRAL            & $6{,}939$ & $24.44\pm2.39\%$ & $2.91\times$ ($1.21$) & $-23$ & $8{,}045$ & $-0.63$ \\
 & Offline $\tau{=}1{,}600$ & $1{,}600$ & $20.97\pm8.89\%$ & $2.50\times$ ($1.16$) & $0$ & $6{,}286$ & $0.42$ \\
 & Offline $\tau{=}5{,}000$ & $5{,}000$ & $27.29\pm5.17\%$ & $3.25\times$ ($1.26$) & $0$ & $9{,}449$ & $-2.00$ \\
\midrule
\multirow{4}{*}{Flan-T5-large}
 & Gate (Alg.\ 1)    & $1{,}597$ & $27.64\pm5.56\%$ & $3.29\times$ ($1.27$) & $53$ & $9{,}571$ & --- \\
 & CORRAL            & $7{,}488$ & $25.42\pm2.31\%$ & $3.03\times$ ($1.23$) & $50$ & $8{,}460$ & $1.49$ \\
 & Offline $\tau{=}1{,}600$ & $1{,}600$ & $26.27\pm7.61\%$ & $3.13\times$ ($1.24$) & $0$ & $8{,}937$ & $0.42$ \\
 & Offline $\tau{=}5{,}000$ & $5{,}000$ & $25.76\pm6.53\%$ & $3.07\times$ ($1.23$) & $0$ & $8{,}681$ & $0.67$ \\
\bottomrule
\end{tabular}
\end{table}

Three things follow. First, no baseline beats the gate significantly on any
expert, and the gate beats CORRAL on Flan-T5-small ($t=2.64$) while issuing
$16\times$ fewer expert queries; on the other two experts CORRAL matches the
gate at four to five times the query budget. Second, the offline arm at a
matched budget of $1{,}600$ pulls is statistically indistinguishable from the
gate on all three experts even though its acquisition is free, so the online
routing decision costs essentially nothing---which is the content of
$A_\tau\approx0$, now visible from the other side. Tripling the free budget to
$5{,}000$ still does not separate the two. Where an alternative does come out
ahead of the gate in the mean, the margin is small, never significant, and
paid for with a much larger budget: CORRAL on Flan-T5-base gains $1.6$
points ($t=-0.63$) with $4.4\times$ the gate's expert queries, and the
offline arm at $\tau=5{,}000$ gains $1.6$ and $4.5$ points on Flan-T5-small
and Flan-T5-base ($t=-0.59$, $-2.00$) only by drawing $3$--$4\times$ as many
expert answers, all collected before the online run begins. None of these
trades buys a reliable improvement, so the gate is the cheapest router in
the comparison and no alternative is significantly better on any expert.

Third, and most directly, the decomposition survives the change of master. In
all twelve rows $V_\tau$ carries the value and $A_\tau$ is negligible beside
it, and $A_\tau$ is exactly where the choice of router shows up: it is $0$ by
construction offline, a few dozen for the gate, and \emph{negative} for CORRAL
on the two weaker experts ($-508$ and $-23$), which is the price of routing
$22{,}860$ rounds to a $5.2\%$ expert while cold LinUCB is improving. The
comparison also shows that coverage is not bought by volume: CORRAL routes
$16\times$ as many rounds as the gate on Flan-T5-small and ends with
$V_\tau=744$ against the gate's $2{,}936$. What the gate supplies is not more
expert pulls but pulls on the directions cold LinUCB is trapped on, which is
what $\pi(u)$ in Theorem~\ref{thm:general-price} measures and what
Proposition~\ref{prop:response-sum} says is the whole of the coverage channel.

\section{Conclusion}
\label{sec:conclusion}

Reward multiplication in LLM--bandit hybrids is an accounting statement: the
hybrid beats cold LinUCB by the expert's direct advantage plus the value of the
data it leaves behind, and beats the expert by the learner's tail advantage.
The data-value term is governed by whether the routed actions push the
confidence widths of dangerous competitors below their gaps; the
determinant-ratio term that controls $\sqrt T$-rates cannot produce more than
a constant factor. When the directions that trap cold LinUCB recur and the
expert has a cheaper way of covering them, an online handoff multiplies, and
the multiplier is the ratio of the two prices; when any one of recurrence,
cheapness, correct targeting, or late suboptimality of the expert is removed,
it does not.

A warm start, however, carries two objects, and the price accounting governs
one of them. Under the linear model the second is redundant: the response sum
is the Gram matrix applied to $\theta^\star$ plus noise whose law depends only
on that matrix, so experts of equal coverage are interchangeable and the price
ratio is a complete account of information value
(Proposition~\ref{prop:response-sum}). Real rewards are not linear, and there
the estimate at handoff acquires a transfer---the ridge projection of the
residuals onto the expert's own picks---which we derive in closed form and
bound (Propositions~\ref{prop:transfer} and~\ref{prop:transfer-regret}). On
AmazonCat it is the larger channel. A linear function of the joint features
assigns a correct action a score of $2.5\times10^{-3}$ against a reward of $1$,
so the transfer is essentially the entire estimate at handoff. Coverage alone
saves about $5\%$ of cold LinUCB's regret, the same for every expert, while
the full router saves up to $21\%$; and every significant difference between
three experts of equal accuracy disappears when the transfer is removed,
whether by suppressing the response sum, by replacing the routed-round signal
with its expert-independent part, or by linearizing the reward. The practical
consequence is that a warm-start expert cannot be chosen by its accuracy, nor
by the coverage its actions provide: what decides its value is the residuals
its choices carry---on AmazonCat, a handful of correct answers on the most
frequent category. Information value is not policy value; and on real data it
is not coverage alone either.

\paragraph{Limitations.}
The theory assumes linear reward, and on our main dataset the channel it
leaves out is the larger one: the transfer has a closed form and a regret bound
(Proposition~\ref{prop:transfer-regret}), but no price-style quantity that
predicts its size before the expert is run. Reward models suited to binary
feedback, such as generalized linear bandits, may shrink the residuals the
transfer is built from; we have not tested whether the two-channel picture
changes under them. The price ratio we report on real data is measured after
the run, from the trajectories it is compared against, on streams outside the
fixed-arm setting of Theorem~\ref{thm:general-price}; it is an explanatory
statistic, not a prediction. Theorem~\ref{thm:general-price} bounds the
online multiplier on any finite-arm instance with the optimal arm always
available; real streams satisfy this only approximately, and the bound has a
residual---cold LinUCB's regret on arms the warm start does not certify---that
is not itself governed by a price. We have no lower bound showing that the
residual can be large without a large price ratio, and no lower bound on
$R_T(\Hyb_\tau)$ beyond the acquisition regret, so the theorem characterizes
the certified saving rather than the multiplier outright. The positive
instances (Theorem~\ref{thm:price}, Proposition~\ref{prop:cat-tail-saving})
use a fixed confidence radius and a regularization scaled to the dimension,
and their constants are not optimized.
The constants in the catalog are not optimized and the lower bound on cold
LinUCB's regret there counts only the rounds that optimism forces. The pre-runtime estimates require a labeled held-out sample and the
expert's proposals on it; their thresholds are provisional and were not tuned
across datasets. The extension of the width-based analysis to other
confidence-based learners (OFUL, SupLinUCB, Thompson sampling) is immediate
for Lemmas~\ref{lem:loewner} and~\ref{lem:transfer}, which are statements
about design matrices, but Lemma~\ref{lem:width-limited} uses the optimistic
selection rule and would need an analogue for posterior sampling. The large
AmazonCat multipliers depend on a dominant category that the learner can be
steered to (Appendix~\ref{app:dominant}); where no category dominates, as on
WikiLinks, we see no multiplication, and we have not tested other datasets
with a skewed label distribution. Only two
real datasets are used; the effect is large in one feature space and modest in
the other, where cold LinUCB is already strong; multi-expert behaviour is not
stress-tested;
LLM-guided recommendation can inherit biases from pretrained models and
logged data.

\paragraph{Broader impact.}
The work is foundational theory and pre-runtime estimation for contextual
bandits. A positive impact is reducing unnecessary LLM inference in
sequential decision systems by identifying when cold-start LinUCB suffices,
which lowers compute cost and energy use. A potential negative impact is that
LLM-guided contextual-bandit systems deployed in recommendation or ranking
may inherit biases from pretrained models or logged data without auditing;
concrete harms depend on follow-on applications not studied here.

\bibliography{references}
\bibliographystyle{abbrv}


\appendix

\section{Related Work}\label{sec:related}
This appendix expands on the related-work discussion deferred from the
introduction. We cover (i) linear contextual bandits and confidence geometry,
(ii) model selection and expert aggregation, (iii) warm starts and prior
information, (iv) the geometry of information gain in linear bandits, and
(v) large language models for sequential decision making.

Sequential decision making with contextual information has been studied
extensively in the contextual bandit literature, where the learner must
balance exploration and exploitation under partial feedback
\cite{langford2007epoch,agarwal2014taming,foster2021statistical,
lattimore2020bandit}. In the linear setting, algorithms such as LinUCB and
OFUL achieve regret bounds that depend explicitly on the geometry of the
accumulated design matrix through self-normalized confidence ellipsoids
\cite{chu2011contextual,abbasi2011improved}. Subsequent work has developed
statistically efficient contextual bandit algorithms under general function
approximation assumptions \cite{agarwal2012contextual,foster2018practical,
foster2020beyond,simchi2022bypassing,zhu2022contextual,rucker2023infinite,
zhang2024efficient}. A central insight of this line of work is that the
conditioning of the covariance matrix determines how quickly uncertainty
shrinks and therefore how rapidly regret decreases. Our work builds on this
perspective by showing that actions selected by a frozen predictor can
reshape this geometry early in learning and accelerate subsequent adaptation.

\paragraph{Model selection and expert aggregation in bandits.}
A closely related direction studies how to combine multiple candidate
policies or base algorithms online. Classical expert-learning methods provide
regret guarantees relative to the best fixed expert in hindsight
\cite{auer2002nonstochastic}. In the bandit setting, CORRAL
\cite{agarwal2017corralling} uses Online Mirror Descent with a log-barrier
regularizer to achieve regret competitive with the best base algorithm;
subsequent refinements tighten the rates and extend the framework to
contextual and stochastic settings
\cite{pacchiano2020model,marinov2021pareto,zhu2022pareto,zhu2022near,
dann2024data}. These approaches are inherently \emph{randomized} and target
robustness: the master performs nearly as well as the best base algorithm
regardless of which one that is. Our objective is different---we want a
\emph{deterministic} rule that exploits the complementarity between a good
early frozen policy and a strong asymptotic adaptive learner to
\emph{strictly outperform} both. Correspondingly, our analysis does not
bound a relative regret against the best base, but instead tracks how early
actions selected by one policy reshape the design matrix and accelerate the
statistical efficiency of another.

\paragraph{Warm starts and prior information in contextual bandits.}
Several works have investigated how prior information can be incorporated
into contextual bandit algorithms through initialization, offline datasets,
or pretrained predictors. Such approaches can substantially improve
short-horizon performance by reducing uncertainty before online interaction
begins \cite{bastani2021mostly,zhang2019warm}. Related directions in
transfer learning and structured bandits further show that leveraging
information across tasks or user populations can accelerate adaptation and
reduce exploration requirements \cite{lazaric2012transfer,gentile2014online}.
Meta-learning approaches similarly learn priors that enable faster
exploration in new environments \cite{sharaf2019meta}. Empirical work on
neural-linear bandit architectures also demonstrates that representation
pretraining can significantly improve early decision quality
\cite{riquelme2018deep}. Our analysis formalizes these intuitions in a
linear bandit framework: informative early actions contract the confidence
ellipsoid on specific competitors (Lemmas~\ref{lem:loewner}
and~\ref{lem:transfer}), and the resulting saving is a difference of
width-limited regrets rather than an improvement of the
determinant-ratio rate, which Proposition~\ref{thm:handoff} shows cannot by
itself exceed a constant factor.

\paragraph{Geometry of information gain in linear bandits.}
A central theme in the analysis of linear bandit algorithms is that learning
efficiency is governed by the geometry of the accumulated design matrix.
Early work on stochastic linear optimization under bandit feedback showed
that regret depends explicitly on the growth of the covariance matrix
associated with the chosen actions and the resulting confidence ellipsoid
\cite{dani2008stochastic}. Subsequent analyses of linearly parameterized
bandits further clarified how the eigenvalue structure of this matrix
determines parameter identifiability and sample complexity
\cite{rusmevichientong2010linearly}. Later work established sharp
self-normalized confidence bounds and finite-time regret guarantees
expressed directly in terms of the design matrix and its inverse
\cite{abbasi2011improved}, while related results connected spectral growth
of the covariance matrix to optimal arm identification in linear settings
\cite{soare2014best}. From an information-theoretic perspective, regret can
also be related to determinant growth and effective dimension of the
covariance operator \cite{srinivas2012information,valko2013finite}. The
price \eqref{eq:price} is the single-direction shadow of the allocation
programs behind asymptotic instance-dependent lower bounds, where regret is
minimized subject to coverage constraints on every direction
\cite{graves1997asymptotically,lattimore2017end}; our use of it is
non-asymptotic and comparative: two policies on the same instance pay
different prices for the same coverage. Prior side information has also been
analyzed as hints or predictable sequences in online learning
\cite{rakhlin2013online,cutkosky2022leveraging}. The comparison with
\cite{cutkosky2022leveraging} is the closest: there the learner receives a
hint about $\theta^\star$ and the regret improves as a function of the hint's
accuracy, with the standard rate recovered when the hint is poor. Our warm
start is not a hint about $\theta^\star$ but data in the design matrix, and
its value is not a function of accuracy: two experts of equal accuracy can
leave design matrices of very different value (Section~\ref{sec:amazon}),
and an expert of zero accuracy can leave one of enormous value
(Proposition~\ref{prop:three-way}). What governs the value is the price ratio
of Theorem~\ref{thm:general-price}. The warm-start analyses of
\cite{bastani2021mostly,zhang2019warm} bound the regret of a learner started
from logged data through the quality of that data; they are offline in our
sense (row~1 of Table~\ref{tab:catalog}), and our contribution relative to
them is the online accounting, in which the acquisition of the data is
charged and can exceed its value. Our
work extends this geometric viewpoint by showing that a fixed policy can
shape the spectrum of the design matrix during early rounds of interaction,
thereby improving conditioning of the confidence ellipsoid and accelerating
the learning rate of subsequent adaptive exploration.

\paragraph{Large language models for sequential decision making.}
Recent work has explored the use of large language models (LLMs) as policies
in contextual bandit environments, demonstrating that frozen pretrained
models can achieve strong short-horizon performance even without online
updates \cite{chen2024llmdecision}. More broadly, LLMs exhibit strong
in-context learning capabilities \cite{brown2020language,xie2021explanation,
garg2022can,akyurek2022learning}, and instruction tuning further improves
their effectiveness in downstream decision tasks
\cite{ouyang2022training,chung2024scaling}. Several recent approaches
interpret sequential decision making itself as a sequence-modeling problem
and adapt transformer architectures to interactive settings
\cite{chen2021decision,zheng2022online,reid2022can,sun2023smart,
raparthy2023generalization,lee2023supervised}.
Closest to our work is the approach of \cite{chen2024llmdecision}, who
combine an LLM policy with an adaptive contextual bandit learner using an
online model-selection master and empirically document the
reward-multiplication phenomenon on large-scale recommendation datasets.
Our contribution relative to that line is theoretical: we give a
regret-based explanation for why such combinations can beat both components,
and replace the randomized CORRAL-style master with a deterministic
confidence-gated rule whose routing decision directly exposes the geometric
warm-start mechanism. A concurrent line of work \cite{huangcontextual} also
uses LLM uncertainty to gate exploration, but does so by calibrating an
embedding-based uncertainty estimate to adjust UCB-style exploration
bonuses, rather than by comparing self-normalized confidence radii of a
frozen predictor and an online learner; their analysis bounds the regret
of a single calibrated agent through the accuracy of the LLM's reward
predictions, whereas we give an exact decomposition of the hybrid's
advantage, bound the multiplier over the cold learner on any finite-arm
instance by a ratio of coverage prices, and show the bound is attained; in
their setting the LLM's value enters through prediction error, in ours
through the design matrix, and the two can disagree by a factor of the
horizon (Proposition~\ref{prop:three-way}).
%% NOTE: the next two citations need references.bib entries:
%%   jumpstart2026   -- "Jump Start or False Start? A Theoretical and Empirical
%%                      Evaluation of LLM-initialized Bandits", arXiv:2604.02527
%%   calibgated2026  -- "Calibration-Gated LLM Pseudo-Observations for Online
%%                      Contextual Bandits", arXiv:2604.14961
Two concurrent works are closer still. \cite{jumpstart2026} warm-start a linear
bandit from LLM-generated preferences and bound the warm-started estimator's
error through a single prior-error term measured in the geometry of the
pretraining design matrix, which treats the design matrix and the response
vector as one object; our analysis separates the two, shows that under linear
rewards the second is determined by the first, and identifies the
expert-dependent remainder in closed form, which is the part that differs
across equally accurate experts on our data. \cite{calibgated2026} inject LLM
pseudo-observations into LinUCB with a calibration-dependent weight on both the
design matrix and the response vector, a routing rule close in spirit to ours,
and analyze it through a bias--variance argument without a regret bound; the
decomposition of Proposition~\ref{prop:decomp} and the price bound of
Theorem~\ref{thm:general-price} apply to any such rule, since both hold for an
arbitrary routed set. Additional evidence that LLMs can
exhibit in-context reinforcement learning behavior
\cite{monea2024llms,krishnamurthy2024can} and that LLM-derived priors can
jump-start exploration \cite{alamdari2024jumpstarting,abbasiyadkori2024believe}
further motivates the study of principled mechanisms for combining frozen
predictors with online learners.


\section{Proofs}
\label{app:proofs}

Section~\ref{app:proof-decomp} proves the decomposition identity.
Section~\ref{app:proof-three-way} reproduces the proof of
Proposition~\ref{prop:three-way} from the original construction.
Section~\ref{app:proof-one-pass} proves Proposition~\ref{prop:one-pass}.
Sections~\ref{app:proof-width}--\ref{app:proof-transfer} prove the three
lemmas of Section~\ref{sec:width}, and Section~\ref{app:proof-handoff}
proves Proposition~\ref{thm:handoff}.

\subsection{Proof of Proposition~\ref{prop:decomp}}
\label{app:proof-decomp}

\begin{proof}
Split $[T]$ into $[1,\tau]$ and $[\tau+1,T]$. On $[1,\tau]$ the hybrid plays
$a_t^{\Exp}$, so $R_{[1,\tau]}(\Hyb_\tau)=R_{[1,\tau]}(\Exp)$. Hence
\[
R_T(\Cold)-R_T(\Hyb_\tau)
=\bigl(R_{[1,\tau]}(\Cold)-R_{[1,\tau]}(\Exp)\bigr)
+\bigl(R_{[\tau+1,T]}(\Cold)-R_{[\tau+1,T]}(\Hyb_\tau)\bigr)
=A_\tau+V_\tau ,
\]
and
\[
R_T(\Exp)-R_T(\Hyb_\tau)
=\bigl(R_{[1,\tau]}(\Exp)-R_{[1,\tau]}(\Exp)\bigr)
+\bigl(R_{[\tau+1,T]}(\Exp)-R_{[\tau+1,T]}(\Hyb_\tau)\bigr)
=B_\tau .
\]
Both are identities between sums of gaps, so they hold for every noise
realization and therefore in expectation. The ``if and only if'' statements
are restatements of the identities.
\end{proof}

\subsection{Proof of Proposition~\ref{prop:three-way}}
\label{app:proof-three-way}

We first record the standard subgaussian concentration lemma used in the
high-probability part of the argument, and then construct the linear bandit
instance and verify the three claims of Proposition~\ref{prop:three-way}.

\begin{lemma}[Subgaussian concentration, Corollary 5.5 of \cite{lattimore2020bandit}]
\label{lem:subgaussian}
Assume that $X_1-\mu,\dots,X_n-\mu$ are independent, $\sigma$-subgaussian random variables. Then for any $\varepsilon \ge 0$,
\[
\mathbb{P}\!\left(\hat{\mu} \ge \mu+\varepsilon\right)
\le
\exp\!\left(-\frac{n\varepsilon^2}{2\sigma^2}\right),
\qquad
\mathbb{P}\!\left(\hat{\mu} \le \mu-\varepsilon\right)
\le
\exp\!\left(-\frac{n\varepsilon^2}{2\sigma^2}\right),
\]
where
\[
\hat{\mu}=\frac{1}{n}\sum_{t=1}^n X_t.
\]
\end{lemma}

\subparagraph{Problem setup.}
Let the dimension satisfy $d \ge n+1$, where $n$ is the horizon. The unknown parameter is
\[
\theta^\star = e_1 \in \mathbb{R}^d,
\]
where $\{e_i\}_{i=1}^d$ denotes the standard basis. Rewards follow the stochastic linear model
\[
X_t = \langle \theta^\star, A_t \rangle + \eta_t,
\]
where $\eta_t$ are independent mean-zero $\sigma$-subgaussian random variables.

At each round $t=1,\dots,n$, the learner is offered two actions,
\[
g:=e_1,
\qquad
b_t:=L e_{t+1},
\]
where $L>1$ is a fixed constant. Their expected rewards are
\[
\langle \theta^\star,g\rangle = 1,
\qquad
\langle \theta^\star,b_t\rangle = 0.
\]
Both LinUCB variants below use the rule
\[
\mathrm{UCB}_t(a)
=
a^\top \hat{\theta}_{t-1}
+
\beta \|a\|_{V_{t-1}^{-1}},
\]
with
\[
V_{t-1}=V_0+\sum_{s<t} A_sA_s^\top,
\qquad
\hat{\theta}_{t-1}
=
V_{t-1}^{-1}\sum_{s<t} A_s X_s.
\]

\subparagraph{Algorithm 1: a constant policy.}
Consider a policy $\pi_{\mathrm{const}}$ that selects an action by maximizing $\langle \hat{\theta}_{\mathrm{const}}, a\rangle$ for some fixed parameter estimate $\hat{\theta}_{\mathrm{const}}\in\mathbb{R}^d$ that is never updated. For concreteness, take
\[
\hat{\theta}_{\mathrm{const}}
=
\frac{1}{L}\sum_{j=2}^{n+1} e_j,
\]
so that
\[
\langle \hat{\theta}_{\mathrm{const}},g\rangle = 0,
\qquad
\langle \hat{\theta}_{\mathrm{const}},b_t\rangle = 1
\quad\text{for every } t\le n.
\]
Then $\pi_{\mathrm{const}}$ selects $b_t$ at every round $t$. Since each pull of $b_t$ yields expected reward $0$ while $g$ has expected reward $1$, the cumulative pseudo-regret of $\pi_{\mathrm{const}}$ on the horizon $n$ is exactly $n$. This proves part (1) of Proposition~\ref{prop:three-way}.

The failure mode is structural rather than statistical: $\pi_{\mathrm{const}}$ never queries direction $e_1$, so no amount of data can refine $\hat{\theta}_{\mathrm{const}}$ in the reward-relevant direction. The policy has no learning mechanism. We will see, however, that the trajectory it produces is precisely what suppresses the misplaced optimism of LinUCB in the bad directions.

\subparagraph{Algorithm 2: LinUCB with isotropic initialization.}
The first LinUCB variant uses
\[
V_0=\lambda I_d,
\qquad
\hat{\theta}_0=0.
\]
At $t=1$,
\[
\|g\|_{V_0^{-1}}^2=\frac{1}{\lambda},
\qquad
\|b_1\|_{V_0^{-1}}^2=\frac{L^2}{\lambda}.
\]
Since $\hat{\theta}_0=0$, we have
\[
\mathrm{UCB}_1(b_1)>\mathrm{UCB}_1(g),
\]
so the algorithm selects $b_1$.

Now suppose inductively that the algorithm has selected $b_1,\dots,b_{t-1}$ in the first $t-1$ rounds. Then
\[
V_{t-1}
=
\lambda I_d+\sum_{s=1}^{t-1} b_s b_s^\top
=
\lambda I_d + L^2\sum_{j=2}^{t} e_j e_j^\top.
\]
Hence $V_{t-1}$ is diagonal in the standard basis. In particular, for the fresh direction $e_{t+1}$,
\[
e_{t+1}^\top V_{t-1}^{-1} e_{t+1}=\frac{1}{\lambda},
\qquad
\|b_t\|_{V_{t-1}^{-1}}^2=\frac{L^2}{\lambda},
\qquad
\|g\|_{V_{t-1}^{-1}}^2=\frac{1}{\lambda}.
\]

The key point is that although the observed rewards on $b_1,\dots,b_{t-1}$ contain noise, the estimate
\[
\hat{\theta}_{t-1}
=
V_{t-1}^{-1}\sum_{s=1}^{t-1} b_s X_s
\]
lies entirely in $\mathrm{span}\{e_2,\dots,e_t\}$. Therefore,
\[
g^\top \hat{\theta}_{t-1}=0,
\qquad
b_t^\top \hat{\theta}_{t-1}=0,
\]
because both $g=e_1$ and the fresh arm $b_t=L e_{t+1}$ are orthogonal to $\mathrm{span}\{e_2,\dots,e_t\}$.

Thus, even in the stochastic-reward setting, the comparison between $g$ and $b_t$ is determined purely by the confidence terms:
\[
\mathrm{UCB}_t(g)=\frac{\beta}{\sqrt{\lambda}},
\qquad
\mathrm{UCB}_t(b_t)=\frac{\beta L}{\sqrt{\lambda}}.
\]
Since $L>1$, we have $\mathrm{UCB}_t(b_t)>\mathrm{UCB}_t(g)$ for all $t\le n$. Hence the isotropically initialized algorithm selects $b_t$ at every round and, like $\pi_{\mathrm{const}}$, incurs cumulative pseudo-regret $n$. This proves part (2) of Proposition~\ref{prop:three-way}.

\subparagraph{Algorithm 3: LinUCB with an anisotropic warm start.}
We now define an initialization that is realized by exactly the trajectory $\pi_{\mathrm{const}}$ would produce. Fix an integer $m\ge 1$, and let
\[
V_0'
=
\lambda I_d + mL^2 \sum_{j=2}^{n+1} e_j e_j^\top,
\qquad
\hat{\theta}_0'=0.
\]
This corresponds to having observed each bad direction $L e_j$, $j=2,\dots,n+1$, exactly $m$ times before the online interaction begins --- precisely the kind of data $\pi_{\mathrm{const}}$ generates if rolled out for $mn$ rounds on a sequence of action sets containing each bad direction. Thus $V_0'$ is a genuine design matrix obtainable from a collection of outer products produced by a constant logging policy.

At $t=1$,
\[
\|g\|_{(V_0')^{-1}}^2=\frac{1}{\lambda},
\qquad
\|b_1\|_{(V_0')^{-1}}^2
=
L^2\cdot \frac{1}{\lambda+mL^2}
=
\frac{L^2}{\lambda+mL^2}.
\]
Hence, if
\[
m > \frac{\lambda(L^2-1)}{L^2},
\]
then $\frac{1}{\lambda}>\frac{L^2}{\lambda+mL^2}$, and therefore
\[
\mathrm{UCB}_1'(g)>\mathrm{UCB}_1'(b_1).
\]
So the warm-started algorithm selects $g$ at round $1$.

Suppose now that the algorithm continues to select $g$ for the first $s$ rounds. Then
\[
V_s'
=
\lambda I_d
+
mL^2 \sum_{j=2}^{n+1} e_j e_j^\top
+
s\, e_1 e_1^\top.
\]
Since each reward observed on $g$ has mean $1$,
\[
\sum_{i=1}^s X_i = s + \sum_{i=1}^s \eta_i,
\]
and therefore
\[
g^\top \hat{\theta}_s'
=
e_1^\top (V_s')^{-1}\sum_{i=1}^s g X_i
=
\frac{1}{\lambda+s}\sum_{i=1}^s X_i.
\]
Applying Lemma~\ref{lem:subgaussian} with $\mu=1$, we obtain that with probability at least $1-\delta$, simultaneously for all $1\le s\le n$,
\[
\frac{1}{s}\sum_{i=1}^s X_i
\ge
1-\sqrt{\frac{2\sigma^2\log(2n/\delta)}{s}}.
\]
On this event,
\[
g^\top \hat{\theta}_s'
\ge
\frac{s}{\lambda+s}
\left(
1-\sqrt{\frac{2\sigma^2\log(2n/\delta)}{s}}
\right).
\]

For every fresh bad arm $b_{s+1}=Le_{s+2}$,
\[
b_{s+1}^\top \hat{\theta}_s' = 0,
\]
since the online trajectory has only sampled the direction $e_1$, and the warm-start contributes only to the design matrix, not to the response vector. Moreover,
\[
\|b_{s+1}\|_{(V_s')^{-1}}^2
=
\frac{L^2}{\lambda+mL^2},
\qquad
\|g\|_{(V_s')^{-1}}^2
=
\frac{1}{\lambda+s}.
\]

On the same high-probability event,
\[
\mathrm{UCB}_{s+1}'(g)
\ge
\frac{s}{\lambda+s}
\left(
1-\sqrt{\frac{2\sigma^2\log(2n/\delta)}{s}}
\right)
+
\frac{\beta}{\sqrt{\lambda+s}},
\]
while
\[
\mathrm{UCB}_{s+1}'(b_{s+1})
=
\frac{\beta L}{\sqrt{\lambda+mL^2}}.
\]
Consequently, if $m$ is chosen large enough so that for every $1\le s\le n$,
\begin{equation}
\label{eq:m-condition}
\frac{s}{\lambda+s}
\left(
1-\sqrt{\frac{2\sigma^2\log(2n/\delta)}{s}}
\right)
+
\frac{\beta}{\sqrt{\lambda+s}}
>
\frac{\beta L}{\sqrt{\lambda+mL^2}},
\end{equation}
then, with probability at least $1-\delta$, the warm-started algorithm selects $g$ at every round, and the cumulative pseudo-regret over the $n$ online rounds is $0$ on this event. This proves part (3) of Proposition~\ref{prop:three-way}; condition~\eqref{eq:m-condition} is the explicit ``$m$ sufficiently large'' requirement referenced in the proposition statement. The condition is satisfiable by a finite $m$ whenever the left-hand side is positive for every $s$, which holds for any $\beta\ge\sigma\sqrt{2\log(2n/\delta)}$ (in particular for the radius of \cite{abbasi2011improved}), since then $\beta/\sqrt{\lambda+s}\ge\sigma\sqrt{2\log(2n/\delta)}\,\sqrt s/(\lambda+s)$.\hfill$\square$


\begin{remark}
The warm-started algorithm in the proof uses $\widehat\theta_0'=0$, i.e., it
carries the offline Gram matrix but not the offline response sum. Because
every offline observation is of an arm with expected reward $0$, the offline
response sum has mean zero in each bad direction and carrying it would only
add a mean-zero perturbation to $\widehat\theta$ in directions orthogonal to
$g$; the comparison between $g$ and $b_{s+1}$ is unchanged up to the noise
term already absorbed in the choice of $m$. We keep the convention of the
original construction.
\end{remark}

\subsection{Proof of Proposition~\ref{prop:one-pass}}
\label{app:proof-one-pass}

\begin{proof}
Work on the instance of Section~\ref{app:proof-three-way} with horizon $T$:
$\theta^\star=e_1$, $\mathcal A_t=\{g,b_t\}$ with $g=e_1$ and
$b_t=Le_{t+1}$. Fix an arbitrary expert $\Exp$ and a handoff time $\tau$.
Whatever the expert plays during rounds $1,\dots,\tau$, and whatever LinUCB
plays afterwards, the set of arms played before round $t$ is contained in
$\{g\}\cup\{b_s:s<t\}$, all of which are orthogonal to $e_{t+1}$. Hence for
every $t$, the design matrix $V_t=\lambda I+\sum_{s<t}x_{s,A_s}x_{s,A_s}^\top$
is block diagonal with the block on $e_{t+1}$ equal to $\lambda$, and
$\widehat\theta_t=V_t^{-1}b_t$ lies in $\mathrm{span}\{e_1,\dots,e_t\}$. Therefore
\[
\|b_t\|_{V_t^{-1}}^2=\frac{L^2}{\lambda},\qquad \langle\widehat\theta_t,b_t\rangle=0,
\qquad\text{for every }t\text{ and every trajectory.}
\]
So $\mathrm{UCB}_t(b_t)=\beta L/\sqrt\lambda$ exactly, for all $t$.

Now bound $\mathrm{UCB}_t(g)$. Let $s_t$ be the number of pulls of $g$ before
round $t$ and $X_1,\dots,X_{s_t}$ the rewards observed on those pulls; since
$V_t$ is diagonal with entry $\lambda+s_t$ on $e_1$ and the rewards on bad
arms enter $b_t$ only in directions orthogonal to $e_1$,
\[
\langle\widehat\theta_t,g\rangle=\frac{1}{\lambda+s_t}\sum_{i=1}^{s_t}X_i,
\qquad
\|g\|_{V_t^{-1}}^2=\frac{1}{\lambda+s_t}\le\frac1\lambda .
\]
By Lemma~\ref{lem:subgaussian} and a union bound over $s\in\{1,\dots,T\}$,
with probability at least $1-\delta$ we have
$\sum_{i\le s}X_i\le s+\sqrt{2\sigma^2s\log(2T/\delta)}$ for all $s\le T$
simultaneously. On that event,
\[
\langle\widehat\theta_t,g\rangle
\le\frac{s_t}{\lambda+s_t}+\frac{\sqrt{2\sigma^2s_t\log(2T/\delta)}}{\lambda+s_t}
\le 1+\frac{\sigma\sqrt{2\log(2T/\delta)}}{\sqrt{\lambda+s_t}}
\le 1+\kappa ,
\]
using $\sqrt{s_t}/(\lambda+s_t)\le1/\sqrt{\lambda+s_t}\le1/\sqrt\lambda$.
Hence $\mathrm{UCB}_t(g)\le1+\kappa+\beta/\sqrt\lambda$. Under the hypothesis
$\beta(L-1)/\sqrt\lambda>1+\kappa$ this is strictly less than
$\beta L/\sqrt\lambda=\mathrm{UCB}_t(b_t)$, so LinUCB, whenever it acts, plays
$b_t$. The argument did not use anything about rounds $1,\dots,\tau$ beyond
the fact that they cannot produce coverage of $e_{t+1}$ for $t>\tau$, so it
holds for every expert and every $\tau$. Since $\Delta_{t,b_t}=1$,
$R_{[\tau+1,T]}(\Hyb_\tau)=T-\tau$. Cold LinUCB plays $b_t$ on every round by
part~(2) of Proposition~\ref{prop:three-way} (no condition on $\beta$ is
needed there because $\widehat\theta_t\perp g$ on its own trajectory), so
$R_{[\tau+1,T]}(\Cold)=T-\tau$ and $V_\tau=0$.
\end{proof}

\subsection{Proof of Lemma~\ref{lem:width-limited}}
\label{app:proof-width}

\begin{proof}
Fix a round $t$ at which LinUCB acts, and write $x^\star=x_{t,a_t^\star}$,
$x=x_{t,A_t}$, $\wid(u)=\|u\|_{V_t^{-1}}$. On $\mathcal E_\delta$, Cauchy--Schwarz
gives $|\langle\widehat\theta_t-\theta^\star,u\rangle|\le\beta_t\wid(u)$ for
every $u$. By the selection rule,
$\langle\widehat\theta_t,x\rangle+\beta_t\wid(x)\ge\langle\widehat\theta_t,x^\star\rangle+\beta_t\wid(x^\star)\ge\langle\theta^\star,x^\star\rangle$,
and $\langle\widehat\theta_t,x\rangle\le\langle\theta^\star,x\rangle+\beta_t\wid(x)$.
Combining, $\Delta_{t,A_t}=\langle\theta^\star,x^\star-x\rangle\le2\beta_t\wid(x)$.
Thus $\Delta_{t,A_t}=\min\{\Delta_{t,A_t},2\beta_t\wid(x_{t,A_t})\}\le\max_a\min\{\Delta_{t,a},2\beta_t\wid(x_{t,a})\}=\Phi_t(V_t)$.
Monotonicity: $V\preceq V'$ implies $V'^{-1}\preceq V^{-1}$, so every width is
non-increasing and so is each term in the maximum.
\end{proof}

\subsection{Proof of Lemma~\ref{lem:loewner}}
\label{app:proof-loewner}

\begin{proof}
$\lambda I+M\succeq\lambda I$ gives $(\lambda I+M)^{-1}\preceq(\lambda I)^{-1}$,
which is the first claim. For the second, $\lambda I+M\succeq\lambda I+mP$, so
$x^\top(\lambda I+M)^{-1}x\le x^\top(\lambda I+mP)^{-1}x$, and for
$x\in\mathrm{Im}(P)$ the right-hand side is $\|x\|_2^2/(\lambda+m)$ because
$(\lambda I+mP)^{-1}$ acts as $(\lambda+m)^{-1}$ on $\mathrm{Im}(P)$. For the
third, if $x\perp\mathrm{Im}(M)$ then $x$ is an eigenvector of $\lambda I+M$
with eigenvalue $\lambda$, so $x^\top(\lambda I+M)^{-1}x=\|x\|_2^2/\lambda$.
\end{proof}

\subsection{Proof of Lemma~\ref{lem:transfer}}
\label{app:proof-transfer}

\begin{proof}
By the Sherman--Morrison formula,
$(V+xx^\top)^{-1}=V^{-1}-\dfrac{V^{-1}xx^\top V^{-1}}{1+x^\top V^{-1}x}$. Multiplying
on both sides by $z$,
\[
z^\top(V+xx^\top)^{-1}z
= z^\top V^{-1}z-\frac{(z^\top V^{-1}x)^2}{1+\|x\|^2_{V^{-1}}}
=\|z\|^2_{V^{-1}}\Bigl(1-\frac{(z^\top V^{-1}x)^2}{(1+\|x\|^2_{V^{-1}})\|z\|^2_{V^{-1}}}\Bigr).
\]
The bracket is at least $1-\dfrac{\|x\|^2_{V^{-1}}\|z\|^2_{V^{-1}}}{(1+\|x\|^2_{V^{-1}})\|z\|^2_{V^{-1}}}>0$
by Cauchy--Schwarz in the $V^{-1}$ inner product, so $\rho_V\in[0,1)$. At
$V=\lambda I$, $z^\top V^{-1}x=\langle z,x\rangle/\lambda$,
$\|x\|^2_{V^{-1}}=\|x\|_2^2/\lambda$, $\|z\|^2_{V^{-1}}=\|z\|_2^2/\lambda$, and
substituting gives
$\rho_{\lambda I}(x,z)=\dfrac{\langle x,z\rangle^2}{\|x\|_2^2\|z\|_2^2}\cdot\dfrac{\|x\|_2^2}{\lambda+\|x\|_2^2}$.
\end{proof}

\subsection{Proof of Proposition~\ref{thm:handoff}}
\label{app:proof-handoff}

\begin{proof}
The acquisition term is the definition of $R_{[1,\tau]}(\Hyb_\tau)=R_{[1,\tau]}(\Exp)$.
For the LinUCB phase we give two bounds and take the minimum.

\emph{First bound.} By Lemma~\ref{lem:width-limited} applied at each
$t>\tau$, $\sum_{t>\tau}\Delta_{t,A_t}\le\sum_{t>\tau}\Phi_t(V_t)$.

\emph{Second bound.} From the proof of Lemma~\ref{lem:width-limited},
$\Delta_{t,A_t}\le2\beta_t\|x_{t,A_t}\|_{V_t^{-1}}\le2\beta_T\|x_{t,A_t}\|_{V_t^{-1}}$
for $t>\tau$ (the radius $\beta_t$ is non-decreasing in $t$). Cauchy--Schwarz
over $t=\tau+1,\dots,T$ gives
\[
\sum_{t>\tau}\Delta_{t,A_t}\le2\beta_T\sqrt{(T-\tau)\sum_{t>\tau}\|x_{t,A_t}\|^2_{V_t^{-1}}} .
\]
By the elliptical potential lemma \cite[Lemma~11]{abbasi2011improved},
$\sum_{t>\tau}\min\{1,\|x_{t,A_t}\|^2_{V_t^{-1}}\}\le2\log\det V_{T+1}/\det V_{\tau+1}$;
since $\|x_{t,A_t}\|^2_{V_t^{-1}}\le L^2/\lambda$, the truncation at $1$ costs at
most a factor $\max\{1,L^2/\lambda\}$ on the sum of squares, which is the
factor $\max\{1,L/\sqrt\lambda\}$ in \eqref{eq:handoff-bound} after the square
root; with $L^2\le\lambda$ nothing is lost. This gives
the second term of \eqref{eq:handoff-bound}.

\emph{Determinant ratio.} Let $V_{\tau+1}=\sum_j\sigma_ju_ju_j^\top$ be an
eigendecomposition and $G=\sum_{t>\tau}x_{t,A_t}x_{t,A_t}^\top$. Then
$\det(V_{\tau+1}+G)/\det V_{\tau+1}=\det(I+V_{\tau+1}^{-1/2}GV_{\tau+1}^{-1/2})\le\prod_j(1+u_j^\top Gu_j/\sigma_j)$
by Hadamard's inequality for positive definite matrices, and
$u_j^\top Gu_j=\sum_{t>\tau}\langle u_j,x_{t,A_t}\rangle^2\le L^2(T-\tau)$.
Taking logarithms gives the stated bound. If all played actions lie in an
$r$-dimensional subspace $\mathrm{Im}(P)$ then $G$ is supported there and only
the $r$ eigenvalues of $V_{\tau+1}$ restricted to $\mathrm{Im}(P)$ appear, which
is how the factor $r$ in the remark after the theorem arises; if in addition
$V_{\tau+1}\succeq(\lambda+\nu\tau)P$ on that subspace, each $\sigma_j\ge\lambda+\nu\tau$.

The cold bound is the case $\tau=0$, $V_1=\lambda I$, $\sigma_j=\lambda$. The
ratio $\Psi$ in the text is the quotient of the two isotropic bounds and is
stated only to show that it does not tend to zero.
\end{proof}

\begin{remark}[Validity of the confidence event with an online warm start]
The hybrid's LinUCB phase uses $(V_{\tau+1},b_{\tau+1})$ built from the
expert-routed rounds together with the observed rewards on those rounds. The
self-normalized bound of \cite{abbasi2011improved} applies to any sequence of
(feature, reward) pairs in which the feature at round $s$ is measurable with
respect to the history before $s$, whoever chose it, so the event
$\mathcal E_\delta$ covers the hybrid's trajectory. If instead the warm start
carried $M_\tau$ but dropped $c_\tau$ (``passive'' updates), the estimate
$\widehat\theta_t$ would be biased toward zero in the directions of $M_\tau$
and the event would not be guaranteed; the Router in our experiments updates
both, and Algorithm~\ref{alg:confidence-gated} reflects this.
\end{remark}

\subsection{Proofs for Section~\ref{sec:two-channels}}
\label{app:proof-two-channels}

\begin{proof}[Proof of Proposition~\ref{prop:two-channel}]
Write $x^\star=x_{t,a_t^\star}$, $x=x_{t,A_t}$, $\wid(u)=\|u\|_{V_t^{-1}}$.
Cauchy--Schwarz in the $V_t$ inner product gives
$|\langle\widehat\theta_t-\theta^\star,u\rangle|\le\rho_t\wid(u)$ for every $u$;
this is the inequality in the proof of Lemma~\ref{lem:width-limited} with the
realized error in place of its bound. The selection rule gives
\[
\langle\widehat\theta_t,x\rangle+\beta_t\wid(x)
\ \ge\ \langle\widehat\theta_t,x^\star\rangle+\beta_t\wid(x^\star)
\ \ge\ \langle\theta^\star,x^\star\rangle+(\beta_t-\rho_t)\wid(x^\star)
\ \ge\ \langle\theta^\star,x^\star\rangle,
\]
using $\rho_t\le\beta_t$, and
$\langle\widehat\theta_t,x\rangle\le\langle\theta^\star,x\rangle+\rho_t\wid(x)$.
Subtracting, $\Delta_{t,A_t}\le(\beta_t+\rho_t)\wid(x)$, and taking the maximum
over $a$ of the minimum with the gap gives \eqref{eq:phi-two}. Monotonicity in
$V$ is Lemma~\ref{lem:width-limited}; in $\rho$ it is immediate. Finally
$\min\{\Delta,(\beta+\rho)w\}\ge\min\{\Delta,\beta w\}\ge\tfrac12\min\{\Delta,2\beta w\}$.
\end{proof}

\begin{proof}[Proof of Proposition~\ref{prop:response-sum}]
$c=\sum_sr_sx_s=\sum_s(\langle\theta^\star,x_s\rangle+\eta_s)x_s
=\bigl(\sum_sx_sx_s^\top\bigr)\theta^\star+\sum_s\eta_sx_s=M\theta^\star+\xi$.
Conditionally on the features, $\xi$ is Gaussian with mean zero and covariance
$\sigma^2\sum_sx_sx_s^\top=\sigma^2M$. Then
$\widehat\theta_1-\theta^\star=V_1^{-1}(M\theta^\star+\xi)-\theta^\star
=V_1^{-1}(M-V_1)\theta^\star+V_1^{-1}\xi=-\lambda V_1^{-1}\theta^\star+V_1^{-1}\xi$.
Both $V_1=\lambda I+M$ and the law of $\xi$ are functions of $M$ alone, so the
conditional law of $(V_1,b_1)$ given $M$ is the same whatever policy produced
the actions. The LinUCB phase is a deterministic function of $(V_1,b_1)$ and of
the future contexts and noise, which are independent of the warm start; hence
the tail trajectories agree in law, and so do $V_\tau$ and $B_\tau$.
\end{proof}

\begin{proof}[Proof of Proposition~\ref{prop:transfer}]
$c=\sum_s(\langle\theta^\star,x_s\rangle+\varepsilon(x_s)+\eta_s)x_s
=M\theta^\star+\sum_s\varepsilon(x_s)x_s+\xi$, and
$V_1^{-1}M\theta^\star=\theta^\star-\lambda V_1^{-1}\theta^\star$.
\end{proof}

\begin{proof}[Proof of Proposition~\ref{prop:transfer-regret}]
Fix $\widetilde\theta$ and write $\wid(u)=\|u\|_{V_t^{-1}}$,
$x^\star=x_{t,a^\star_t}$, $x=x_{t,A_t}$. Cauchy--Schwarz gives
$|\langle\widehat\theta_t-\widetilde\theta,u\rangle|\le\widetilde\rho_t\wid(u)$.
By the definition of $\widetilde\varepsilon$,
\[
\bar r(x^\star)=\langle\widetilde\theta,x^\star\rangle+\widetilde\varepsilon(x^\star)
\le\langle\widehat\theta_t,x^\star\rangle+\widetilde\rho_t\wid(x^\star)+\widetilde\varepsilon(x^\star).
\]
The selection rule gives
$\langle\widehat\theta_t,x^\star\rangle\le\langle\widehat\theta_t,x\rangle+\beta_t\wid(x)-\beta_t\wid(x^\star)$,
and
$\langle\widehat\theta_t,x\rangle\le\langle\widetilde\theta,x\rangle+\widetilde\rho_t\wid(x)
=\bar r(x)-\widetilde\varepsilon(x)+\widetilde\rho_t\wid(x)$.
Chaining the three displays,
\[
\bar r(x^\star)-\bar r(x)\le(\beta_t+\widetilde\rho_t)\wid(x)
+(\widetilde\rho_t-\beta_t)\wid(x^\star)
+\widetilde\varepsilon(x^\star)-\widetilde\varepsilon(x),
\]
and $(\widetilde\rho_t-\beta_t)\wid(x^\star)\le(\widetilde\rho_t-\beta_t)_+\wid(x^\star)$.
\end{proof}

\subsection{Proof of Theorem~\ref{thm:general-price}}
\label{app:proof-general-price}

Throughout, work on the confidence event $\mathcal E_\delta$ of cold LinUCB
(Section~\ref{sec:width}), so that Lemma~\ref{lem:width-limited} applies at
every round, and write $\beta=\beta_T\ge\beta_t$.

\paragraph{(a)}
Fix $x\in S$ and suppose cold LinUCB plays $x$ at round $t$ after $m$ earlier
pulls of $x$. From the proof of Lemma~\ref{lem:width-limited}, the played arm
satisfies $\Delta(x)=\Delta_{t,A_t}\le2\beta_t\|x\|_{V_t^{-1}}\le2\beta\|x\|_{V_t^{-1}}$.
Since $V_t\succeq\lambda I+m\,xx^\top$ (the $m$ earlier pulls are among the
summands of $V_t$), Lemma~\ref{lem:loewner} and the Sherman--Morrison identity
give $\|x\|^2_{V_t^{-1}}\le\|x\|^2_{(\lambda I+mxx^\top)^{-1}}=\|x\|_2^2/(\lambda+m\|x\|_2^2)$.
Hence $\Delta(x)^2(\lambda+m\|x\|_2^2)\le4\beta^2\|x\|_2^2$, i.e.\
$\lambda+m\|x\|_2^2\le c_x$. If $c_x<\lambda$ this is impossible already for
$m=0$, so $x$ is never played. If $c_x\ge\lambda$, every pull of $x$ has
$m\le(c_x-\lambda)/\|x\|_2^2$, so the number of pulls is at most
$(c_x-\lambda)/\|x\|_2^2+1=m_x$. The number of pulls is also at most $N_x$,
the number of rounds at which $x$ is offered. Each pull costs exactly
$\Delta(x)$ because $x^\star\in\mathcal A_t$, and pulls of optimal arms cost
nothing. Summing over $x\in S$ gives (a).

\paragraph{(b)}
Fix $u_x\in U\subseteq U_{\mathrm{cert}}(t_0)$ and write
$C(u_x):=\sum_{t\le t_0}\langle x^\pi_t,u_x\rangle^2$ for the coverage of
$u_x$ after round $t_0$. Since $x^\top V^\pi_{t_0+1}x=\lambda\|x\|_2^2+\sum_{t\le t_0}\langle x^\pi_t,x\rangle^2=\|x\|_2^2(\lambda+C(u_x))$,
Cauchy--Schwarz in the $V^\pi_{t_0+1}$ inner product (as in the proof of
Lemma~\ref{lem:no-free-coverage}) gives
$\|x\|^2_{(V^\pi_{t_0+1})^{-1}}\ge\|x\|_2^4/(x^\top V^\pi_{t_0+1}x)=\|x\|_2^2/(\lambda+C(u_x))$.
Certification, $4\beta^2\|x\|^2_{(V^\pi_{t_0+1})^{-1}}<\Delta(x)^2$, therefore forces
$\lambda+C(u_x)>4\beta^2\|x\|_2^2/\Delta(x)^2=c_x$, i.e.\ $C(u_x)>c_x-\lambda$;
and trivially $C(u_x)\ge0\ge(c_x-\lambda)_+$ when $c_x<\lambda$. This is the
first inequality. For the second, exchange the order of summation:
$\sum_{u_x\in U}C(u_x)=\sum_{t\le t_0}\mathrm{cov}_U(x^\pi_t)$. Split the
rounds $t\le t_0$ into those with $\Delta(x^\pi_t)=0$, which contribute
$C^{\mathrm{free}}_U(t_0)$ by definition, and those with
$\Delta(x^\pi_t)>0$. At a round of the second kind with
$\mathrm{cov}_U(x^\pi_t)>0$, the definition of $\pi_U(t_0)$ as a
minimum gives $\mathrm{cov}_U(x^\pi_t)\le\Delta(x^\pi_t)/\pi_U(t_0)$;
at a round with $\mathrm{cov}_U(x^\pi_t)=0$ the same inequality holds
trivially. Summing, the rounds of the second kind contribute at most
$\sum_{t\le t_0}\Delta(x^\pi_t)/\pi_U(t_0)=R_{[1,t_0]}(\pi)/\pi_U(t_0)$
(if there are no rounds of the second kind with positive coverage, the
contribution is $0$, which matches the convention $1/\pi_U(t_0)=0$).

\paragraph{(c)}
Split the sum in (a) into $u_x\in U$ and $u_x\notin U$. For $u_x\in U$, use
$\min\{m_x,N_x\}\le m_x\le(c_x-\lambda)_+/\|x\|_2^2+1$ to get
$\Delta(x)\min\{m_x,N_x\}\le p(x)(c_x-\lambda)_++\Delta(x)\le p_U(c_x-\lambda)_++\Delta(x)$.
Summing over $u_x\in U$ and applying (b),
\[
\sum_{u_x\in U}\Delta(x)\min\{m_x,N_x\}\ \le\ p_U\Bigl(\frac{R_{[1,t_0]}(\pi)}{\pi_U(t_0)}+C^{\mathrm{free}}_U(t_0)\Bigr)+\sum_{u_x\in U}\Delta(x),
\]
and $p_U/\pi_U(t_0)=P_U$. The terms with $u_x\notin U$ are carried
unchanged. Finally $R_T(\pi)\ge R_{[1,t_0]}(\pi)$ because pseudo-regret is
nonnegative on every round; dividing gives \eqref{eq:general-price}. For
$\pi=\Hyb_\tau$ and $t_0=\tau$, $R_{[1,\tau]}(\Hyb_\tau)=R_{[1,\tau]}(\Exp)$
because the hybrid plays the expert's actions on $[1,\tau]$.
\qed

\paragraph{The template as a check.}
On the instance of Theorem~\ref{thm:price} take $U=\{e_j:j\in J\}$, all of
which are certified at $\tau_0$ by (iv). Then $p(b_j)=1/L^2$, so $p_U=1/L^2$;
the expert plays only probes in the window, $\mathrm{cov}_U(p_j)=c^2$ and
$\Delta(p_j)=\eps$, so $\pi_U(\tau_0)=\eps/c^2$ and $P_U=c^2/(\eps L^2)$;
$C^{\mathrm{free}}_U=0$ because $g\perp U$; $\sum_U\Delta=K$; and the
uncertified arms are the probes, for which
$\Delta(p_j)\min\{m_{p_j},N_{p_j}\}\le\eps n$. The theorem gives
$R_T(\Cold)/R_T(\Hyb_{\tau_0})\le c^2/(\eps L^2)+(K+\eps Kn)/(\eps Kn)=c^2/(\eps L^2)+1+1/(\eps n)$,
consistent with the upper bound $1+(8\beta^2+1)/(\eps n)\le1+c^2/(\eps L^2)$
proved directly. The same computation applies verbatim to
Proposition~\ref{prop:cat-tail-saving}, where the distractors have gap $0$
and contribute nothing.

\section{Proofs for the Catalog of Section~\ref{sec:catalog}}
\label{app:catalog}

Throughout this appendix the conventions \eqref{eq:conventions} are in force,
$\alpha:=1-\eps$, $m_0:=\lfloor\lambda/30\rfloor$, and ``LinUCB acts at round
$t$'' means that the action at $t$ is chosen by the LinUCB rule (every round
for $\Cold$; rounds $t>\tau$ for $\Hyb_\tau$). Ties in the LinUCB rule are
broken arbitrarily; every comparison we use is strict. For a policy $\pi$ and
a round $t$ we write $s(t)$, $m_j(t)$, $n_j(t)$ for the numbers of pulls of
$g$, $b$ in direction $e_j$, and $p$ in direction $e_j$ before round $t$
(all of them counted from the start, including expert-routed rounds), and we
drop the argument $t$ when no confusion arises. In the variant of
Proposition~\ref{prop:cat-gap} the roles of $g$ and $p$ are played by
$g_\eta$ and $p_{t,\eta}$ and $\alpha:=(1-\eps)\eta$; we prove everything for
general $\eta\in(0,1]$ and set $\eta=1$ outside that proposition.

\subsection{Events and their probabilities}

\begin{lemma}[Good events]
\label{lem:cat-events}
Fix a policy $\pi$ on any instance of the template. Let
$\mathcal E_{\mathrm{conf}}(\pi)$ be the event that
$\|\widehat\theta_t-\theta^\star\|_{V_t}\le\beta$ for every $t\in[T]$, where
$(V_t,b_t)$ are formed from all of $\pi$'s data before $t$. Let
$\mathcal E_{\mathrm{noise}}(\pi)$ be the event that for every direction
$j\in\{2,\dots,d\}$ and every $m\in[T]$, the sum $S_j^{(m)}$ of the noises on
the first $m$ pulls of $b$ in direction $e_j$ by $\pi$ satisfies
$|S_j^{(m)}|\le\sigma\sqrt{2m\ell_1}$. Then
$\Prb(\mathcal E_{\mathrm{conf}}(\pi))\ge1-\delta$ and
$\Prb(\mathcal E_{\mathrm{noise}}(\pi))\ge1-\delta$. Moreover, under
\eqref{eq:conventions}, $\beta\ge4\sigma\sqrt{2\ell_1}$.
\end{lemma}

\begin{proof}
All arm vectors of the template have Euclidean norm at most $L$: $\|g_\eta\|=\eta\le1$,
$\|b_t\|=L$, $\|p_{t,\eta}\|^2=\alpha^2+c^2\le2\le L^2$. Also $\|\theta^\star\|=1$.
By Theorem~2 of \cite{abbasi2011improved}, with probability at least
$1-\delta$, simultaneously for all $t$,
$\|\widehat\theta_t-\theta^\star\|_{V_t}\le\sigma\sqrt{2\log\bigl(\det(V_t)^{1/2}\det(\lambda I)^{-1/2}/\delta\bigr)}+\sqrt\lambda$.
Since $V_t\preceq(\lambda+tL^2)I$, $\det(V_t)^{1/2}\det(\lambda I)^{-1/2}\le(1+tL^2/\lambda)^{d/2}$,
so the radius is at most
$\sigma\sqrt{d\log(1+TL^2/\lambda)+2\log(1/\delta)}+\sqrt\lambda\le\sigma\sqrt{d\,\ell_0}+\sqrt\lambda\le2\sqrt\lambda=\beta$,
using $\lambda\ge1$ and $\lambda\ge\sigma^2d\ell_0$. This proves the first claim.
The theorem of \cite{abbasi2011improved} only requires the noise to be
conditionally sub-Gaussian given the past and the current action, which holds
in the reward-table model, and it applies to the data of any policy, not only
to LinUCB's own.

For the second claim, in the reward-table model the noises on successive pulls
of $b$ in direction $e_j$ form an i.i.d.\ $\sigma$-sub-Gaussian sequence
$(\eta_{j,i})_{i\ge1}$ regardless of the policy, and $S_j^{(m)}=\sum_{i\le m}\eta_{j,i}$.
By Hoeffding's inequality for sub-Gaussian sums,
$\Prb(|S_j^{(m)}|>\sigma\sqrt{2m\ell_1})\le2e^{-\ell_1}=\delta/(dT)$. A union
bound over $d-1<d$ directions and $T$ values of $m$ gives probability at most
$\delta$ for the complement. Finally $\beta=2\sqrt\lambda\ge2\sqrt{8\sigma^2\ell_1}=4\sigma\sqrt{2\ell_1}$.
\end{proof}

\begin{lemma}[Confidence consequences]
\label{lem:cat-conf}
On $\mathcal E_{\mathrm{conf}}(\pi)$, for every $t$ and every $x\in\R^d$,
$|\langle\widehat\theta_t-\theta^\star,x\rangle|\le\beta\|x\|_{V_t^{-1}}$.
Hence the index $U_t(x):=\langle\widehat\theta_t,x\rangle+\beta\|x\|_{V_t^{-1}}$
satisfies $\mu(x)\le U_t(x)\le\mu(x)+2\beta\|x\|_{V_t^{-1}}$, where
$\mu(x)=\langle\theta^\star,x\rangle$.
\end{lemma}

\begin{proof}
Cauchy--Schwarz in the $V_t$ inner product:
$|\langle\widehat\theta_t-\theta^\star,x\rangle|\le\|\widehat\theta_t-\theta^\star\|_{V_t}\|x\|_{V_t^{-1}}\le\beta\|x\|_{V_t^{-1}}$.
\end{proof}

\subsection{The state of LinUCB on the template}

Every action vector of the template is $g_\eta$, some $b$ in direction $e_j$,
or some $p_\eta$ in direction $e_j$, and the probe is the only one that couples
two coordinates. The design matrix therefore has an arrowhead structure whose
inverse is explicit.

\begin{lemma}[Arrowhead inverse]
\label{lem:arrowhead}
Let $V=\lambda I+Se_1e_1^\top+\sum_{j=2}^{d}\bigl(M_je_je_j^\top+n_jq_jq_j^\top\bigr)$
with $S,M_j,n_j\ge0$ and $q_j=\alpha e_1+ce_j$. Put
\[
D_j:=\lambda+M_j+n_jc^2,\qquad u_j:=\alpha cn_j,
\]
\[
Q:=\lambda+S+\alpha^2\sum_{j\ge2}n_j-\sum_{j\ge2}\frac{u_j^2}{D_j}
=\lambda+S+\alpha^2\sum_{j\ge2}\frac{n_j(\lambda+M_j)}{D_j}.
\]
Then $Q\ge\lambda$ and
\begin{align*}
(V^{-1})_{11}&=\frac1Q,&
(V^{-1})_{1j}&=-\frac{u_j}{D_jQ},\\
(V^{-1})_{jj}&=\frac1{D_j}+\frac{u_j^2}{D_j^2Q},&
(V^{-1})_{jj'}&=\frac{u_ju_{j'}}{D_jD_{j'}Q}\ \ (j\ne j'),
\end{align*}
for $j,j'\ge2$. In particular, if $n_j=0$ then
$(V^{-1})_{jj}=1/D_j$ and $(V^{-1})_{1j}=(V^{-1})_{jj'}=0$ for all $j'\ne j$; and
if $n_j\ge1$ and $M_j\ge1$ then $(V^{-1})_{jj}\le 1/D_j+1/M_j$.
\end{lemma}

\begin{proof}
Write $V=\begin{pmatrix}a&u^\top\\ u&D\end{pmatrix}$ with
$a=\lambda+S+\alpha^2\sum_jn_j$, $u=(u_2,\dots,u_d)^\top$, and
$D=\mathrm{diag}(D_2,\dots,D_d)$; this is the matrix in the statement because
$q_jq_j^\top=\alpha^2e_1e_1^\top+\alpha c(e_1e_j^\top+e_je_1^\top)+c^2e_je_j^\top$.
The Schur complement of $D$ is $a-u^\top D^{-1}u=Q$, and the block inverse
formula gives
$(V^{-1})_{11}=Q^{-1}$, $(V^{-1})_{1,\cdot}=-Q^{-1}u^\top D^{-1}$, and
$(V^{-1})_{\cdot,\cdot}=D^{-1}+D^{-1}uu^\top D^{-1}Q^{-1}$ on the lower block,
which are the displayed entries. The second expression for $Q$ follows from
$\alpha^2n_j-u_j^2/D_j=\alpha^2n_j(D_j-c^2n_j)/D_j=\alpha^2n_j(\lambda+M_j)/D_j$,
and it shows $Q\ge\lambda$. If $n_j=0$ then $u_j=0$ and the special case
follows. For the last claim, $Q\ge\alpha^2n_j(\lambda+M_j)/D_j\ge\alpha^2n_jM_j/D_j$,
so $u_j^2/(D_j^2Q)\le\alpha^2c^2n_j^2D_j/(D_j^2\alpha^2n_jM_j)=c^2n_j/(D_jM_j)\le1/M_j$,
using $c^2n_j\le D_j$.
\end{proof}

\begin{lemma}[State of a policy on the template]
\label{lem:cat-state}
Let $\pi$ be any policy on an instance of the template and $t$ a round. Then
$V_t$ is the matrix of Lemma~\ref{lem:arrowhead} with $S=\eta^2s$, $M_j=m_jL^2$,
$n_j$ the probe counts, and $q_j=p_\eta$'s direction vector. For a direction
$j\ge2$ with $n_j=0$:
\begin{enumerate}[itemsep=1pt,topsep=1pt,leftmargin=*]
\item[(a)] $\|e_j\|^2_{V_t^{-1}}=1/(\lambda+m_jL^2)$ and
$\widehat\theta_{t,j}=LS_j^{(m_j)}/(\lambda+m_jL^2)$, where $S_j^{(m_j)}$ is the
sum of the noises on the $m_j$ pulls of $b$ in direction $e_j$;
\item[(b)] $\|g_\eta\|^2_{V_t^{-1}}=\eta^2/Q\le1/\lambda$ and
$\|p_\eta\|^2_{V_t^{-1}}=\alpha^2/Q+c^2/(\lambda+m_jL^2)\le2/\lambda$ for the probe
in direction $e_j$.
\end{enumerate}
For any direction with $m_j\ge1$, $\|e_j\|^2_{V_t^{-1}}\le2/(m_jL^2)$.
\end{lemma}

\begin{proof}
The form of $V_t$ is immediate from the definitions. For (a), Lemma~\ref{lem:arrowhead}
with $n_j=0$ gives $(V_t^{-1})_{jj}=1/D_j=1/(\lambda+m_jL^2)$ and that the $j$-th
row of $V_t^{-1}$ vanishes off the diagonal. The $j$-th coordinate of the
response vector $b_t$ is $L\sum r$ over the pulls of $b$ in direction $e_j$
plus $c\sum r$ over the pulls of $p$ in that direction; the second sum is
empty, and since $\mu(b)=0$ the first is $LS_j^{(m_j)}$. Hence
$\widehat\theta_{t,j}=(V_t^{-1})_{jj}LS_j^{(m_j)}$. For (b),
$\|g_\eta\|^2_{V_t^{-1}}=\eta^2(V_t^{-1})_{11}=\eta^2/Q$ and
$\|p_\eta\|^2_{V_t^{-1}}=\alpha^2(V_t^{-1})_{11}+2\alpha c(V_t^{-1})_{1j}+c^2(V_t^{-1})_{jj}
=\alpha^2/Q+c^2/D_j$, using $(V_t^{-1})_{1j}=0$; both bounds use $Q\ge\lambda$,
$D_j\ge\lambda$, $\alpha\le1$, $c\le1$, $\eta\le1$. The last claim is the last
claim of Lemma~\ref{lem:arrowhead} with $M_j=m_jL^2\ge L^2\ge1$ together with
$1/D_j\le1/(m_jL^2)$, and it holds trivially when $n_j=0$.
\end{proof}

\subsection{Three reusable facts}

\begin{lemma}[Optimism trap]
\label{lem:cat-trap}
Let $\pi$ be a policy on an instance of the template and suppose LinUCB acts
for $\pi$ at round $t$. On $\mathcal E_{\mathrm{conf}}(\pi)\cap\mathcal E_{\mathrm{noise}}(\pi)$,
if the direction $j=j(t)$ satisfies $n_j(t)=0$ and $m_j(t)<m_0$, then
\[
U_t(b_t)\ \ge\ \tfrac34\,\frac{\beta L}{\sqrt{\lambda+m_j(t)L^2}}\ >\ \max\{U_t(g_\eta),U_t(p_{t,\eta})\},
\]
where the probe term is omitted if the probe is not offered at $t$. In
particular LinUCB plays $b_t$.
\end{lemma}

\begin{proof}
Write $m=m_j(t)$ and $D=\lambda+mL^2$. By Lemma~\ref{lem:cat-state}(a),
$U_t(b_t)=L\widehat\theta_{t,j}+\beta L\|e_j\|_{V_t^{-1}}=L^2S_j^{(m)}/D+\beta L/\sqrt D$.
On $\mathcal E_{\mathrm{noise}}$, $|L^2S_j^{(m)}/D|\le L^2\sigma\sqrt{2m\ell_1}/D\le L\sigma\sqrt{2\ell_1}/\sqrt D$,
since $L\sqrt m\le\sqrt D$; by Lemma~\ref{lem:cat-events} this is at most
$(\beta/4)L/\sqrt D$. This gives the first inequality. For the second, by
Lemma~\ref{lem:cat-conf} and Lemma~\ref{lem:cat-state}(b),
$U_t(g_\eta)\le\eta+2\beta\eta/\sqrt Q\le1+2\beta/\sqrt\lambda$ and
$U_t(p_{t,\eta})\le(1-\eps)\eta+2\beta\sqrt{2/\lambda}\le1+2\sqrt2\,\beta/\sqrt\lambda$.
Since $\beta=2\sqrt\lambda$, $1\le\beta/(2\sqrt\lambda)$, so both indices are at
most $(\tfrac12+2\sqrt2)\beta/\sqrt\lambda$. It remains to check
$\tfrac34L/\sqrt D>(\tfrac12+2\sqrt2)/\sqrt\lambda$, i.e.\
$D<\lambda L^2\,(3/4)^2/(\tfrac12+2\sqrt2)^2$. Numerically
$(3/4)^2/(\tfrac12+2\sqrt2)^2=0.5625/11.0784>0.0507$. Since $m\le m_0-1<\lambda/30$,
$D<\lambda+\lambda L^2/30=\lambda L^2(1/L^2+1/30)\le\lambda L^2(1/64+1/30)<0.0491\,\lambda L^2$,
using $L\ge8$. The inequality holds with room to spare.
\end{proof}

\begin{lemma}[Bad-pull cap]
\label{lem:cat-cap}
Let $\pi$ be a policy on an instance of the template. On
$\mathcal E_{\mathrm{conf}}(\pi)$, at any round $t$ at which LinUCB acts for
$\pi$ and plays $b_t$ with direction $j$, $m_j(t)\le8\beta^2/\eta^2$; if
moreover $n_j(t)=0$, then $m_j(t)\le4\beta^2/\eta^2$. Consequently the number
of rounds at which LinUCB acting for $\pi$ plays $b$ in direction $e_j$ is at
most $8\beta^2/\eta^2+1$, and at most $4\beta^2/\eta^2+1$ if no probe in
direction $e_j$ is ever played by $\pi$.
\end{lemma}

\begin{proof}
By Lemma~\ref{lem:width-limited}, a bad arm with gap $\eta$ is played only if
$\eta\le2\beta L\|e_j\|_{V_t^{-1}}$, i.e.\ $\|e_j\|^2_{V_t^{-1}}\ge\eta^2/(4\beta^2L^2)$.
If $m_j(t)\ge1$, Lemma~\ref{lem:cat-state} gives $\|e_j\|^2_{V_t^{-1}}\le2/(m_j(t)L^2)$,
so $m_j(t)\le8\beta^2/\eta^2$; with $n_j(t)=0$ the sharper bound
$\|e_j\|^2_{V_t^{-1}}=1/(\lambda+m_j(t)L^2)\le1/(m_j(t)L^2)$ gives $m_j(t)\le4\beta^2/\eta^2$.
Since $m_j(\cdot)$ increases by one at each such round, the number of rounds
at which $b$ in direction $e_j$ is played with $m_j(t)\le M$ is at most $M+1$.
\end{proof}

\begin{lemma}[Expert routing]
\label{lem:expert-routes}
Let $\eps L<c$ and let $\Exp_w$ be the frozen linear policy with
$\widehat\theta_{\Exp}=e_1+w\sum_{j\ge2}e_j$ on the instance with $g_\eta,b_t,p_{t,\eta}$.
\begin{enumerate}[itemsep=1pt,topsep=1pt,leftmargin=*]
\item[(a)] If $w\in(\eta/L,(1-\eps)\eta/(L-c))$, a nonempty interval, then
$\Exp_w$ plays $p_{t,\eta}$ whenever the probe is offered and $b_t$ otherwise.
\item[(b)] If $w\in(\eps\eta/c,\eta/L)$, a nonempty interval, then $\Exp_w$ plays
$p_{t,\eta}$ whenever the probe is offered and $g_\eta$ otherwise.
\end{enumerate}
All three comparisons are strict.
\end{lemma}

\begin{proof}
The scores are $\langle\widehat\theta_{\Exp},g_\eta\rangle=\eta$,
$\langle\widehat\theta_{\Exp},b_t\rangle=wL$, and
$\langle\widehat\theta_{\Exp},p_{t,\eta}\rangle=(1-\eps)\eta+wc$. The probe beats
$g_\eta$ iff $wc>\eps\eta$ and beats $b_t$ iff $w<(1-\eps)\eta/(L-c)$; $b_t$
beats $g_\eta$ iff $w>\eta/L$. The interval in (a) is nonempty iff
$(1-\eps)\eta L>\eta(L-c)$, i.e.\ $c>\eps L$, and every $w$ in it satisfies
$w>\eta/L>\eps\eta/c$ because $c>\eps L$. The interval in (b) is nonempty iff
$\eps\eta/c<\eta/L$, again $c>\eps L$; for $w$ in it, $w<\eta/L<(1-\eps)\eta/(L-c)$.
\end{proof}

\begin{lemma}[Widths at handoff after $n$ probes in each direction]
\label{lem:cat-handoff}
Let $V=\lambda I+n\sum_{j=2}^{K+1}q_jq_j^\top$ with $q_j=\alpha e_1+ce_j$, $\alpha>0$,
$n\ge1$. Then for every $j\in\{2,\dots,K+1\}$,
\[
\|e_j\|^2_{V^{-1}}\ \le\ \frac1{nc^2}+\frac1{K\lambda},
\qquad\text{and}\qquad
Q=\lambda+\frac{\alpha^2Kn\lambda}{\lambda+nc^2}\ \ge\ \lambda\Bigl(1+\frac{\alpha^2K}{2c^2}\Bigr)\ \text{ if } nc^2\ge\lambda .
\]
\end{lemma}

\begin{proof}
Apply Lemma~\ref{lem:arrowhead} with $S=0$, $M_j=0$, $n_j=n$: $D_j=D:=\lambda+nc^2$,
$u_j=\alpha cn$, and $Q=\lambda+\alpha^2Kn\lambda/D$. Then
$(V^{-1})_{jj}=1/D+(\alpha cn)^2/(D^2Q)$. The first term is at most $1/(nc^2)$.
For the second, $Q\ge\alpha^2Kn\lambda/D$, so
$(\alpha cn)^2/(D^2Q)\le(\alpha cn)^2D/(D^2\alpha^2Kn\lambda)=c^2n/(DK\lambda)\le1/(K\lambda)$,
using $c^2n\le D$. The bound on $Q$ follows from $D\le2nc^2$ when $nc^2\ge\lambda$.
\end{proof}

\subsection{Proof of Lemma~\ref{lem:no-free-coverage}}

Fix $u$. At a round $t\in I$ with $\langle x_{t,\pi_t},u\rangle=0$ the
contribution to the coverage is zero. At a round with
$\langle x_{t,\pi_t},u\rangle\ne0$, the definition \eqref{eq:price} gives
$\Delta_{t,\pi_t}/\langle x_{t,\pi_t},u\rangle^2\ge\pi_t(u)\ge\pi(u)$, i.e.\
$\langle x_{t,\pi_t},u\rangle^2\le\Delta_{t,\pi_t}/\pi(u)$ (if $\pi(u)=0$ the bound
is vacuous). Summing over $I$ gives the first claim. For the second, let
$u$ be a unit vector and $V=V_{\tau+1}=\lambda I+M_\tau$. By Cauchy--Schwarz in
the $V$ inner product,
$1=(u^\top u)^2=(u^\top V^{1/2}V^{-1/2}u)^2\le(u^\top Vu)(u^\top V^{-1}u)$, so
$\|u\|^2_{V^{-1}}\ge1/(\lambda+u^\top M_\tau u)$, and
$u^\top M_\tau u=\sum_{t\le\tau}\langle x_{t,A_t},u\rangle^2\le R_{[1,\tau]}(\Exp)/\pi(u)$
by the first claim applied to the expert on $[1,\tau]$. \qed

\subsection{Proof of Theorem~\ref{thm:price}}

We work on $\mathcal E_{\mathrm{conf}}(\Cold)\cap\mathcal E_{\mathrm{noise}}(\Cold)\cap\mathcal E_{\mathrm{conf}}(\Hyb_{\tau_0})$,
which has probability at least $1-3\delta$ by Lemma~\ref{lem:cat-events}. Set
$\eta=1$.

\paragraph{(i) The expert.}
By Lemma~\ref{lem:expert-routes}(a), $\Exp_w$ plays $p_t$ at every $t\le\tau_0$
(gap $\eps$) and $b_t$ at every $t>\tau_0$ (gap $1$), so $R_T(\Exp)=\eps Kn+(T-Kn)$.

\paragraph{(ii) Cold LinUCB, lower bound.}
We claim that at every round $t\in[T]$ with $m_{j(t)}(t)<m_0$, cold LinUCB
plays $b_t$. We argue by induction on $t$. Suppose the claim holds for all
rounds before $t$ and let $j=j(t)$ with $m_j(t)<m_0$. If cold LinUCB had
played the probe in direction $e_j$ at some earlier round $t'$ (so $j(t')=j$),
then $m_j(t')\le m_j(t)<m_0$ and the claim at $t'$ says it played $b_{t'}$,
a contradiction; hence $n_j(t)=0$. Lemma~\ref{lem:cat-trap} now gives that
LinUCB plays $b_t$. (The base case $t=1$ is the same argument with no earlier
rounds.) Direction $e_j$ is offered at exactly $n$ rounds in $[1,Kn]$, and at
each of the first $\min\{n,m_0\}$ of them $m_j<m_0$, so cold LinUCB plays the
bad arm, paying $1$ each time. With $n\ge m_0$ this gives $R_{[1,Kn]}(\Cold)\ge Km_0$,
and $A_{\tau_0}=R_{[1,Kn]}(\Cold)-R_{[1,Kn]}(\Exp)\ge Km_0-\eps Kn$ by (i).

\paragraph{(iii) Cold LinUCB, upper bound.}
Cold LinUCB acts at every round. By Lemma~\ref{lem:cat-cap} it plays $b$ in
each of the $K$ directions at most $8\beta^2+1$ times over the horizon, each
costing $1$. Probes exist only at rounds $t\le Kn$ and cost $\eps$ each, so
probe pulls cost at most $\eps Kn$. Pulls of $g$ cost nothing. Hence
$R_T(\Cold)\le K(8\beta^2+1)+\eps Kn$.

\paragraph{(iv) The hybrid.}
During $t\le\tau_0$ the hybrid plays the expert's actions, which by (i) are
the probes, each direction exactly $n$ times. Hence
$V_{\tau_0+1}=\lambda I+n\sum_{j=2}^{K+1}p^{(j)}p^{(j)\top}$ with $p^{(j)}=\alpha e_1+ce_j$,
and Lemma~\ref{lem:cat-handoff} gives, for every bad direction,
$\|e_j\|^2_{V_{\tau_0+1}^{-1}}\le1/(nc^2)+1/(K\lambda)$. With $nc^2\ge36\lambda L^2=9\beta^2L^2$
and $K\lambda\ge36\lambda L^2=9\beta^2L^2$ this is at most $2/(9\beta^2L^2)$, so
$2\beta L\|e_j\|_{V_{\tau_0+1}^{-1}}\le2\sqrt2/3<1$. Now fix $t>\tau_0$. The
hybrid's design matrix satisfies $V_t\succeq V_{\tau_0+1}$, so by
Lemma~\ref{lem:loewner} $\|e_{j(t)}\|_{V_t^{-1}}\le\|e_{j(t)}\|_{V_{\tau_0+1}^{-1}}$,
and by Lemma~\ref{lem:cat-conf}, $U_t(b_t)\le0+2\beta L\|e_{j(t)}\|_{V_t^{-1}}<1\le U_t(g)$.
Since $\mathcal A_t=\{g,b_t\}$, LinUCB plays $g$, with gap $0$. Therefore
$R_{[\tau_0+1,T]}(\Hyb_{\tau_0})=0$, $R_T(\Hyb_{\tau_0})=R_{[1,\tau_0]}(\Exp)=\eps Kn$,
$V_{\tau_0}=R_{[Kn+1,T]}(\Cold)-0\ge0$, and $B_{\tau_0}=(T-Kn)-0$.

\paragraph{The ratios.}
From (ii), (iii), (iv): $R_T(\Cold)/R_T(\Hyb_{\tau_0})\ge Km_0/(\eps Kn)=m_0/(\eps n)$
and $\le(K(8\beta^2+1)+\eps Kn)/(\eps Kn)=1+(8\beta^2+1)/(\eps n)$; the
expert ratio is (i) over (iv). At $n=\lceil36\lambda L^2/c^2\rceil$, since
$36\lambda L^2/c^2\ge36\cdot60\cdot64>1$ we have $n\le37\lambda L^2/c^2$, and
$m_0\ge\lambda/30-1\ge\lambda/60$ because $\lambda\ge60$; hence
$m_0/(\eps n)\ge(\lambda/60)c^2/(37\eps\lambda L^2)=c^2/(2220\,\eps L^2)$. For the
upper bound, $8\beta^2+1=32\lambda+1$ and $n\ge36\lambda L^2/c^2$ give
$(8\beta^2+1)/(\eps n)\le(32\lambda+1)c^2/(36\eps\lambda L^2)\le c^2/(\eps L^2)$
since $\lambda\ge1$. \qed

\subsection{Proof of Corollary~\ref{cor:cat-dear}}

Let $\eps\ge c^2/L^2$ and $n\ge36\lambda L^2/c^2$, so that
$\eps n\ge(c^2/L^2)(36\lambda L^2/c^2)=36\lambda$. With $\beta=2\sqrt\lambda$
from \eqref{eq:conventions} we have $8\beta^2+1=32\lambda+1$, hence
\[
\frac{8\beta^2+1}{\eps n}\ \le\ \frac{32\lambda+1}{36\lambda}
\ =\ \frac89+\frac{1}{36\lambda}\ \le\ \frac89+\frac{1}{2160}\ <\ 1 ,
\]
using $\lambda\ge60$. The upper bound of \eqref{eq:price-ratio} then gives
$R_T(\Cold)/R_T(\Hyb_{\tau_0})\le1+(8\beta^2+1)/(\eps n)<2$. \qed

\subsection{Proof of Proposition~\ref{prop:cat-tail-saving}}
\label{app:proof-tail-saving}

\paragraph{Events.}
All arm vectors of $\mathcal J$ have norm at most $L$: $\|b'_{j'}\|_2^2=1+L'^2=145\le L^2$
since $L\ge13$, and the others are as in Lemma~\ref{lem:cat-events}. Let
$\mathcal E'_{\mathrm{conf}}(\Cold)$ be the event that
$\|\widehat\theta_t-\theta^\star\|_{V_t}\le\tfrac54\sqrt\lambda$ for every $t$
on cold LinUCB's trajectory. By Theorem~2 of \cite{abbasi2011improved} and
the computation in the proof of Lemma~\ref{lem:cat-events}, the
self-normalized radius is at most $\sigma\sqrt{d\ell_0}+\sqrt\lambda\le\tfrac14\sqrt\lambda+\sqrt\lambda$
under $\lambda\ge16\sigma^2d\ell_0$, so
$\Prb(\mathcal E'_{\mathrm{conf}}(\Cold))\ge1-\delta$, and
$\mathcal E'_{\mathrm{conf}}(\Cold)\subseteq\mathcal E_{\mathrm{conf}}(\Cold)$
since $\tfrac54\sqrt\lambda\le\beta=2\sqrt\lambda$. We work on
$\mathcal E'_{\mathrm{conf}}(\Cold)\cap\mathcal E_{\mathrm{noise}}(\Cold)\cap\mathcal E_{\mathrm{conf}}(\Hyb_{\tau_0})$,
of probability at least $1-3\delta$ (Lemma~\ref{lem:cat-events}; the noise
event concerns the pulls of $b$ in the directions $e_j$, $j\in J$, and its
proof is unchanged). On $\mathcal E'_{\mathrm{conf}}(\Cold)$, Cauchy--Schwarz
gives $|\langle\widehat\theta_t-\theta^\star,x\rangle|\le\tfrac54\sqrt\lambda\|x\|_{V_t^{-1}}=\tfrac58\beta\|x\|_{V_t^{-1}}$
for every $x$, hence for the index $U_t(x)=\langle\widehat\theta_t,x\rangle+\beta\|x\|_{V_t^{-1}}$,
\begin{equation}
\label{eq:sharp-index}
\mu(x)+\tfrac38\beta\|x\|_{V_t^{-1}}\ \le\ U_t(x)\ \le\ \mu(x)+\tfrac{13}8\beta\|x\|_{V_t^{-1}} .
\end{equation}

\paragraph{(i) The expert.}
The expert plays $p_{j(t)}$ (gap $\eps$) at every $t\le\tau_0$ and $b_{j(t)}$
(gap $1$) at every $t>\tau_0$, so $R_T(\Exp)=\eps Kn+(T-Kn)$.

\paragraph{(ii) Cold LinUCB in the window.}
We show by induction on $t\le\tau_0$ that cold LinUCB plays $b'_{j'(t)}$ at
round $t$. Suppose this holds before $t$; then no probe and no bad arm has
been played, so for every $j\in J$ no played vector has a component along
$e_j$, $V_te_j=\lambda e_j$, and $V_t^{-1}e_j=e_j/\lambda$. Three width bounds
follow. First, $\|g\|^2_{V_t^{-1}}=(V_t^{-1})_{11}\le1/\lambda$ because
$V_t\succeq\lambda I$. Second, with $j=j(t)$,
$\|p_j\|^2_{V_t^{-1}}=(1-\eps)^2(V_t^{-1})_{11}+2(1-\eps)c\,e_1^\top V_t^{-1}e_j+c^2(V_t^{-1})_{jj}
=(1-\eps)^2(V_t^{-1})_{11}+c^2/\lambda\le(1+c^2)/\lambda\le2/\lambda$, using
$e_1^\top V_t^{-1}e_j=e_1^\top e_j/\lambda=0$. Third, with $j'=j'(t)$ and $m$
the number of earlier pulls of $b'_{j'}$: the only played vector with a
component along $e_{j'}$ is $b'_{j'}$ itself, so
$e_{j'}^\top V_te_{j'}=\lambda+mL'^2$, and Cauchy--Schwarz in the $V_t$ inner
product gives
$\|b'_{j'}\|^2_{V_t^{-1}}\ge\langle b'_{j'},e_{j'}\rangle^2/(e_{j'}^\top V_te_{j'})=L'^2/(\lambda+mL'^2)$.
Direction $j'$ is offered at $\lceil\tau_0/K'\rceil$ rounds in the window, so
$m\le\tau_0/K'\le\lambda/L'^2$ by the condition on $K'$, and
$\|b'_{j'}\|^2_{V_t^{-1}}\ge L'^2/(2\lambda)$. Now apply
\eqref{eq:sharp-index} with $\beta=2\sqrt\lambda$:
\[
U_t(b'_{j'})\ \ge\ 1+\tfrac38\cdot2\sqrt\lambda\cdot\frac{L'}{\sqrt{2\lambda}}=1+\frac{3L'}{4\sqrt2}>7.3,
\qquad
U_t(g)\ \le\ 1+\tfrac{13}8\cdot2\sqrt\lambda\cdot\frac1{\sqrt\lambda}=4.25,
\qquad
U_t(p_j)\ \le\ 1-\eps+\tfrac{13}8\cdot2\sqrt\lambda\cdot\frac{\sqrt2}{\sqrt\lambda}<5.6,
\]
using $L'=12$. So cold LinUCB plays $b'_{j'(t)}$, whose gap is $0$. The base
case $t=1$ is the same computation with $m=0$. Hence $R_{[1,\tau_0]}(\Cold)=0$,
and at round $\tau_0+1$ every direction $e_j$, $j\in J$, is untouched:
$m_j=n_j=0$.

\paragraph{(ii) Cold LinUCB on the tail.}
For $t>\tau_0$ the action set is $\{g,b_{j(t)}\}$ and the directions $e_j$,
$j\in J$, are touched only by the bad arms, so for each of them $e_j$ is an
eigenvector of $V_t$ with eigenvalue $\lambda+m_j(t)L^2$, the $j$-th
coordinate of the response vector is $LS_j^{(m_j(t))}$, and the conclusions of
Lemma~\ref{lem:cat-state}(a) hold verbatim; Lemma~\ref{lem:cat-state}(b)'s
bound $U_t(g)\le1+2\beta/\sqrt\lambda$ holds on $\mathcal E_{\mathrm{conf}}(\Cold)$
by Lemma~\ref{lem:cat-conf} and $\|g\|_{V_t^{-1}}\le1/\sqrt\lambda$. These are
the only properties of the template that the proofs of
Lemmas~\ref{lem:cat-trap} and~\ref{lem:cat-cap} use for a direction with
$n_j=0$, so both lemmas apply to cold LinUCB on $\mathcal J$ at every tail
round (with no probe term). The induction of Theorem~\ref{thm:price}(ii) now
shows that cold LinUCB plays $b_t$ at every tail round with $m_{j(t)}(t)<m_0$;
since $T-\tau_0\ge Km_0$ and the tail directions cycle, each direction is
offered at least $m_0$ times on the tail, giving $R_{[\tau_0+1,T]}(\Cold)\ge Km_0$.
Lemma~\ref{lem:cat-cap} with $n_j=0$ bounds the pulls of $b$ in each
direction by $4\beta^2+1$, and $g$ costs nothing, so
$R_{[\tau_0+1,T]}(\Cold)\le K(4\beta^2+1)=K(16\lambda+1)$.

\paragraph{(iii) The hybrid.}
During $t\le\tau_0$ the hybrid plays the expert's actions, which by (i) are
the probes, each direction $j\in J$ exactly $n$ times; no distractor and no
bad arm is played. Hence $V_{\tau_0+1}=\lambda I+n\sum_{j\in J}p_jp_j^\top$,
which is the matrix of Theorem~\ref{thm:price}(iv), and the argument there
(Lemma~\ref{lem:cat-handoff}, Lemma~\ref{lem:loewner}, Lemma~\ref{lem:cat-conf}
on $\mathcal E_{\mathrm{conf}}(\Hyb_{\tau_0})$) shows that
$U_t(b_t)<1\le U_t(g)$ at every $t>\tau_0$, so the hybrid plays $g$ and
$R_{[\tau_0+1,T]}(\Hyb_{\tau_0})=0$. Therefore $R_T(\Hyb_{\tau_0})=\eps Kn$.

\paragraph{The terms and the ratios.}
$A_{\tau_0}=R_{[1,\tau_0]}(\Cold)-R_{[1,\tau_0]}(\Exp)=0-\eps Kn$;
$V_{\tau_0}=R_{[\tau_0+1,T]}(\Cold)-0\in[Km_0,K(16\lambda+1)]$;
$B_{\tau_0}=(T-Kn)-0$. Since $R_T(\Cold)=R_{[\tau_0+1,T]}(\Cold)$ and
$R_T(\Hyb_{\tau_0})=\eps Kn$, the ratio lies between $m_0/(\eps n)$ and
$(16\lambda+1)/(\eps n)$. At $n=\lceil36\lambda L^2/c^2\rceil$ the lower bound
is at least $c^2/(2220\,\eps L^2)$ exactly as in the proof of
Theorem~\ref{thm:price}, and $(16\lambda+1)/(\eps n)\le(16\lambda+1)c^2/(36\eps\lambda L^2)\le c^2/(2\eps L^2)$
since $\lambda\ge1$. \qed

\subsection{Proof of Proposition~\ref{prop:cat-mimic}}

Work on $\mathcal E_{\mathrm{conf}}(\Cold)\cap\mathcal E_{\mathrm{noise}}(\Cold)$
(probability $\ge1-2\delta$). There are no probes, so $n_j\equiv0$ for every
policy.

\emph{(i)} Let $t\le Km_0$ and $j=j(t)$. Direction $e_j$ was offered at
$\lceil t/K\rceil-1\le m_0-1$ earlier rounds, so $m_j(t)\le m_0-1<m_0$ for cold
LinUCB, and Lemma~\ref{lem:cat-trap} (with the probe term absent) says cold
LinUCB plays $b_t$. Hence cold LinUCB plays $b_t$ at every $t\le Km_0$, and
so does $\Exp$. Fix $\tau\le Km_0$. The two policies $\Cold$ and $\Hyb_\tau$
play the same action vectors in the same order during $[1,\tau]$; in the
reward-table model they therefore observe the same rewards, so their states
$(V_{\tau+1},b_{\tau+1})$ coincide. From round $\tau+1$ on both run LinUCB
from that state on the same action sets and the same reward table, so they
play the same actions at every round. Thus $R_I(\Cold)=R_I(\Hyb_\tau)$ for every
interval $I$, $A_\tau=R_{[1,\tau]}(\Cold)-R_{[1,\tau]}(\Exp)=\tau-\tau=0$, and $V_\tau=0$.

\emph{(ii)} For any $\tau$, $R_T(\Hyb_\tau)\ge R_{[1,\tau]}(\Exp)=\tau$. By
Lemma~\ref{lem:cat-cap} with $n_j\equiv0$, cold LinUCB plays each bad
direction at most $4\beta^2+1$ times, so $R_T(\Cold)\le K(4\beta^2+1)=K(16\lambda+1)$.
For $\tau>Km_0$, $R_T(\Cold)/R_T(\Hyb_\tau)\le K(16\lambda+1)/(Km_0)\le(16\lambda+1)/(\lambda/60)=960+60/\lambda\le961$. \qed

\subsection{Proof of Proposition~\ref{prop:cat-tail}}

Work on the event of Theorem~\ref{thm:price}. By Lemma~\ref{lem:expert-routes}(b),
$\Exp_w$ plays every probe during $[1,\tau_0]$ and $g$ afterwards, so
$R_T(\Exp)=\eps Kn$. The hybrid's routed rounds are the same probes as in
Theorem~\ref{thm:price}(iv), so its state at handoff and its subsequent play
are identical to those established there: $R_T(\Hyb_{\tau_0})=\eps Kn$ and
$R_{[\tau_0+1,T]}(\Hyb_{\tau_0})=0$. Hence $B_{\tau_0}=0-0=0$. Cold LinUCB does not
depend on the expert, so $A_{\tau_0}$ and $V_{\tau_0}$ are as in
Theorem~\ref{thm:price}. \qed

\subsection{Proof of Proposition~\ref{prop:cat-wrong}}

Work on $\mathcal E_{\mathrm{conf}}(\Cold)\cap\mathcal E_{\mathrm{conf}}(\Hyb_{\tau_0})\cap\mathcal E_{\mathrm{noise}}(\Hyb_{\tau_0})$
(probability $\ge1-3\delta$). Here $d=2K+1$; coordinates $2,\dots,K+1$ are the
bad directions and $K+2,\dots,2K+1$ the nuisance directions. The probe $p'_t$
couples $e_1$ with a nuisance coordinate only, so Lemma~\ref{lem:arrowhead}
applies with the nuisance coordinates playing the role of the probed $j$'s
and the bad coordinates having $n_j=0$ for every policy at every round.

\emph{Cold LinUCB.} Since $n_j\equiv0$ on bad coordinates,
Lemma~\ref{lem:cat-cap} gives at most $4\beta^2+1$ bad pulls in each of the $K$
directions, costing $1$ each; probe pulls occur only at $t\le Kn$ and cost
$\eps$ each; so $R_T(\Cold)\le K(4\beta^2+1)+\eps Kn$.

\emph{The hybrid.} During $[1,\tau_0]$ it plays the $Kn$ probes, paying $\eps Kn$.
At handoff, $V_{\tau_0+1}=\lambda I+n\sum_{i}q_iq_i^\top$ over the $K$ nuisance
directions with $q_i=\alpha e_1+ce_{K+1+i}$, and the bad coordinates are
untouched: by Lemma~\ref{lem:arrowhead} (with $M_j=n_j=0$ on them),
$\|e_j\|^2_{V_{\tau_0+1}^{-1}}=1/\lambda$ for every bad direction and the estimate
has zero $j$-th coordinate. By Lemma~\ref{lem:cat-handoff} applied to the
nuisance block, $Q\ge\lambda(1+\alpha^2K/(2c^2))$ once $nc^2\ge\lambda$, which holds
since $n\ge36\lambda L^2/c^2$. As $\alpha=1-\eps\ge1-1/L\ge7/8$ and $c\le1$,
$Q\ge\lambda(1+49K/128)\ge16\lambda=4\beta^2$ because $K\ge36L^2\ge2304$. Hence at
every $t>\tau_0$, since $Q$ only grows, $U_t(g)\le1+2\beta/\sqrt Q\le2$ by
Lemma~\ref{lem:cat-conf}. On the other hand, for $t>\tau_0$ and $j=j(t)$ with
$m_j(t)=m$, Lemma~\ref{lem:cat-state}(a) and $\mathcal E_{\mathrm{noise}}(\Hyb_{\tau_0})$
give, exactly as in the proof of Lemma~\ref{lem:cat-trap},
$U_t(b_t)\ge\tfrac34\beta L/\sqrt{\lambda+mL^2}$, which exceeds $2$ as soon as
$\lambda+mL^2<\tfrac9{64}\beta^2L^2=\tfrac9{16}\lambda L^2$, i.e.\ $m<\tfrac9{16}\lambda-\lambda/L^2$.
With $L\ge8$, $\tfrac9{16}\lambda-\lambda/L^2\ge\tfrac{35}{64}\lambda>\lambda/2$, so this holds
whenever $m\le\lfloor\lambda/2\rfloor-1$. Thus at every round $t>\tau_0$ at which
$m_{j(t)}(t)<\lfloor\lambda/2\rfloor$ the hybrid plays the bad arm. Each direction
is offered at least $\lfloor(T-\tau_0)/K\rfloor\ge\lfloor\lambda/2\rfloor$ times after
$\tau_0$, so each is pulled at least $\lfloor\lambda/2\rfloor$ times, and
$R_{[\tau_0+1,T]}(\Hyb_{\tau_0})\ge K\lfloor\lambda/2\rfloor$. Hence
$R_T(\Hyb_{\tau_0})\ge\eps Kn+K\lfloor\lambda/2\rfloor$.

\emph{The ratio and $V_{\tau_0}$.} Since $\lambda\ge60$, $\lfloor\lambda/2\rfloor\ge\lambda/4$,
and for $x\ge0$, $(a+x)/(b+x)\le\max\{a/b,1\}$; with $a=K(16\lambda+1)$,
$b=K\lambda/4$, $x=\eps Kn$ this gives $R_T(\Cold)/R_T(\Hyb_{\tau_0})\le64+4/\lambda\le65$.
Finally $V_{\tau_0}=R_{[\tau_0+1,T]}(\Cold)-R_{[\tau_0+1,T]}(\Hyb_{\tau_0})\le R_{[\tau_0+1,T]}(\Cold)-K\lfloor\lambda/2\rfloor$.
The expert's Gram matrix $n\sum_iq_iq_i^\top$ has range in
$\mathrm{span}\{e_1,e_{K+2},\dots,e_{2K+1}\}$, which is orthogonal to the span of
the bad directions, so the principal angle is $90^\circ$. \qed

\subsection{Proof of Proposition~\ref{prop:cat-cheap}}

On $\mathcal E_{\mathrm{conf}}(\Cold)$, Lemma~\ref{lem:cat-cap} with $n_j\equiv0$
gives $R_T(\Cold)\le K(4\beta^2+1)=K(16\lambda+1)$; since $R_T(\pi)\ge0$ for every
$\pi$, $(R_T(\Cold)-R_T(\pi))/T\le K(16\lambda+1)/T$, and with $\lambda=\sigma^2d\ell_0$
and $d=K+1$ this is $O(\sigma^2K^2\ell_0/T)$. For the second claim, the
certificate $2\beta L\|e_j\|_{V_{\tau+1}^{-1}}<1$ means $\|e_j\|^2_{V_{\tau+1}^{-1}}<1/(4\beta^2L^2)$,
while Lemma~\ref{lem:no-free-coverage} gives
$\|e_j\|^2_{V_{\tau+1}^{-1}}\ge1/(\lambda+R_{[1,\tau]}(\Exp)/\pi(e_j))$; combining,
$R_{[1,\tau]}(\Exp)/\pi(e_j)>4\beta^2L^2-\lambda$. On the template,
$\pi(e_j)=\min\{1/L^2,\eps/c^2\}$ while probes are offered; the acquisition
regret must exceed $\pi(e_j)(4\beta^2L^2-\lambda)$ for each of the $K$ directions
that has to be certified, and the contributions of different directions come
from disjoint sets of rounds because no action vector of the template has
nonzero components on two bad directions. \qed

\subsection{Proof of Proposition~\ref{prop:cat-gap}}

The proof of Theorem~\ref{thm:price} goes through with $\eta$ carried along;
we indicate the changes. Work on the same three events. (i)~By
Lemma~\ref{lem:expert-routes}(a), $R_T(\Exp)=\eps\eta Kn+\eta(T-Kn)$.
(ii)~Lemma~\ref{lem:cat-trap} was proved for general $\eta$, and the induction
is unchanged; each forced bad pull now costs $\eta$, so
$R_{[1,Kn]}(\Cold)\ge\eta Km_0$. (iii)~Lemma~\ref{lem:cat-cap} gives at most
$8\beta^2/\eta^2+1$ bad pulls in each direction at cost $\eta$, and probe pulls
cost $\eps\eta$, so $R_T(\Cold)\le\eta K(8\beta^2/\eta^2+1)+\eps\eta Kn$.
(iv)~At handoff, Lemma~\ref{lem:cat-handoff} with $\alpha=(1-\eps)\eta>0$ gives
$\|e_j\|^2_{V_{\tau_0+1}^{-1}}\le1/(nc^2)+1/(K\lambda)\le2\eta^2/(9\beta^2L^2)$ under
$nc^2\ge9\beta^2L^2/\eta^2$ and $K\lambda\ge9\beta^2L^2/\eta^2$, i.e.\ $n\ge36\lambda L^2/(\eta^2c^2)$
and $K\ge36L^2/\eta^2$. Then $U_t(b_t)\le2\beta L\|e_j\|_{V_t^{-1}}\le(2\sqrt2/3)\eta<\eta\le U_t(g_\eta)$
for all $t>\tau_0$ by Lemmas~\ref{lem:loewner} and~\ref{lem:cat-conf}, so the
hybrid plays $g_\eta$ and $R_T(\Hyb_{\tau_0})=\eps\eta Kn$. The ratio bounds
follow: the lower bound is $\eta Km_0/(\eps\eta Kn)=m_0/(\eps n)$, and at
$n=\lceil36\lambda L^2/(\eta^2c^2)\rceil\le37\lambda L^2/(\eta^2c^2)$ this is at least
$(\lambda/60)\eta^2c^2/(37\eps\lambda L^2)=\eta^2c^2/(2220\,\eps L^2)$; the upper bound
is $1+(8\beta^2/\eta^2+1)/(\eps n)\le1+(32\lambda+\eta^2)c^2/(36\eps\lambda L^2)\le1+c^2/(\eps L^2)$. \qed

\subsection{Proof of Corollary~\ref{cor:cat-oracle}}

On the event of Proposition~\ref{prop:one-pass}, the hybrid plays $b_t$ at
every $t>\tau$, so $R_{[\tau+1,T]}(\Hyb_\tau)=T-\tau$, while the routed rounds
cost nothing because the expert plays the optimal arm; hence
$R_T(\Hyb_\tau)=T-\tau$ and $R_T(\Exp)=0$. Cold LinUCB also plays $b_t$ at
every round by Proposition~\ref{prop:one-pass} with $\tau=0$, so
$R_{[1,\tau]}(\Cold)=\tau$ and $R_{[\tau+1,T]}(\Cold)=T-\tau$. Therefore
$A_\tau=\tau-0$, $V_\tau=(T-\tau)-(T-\tau)=0$, and $B_\tau=0-(T-\tau)$. \qed

\subsection{Remarks on the constants}

The constants $36$, $30$, $65$, $961$, $2220$ come from three places: the
factor $\tfrac34$ by which noise can reduce a bad arm's index below its bonus,
the factor $\tfrac12+2\sqrt2$ by which the good arm's or the probe's index can
exceed its mean at a cold start, and the choice $\beta=2\sqrt\lambda$. None of
them is tight. What matters in Theorem~\ref{thm:price} is the structure of
\eqref{eq:price-ratio}: the lower bound is (forced cold pulls per direction)
divided by (probe cost per direction), the upper bound is (maximal cold pulls
per direction) divided by the same, and both numerators are independent of
$\eps$, $c$, and $K$. The gap between $m_0\asymp\lambda$ and $8\beta^2\asymp\lambda$
is a constant factor, and that is why the multiplier is pinned to the price
ratio up to constants. The same structure with $\eta$ gives the bounds in
Proposition~\ref{prop:cat-gap}, where the looseness is a factor $\eta^2$ rather
than a constant because the forced pulls in Lemma~\ref{lem:cat-trap} are
counted before the good arm's index has come down to its mean; a sharper
count of cold LinUCB's pulls after the good arm is well estimated, as in the
proof of Proposition~\ref{prop:cat-wrong}, would remove it, but with probes
available to cold LinUCB the required control of the estimate in a coupled
coordinate is more delicate and we have not carried it out.


\section{Numerical Validation of the Catalog}
\label{app:synthetic}

Every claim in Section~\ref{sec:catalog} is a statement about an explicitly
constructed instance, so each can be checked by simulating that instance. We
do so for the arrowhead identity and for each catalog row whose proof asserts
a numerical bound. The simulations use the parameters named in the
corresponding proof, $10$ seeds, and no tuning; they are reported in
Table~\ref{tab:synthetic}. Their purpose is to catch an error in a constant or
an inequality direction, not to estimate anything: a proved claim that fails
here is a bug in the proof.

\begin{table}[h]
\centering
\caption{Simulated instances against their proved claims. ``Claim'' is what
the cited result asserts; ``simulated'' is what the instance produces. Row
numbers refer to Table~\ref{tab:catalog}. Each block uses the instance and
parameters named in the corresponding proof, so the absolute magnitudes are
not comparable across blocks; the first four row-5 entries are the instance at
$\eps=0.002$ and the last two the $\eps$-sweep instance.}
\label{tab:synthetic}
\small
\setlength{\tabcolsep}{4pt}
\begin{tabular}{@{}llll@{}}
\toprule
Result & Quantity & Claim & Simulated \\
\midrule
Lemma~\ref{lem:arrowhead} & max entrywise error, $200$ trials & $0$ & $<10^{-10}$ \\
\midrule
Prop.~\ref{prop:three-way} & cold regret & $T=2{,}000$ & $2{,}000$ \\
 & warm regret at $m=20$ & $0$ & $0$ (all seeds) \\
Prop.~\ref{prop:one-pass} & tail regret, good expert & $T-\tau$ & $T-\tau$ exactly \\
 & tail regret, bad expert & $T-\tau$ & $T-\tau$ exactly \\
\midrule
Prop.~\ref{prop:cat-mimic} (row 4) & hybrid vs.\ cold, $\tau\le Km_0$ & identical & identical \\
 & multiplier & $\le4\beta^2+1=961$ & $0.343$ \\
\midrule
Thm.~\ref{thm:price} (row 5) & expert regret, part (i) & $\eps Kn+(T{-}Kn)=201{,}244$ & $201{,}244$ \\
 & cold regret, part (iii) & $\le K(8\beta^2{+}1)+\eps Kn=62{,}716$ & $3{,}820$ \\
 & hybrid regret, part (iv) & $\eps Kn=1{,}244$ & $1{,}244$ \\
 & non-good plays after handoff & $0$ & $0$ \\
 & handoff certificate $2\beta L\|e_j\|_{V^{-1}}$ & $<1$ & $0.941$ \\
 & $(\text{ratio}-1)\cdot\eps$, $64\times$ sweep in $\eps$ & constant & $0.0044$--$0.0039$ \\
\midrule
Prop.~\ref{prop:cat-tail-saving} (row $5'$) & regret inside window & $0$ & $0.000$ \\
 & probes played & $0$ & $0$ \\
 & tail bad pulls per direction & $[m_0,\,4\beta^2{+}1]=[2,961]$ & $93$--$94$ \\
 & tail regret & $\ge Km_0=8$ & $374$ \\
\midrule
Prop.~\ref{prop:cat-wrong} (row 7) & multiplier & $\le65$ & $0.99$ \\
 & bad pulls per direction & $\Theta(\beta^2)$ & $214$ \\
\midrule
Cor.~\ref{cor:cat-dear} (row 8) & multiplier at $\eps\ge c^2/L^2$ & $\le2$ & $0.63,0.39,0.06$ \\
\bottomrule
\end{tabular}
\end{table}

Two entries carry more weight than the rest. The certificate row is the only
place where the construction could fail silently: Theorem~\ref{thm:price}
requires a $K$ for which the coverage condition holds, and the simulation
confirms that the $K$ the proof constructs does satisfy it
($0.941<1$) rather than merely being claimed to. The $\eps$-sweep row tests the
mechanism rather than the conclusion. Theorem~\ref{thm:price} predicts that
the multiplier exceeds one by an amount of order $1/\eps$; over a $64$-fold
range of $\eps$ the product $(\text{ratio}-1)\cdot\eps$ moves only between
$0.0044$ and $0.0039$, so the predicted scaling is what the instance exhibits,
not just the predicted sign.

The gaps between claim and simulation are all in the direction the proofs
allow, and are the constants discussed in
Appendix~\ref{app:catalog}, ``Remarks on the constants'': cold regret on the
row-5 template comes in at $3{,}820$ against a bound of $62{,}716$, and the
row-4 and row-7 multipliers at $0.343$ and $0.99$ against caps of $961$ and
$65$. In every case the instance does what the row says it does---rows 4, 7
and 8 kill the multiplier, row $5'$ collects its saving only on the tail with
an empty window---while the numerical bound is loose by the stated constant
factors.

\section{The Dominant Category on AmazonCat-13K}
\label{app:dominant}

This appendix gives the evidence behind the paragraph ``What the transfer
carries'' of Section~\ref{sec:amazon}. All numbers are on BERT-768d features
unless stated.

\paragraph{The label distribution.}
On the $50{,}000$ contexts of the online stream the four most frequent of the
$13{,}330$ categories are present on $29.8\%$, $17.0\%$, $11.3\%$ and $7.8\%$
of contexts (category identifiers $1471$, $7961$, $7892$, $7891$). A policy
that always predicts the first scores $29.77\%$; it uses the label
distribution, which an online learner is not given. Context-free bandits that
must discover it do not: over the same stream, $50{,}000$ rounds and seeds
$1$--$10$, UCB1 (one pull of every category first) scores $0.04\%$ and an
optimistic empirical-mean greedy rule $2.31\%$.

\paragraph{Where a run settles.}
We say a run \emph{settles} on a category if that category accounts for more
than $20\%$ of the learner's plays in the last fifth of the horizon. For the
dominant category the classification does not depend on the threshold: on the
runs of Table~\ref{tab:settle}, moving it anywhere between $10\%$ and $30\%$
changes the number of runs that settle on $1471$ by at most one on BERT
features and not at all on BGE features. The counts for the less frequent
categories are more sensitive on BERT, where runs that settle on $7892$
predict it on about a fifth of late rounds. Table~\ref{tab:settle}
reports where runs settle and how accurate the runs that settle on each
category are. On BERT features cold LinUCB settles on the dominant category
in $1$ of $20$ seeds; the two multiplying experts move it there in $17$ of
$20$. On BGE features cold LinUCB already settles there in $7$ of $20$ seeds,
which is why the expert has less to add; there the runs that settle on the
dominant category reach $41.2\%$, above the $29.8\%$ of always predicting
it, so the learner also uses the context. Knowing only where each run
settles predicts mean Router accuracy well: predicting a configuration's
accuracy as the frequency of each settled category times the mean accuracy
of runs settling there in the \emph{other} configurations gives $R^2=0.93$
over $25$ BERT configurations and $R^2=0.93$ over $7$ BGE configurations.
These are every configuration that routes fewer than half the rounds: on
BERT, the gate at $c\in\{1,10,10^2,10^3\}$, CORRAL, the two offline arms of
Table~\ref{tab:baselines} and the transfer-removed ablation of
Table~\ref{tab:two-channel}, for each of the three experts, plus cold
LinUCB; on BGE, the gate at $c\in\{0.01,0.1\}$ for each expert, plus cold
LinUCB.

\begin{table}[h]
\centering
\caption{Where runs settle on AmazonCat-13K (seeds $1$--$20$; BERT at $c=100$,
BGE at $c=0.1$). Each cell counts runs that settle on that category; the
percentage under each category is its frequency in the stream. The bottom
block gives the mean final accuracy of the runs, pooled over the four rows,
that settle on each category.}
\label{tab:settle}
\small
\setlength{\tabcolsep}{4pt}
\begin{tabular}{@{}llcccccc@{}}
\toprule
& & $1471$ & $7961$ & $7892$ & $7891$ & & \\
Features & Policy & ($29.8\%$) & ($17.0\%$) & ($11.3\%$) & ($7.8\%$) & none & Accuracy \\
\midrule
\multirow{4}{*}{BERT-768d} & Cold LinUCB & $1$ & $0$ & $9$ & $1$ & $9$ & $8.34\%$ \\
& Flan-T5-small & $3$ & $0$ & $14$ & $0$ & $3$ & $13.71\%$ \\
& Flan-T5-base & $17$ & $0$ & $3$ & $0$ & $0$ & $25.10\%$ \\
& Flan-T5-large & $17$ & $0$ & $3$ & $0$ & $0$ & $26.90\%$ \\
\midrule
\multirow{4}{*}{BGE-1024d} & Cold LinUCB & $7$ & $12$ & $0$ & $0$ & $1$ & $33.68\%$ \\
& Qwen-0.5B & $8$ & $9$ & $0$ & $2$ & $1$ & $33.23\%$ \\
& Qwen-3B & $17$ & $2$ & $0$ & $1$ & $0$ & $39.94\%$ \\
& Qwen-7B & $16$ & $3$ & $0$ & $1$ & $0$ & $39.67\%$ \\
\midrule
\multicolumn{2}{@{}l}{Accuracy of settled runs, BERT} & $28.08\%$ & --- & $11.61\%$ & $6.24\%$ & $5.93\%$ & \\
\multicolumn{2}{@{}l}{Accuracy of settled runs, BGE} & $41.21\%$ & $30.84\%$ & --- & $24.66\%$ & $25.96\%$ & \\
\bottomrule
\end{tabular}
\end{table}

\paragraph{A handful of routed answers decides it.}
Table~\ref{tab:dose} groups every Router run that routes fewer than half the
rounds by the number $k$ of routed rounds on which the expert proposed the
dominant category and was correct, and reports the fraction that settle on
it. On BERT features the fraction rises from $7\%$ at $k=0$ to $80\%$ at
$k=4$, and every one of the $54$ runs with $k\ge5$ settles there. On BGE
features the fraction at $k=0$ is already $43\%$, the rate at which the
learner finds the category by itself, and again every run with $k\ge5$
settles there. The pooled runs span experts, values of $c$ and, on BERT, the
CORRAL master, so the curve is a property of the learner and the data rather
than of one expert. The three BERT experts differ in exactly this count: in
the $c=100$ runs the mean $k$ is $0.20$ for Flan-T5-small, $3.70$ for
Flan-T5-base and $4.00$ for Flan-T5-large.

\begin{table}[h]
\centering
\caption{Runs that settle on the dominant category, by the number $k$ of
correct routed answers on it. BERT: $177$ runs ($c\in\{1,10,10^2,10^3\}$ and
CORRAL). BGE: $99$ runs ($c\in\{0.01,0.1\}$).}
\label{tab:dose}
\small
\begin{tabular}{@{}lcccccc@{}}
\toprule
$k$ & $0$ & $1$ & $2$ & $3$ & $4$ & $\ge5$ \\
\midrule
BERT-768d & $3/45$ ($7\%$) & $11/29$ ($38\%$) & $9/17$ ($53\%$) & $13/17$ ($76\%$) & $12/15$ ($80\%$) & $54/54$ ($100\%$) \\
BGE-1024d & $25/58$ ($43\%$) & $3/3$ & $5/9$ & $6/8$ & $8/9$ & $12/12$ ($100\%$) \\
\bottomrule
\end{tabular}
\end{table}

\paragraph{Controlled test.}
Table~\ref{tab:inject} changes the expert's answers on a fixed set of $15$
contexts and nothing else. The runs use an exact re-implementation of the
Router and of cold LinUCB, checked to make the same choice as the
experimental runner on every round of a $3{,}000$-round comparison, on
$10{,}000$ held-out contexts that do not appear in the online stream, with a
horizon of $10{,}000$ rounds and $20$ seeds; the seed fixes the order of the
contexts and is shared across conditions, so every comparison is paired.
Multipliers are relative to cold LinUCB on the same runs ($4.63\%$).
\emph{$+15$ on $1471$} replaces $15$ wrong answers of Flan-T5-small by the
dominant category on contexts that contain it. \emph{$+15$ random} gives the
same number of correct answers, on the same contexts where possible, on a
uniformly chosen true category other than $1471$. \emph{$+15$ on $7961$}
gives $15$ correct answers on the second most frequent category.
\emph{$-15$ on $1471$} replaces the $15$ correct dominant-category answers of
Flan-T5-large by wrong answers. The first three conditions raise held-out
accuracy by the same $0.15$ points; only the category differs.

\begin{table}[h]
\centering
\caption{Changing $15$ of the expert's $10{,}000$ held-out answers. Tests are
paired over $20$ seeds against the unmodified expert: a paired $t$-test on
Router accuracy and an exact McNemar test on whether the run settles on
$1471$.}
\label{tab:inject}
\small
\setlength{\tabcolsep}{4pt}
\begin{tabular}{@{}llccccc@{}}
\toprule
Expert & Change & Expert acc. & Router acc. & Mult. & Settle on $1471$ & $p$ (acc.\ / settle) \\
\midrule
\multirow{4}{*}{Flan-T5-small} & none & $5.34\%$ & $6.14\%$ & $1.33\times$ & $1/20$ & --- \\
& $+15$ random & $5.49\%$ & $7.63\%$ & $1.65\times$ & $2/20$ & $0.15$ / $1.00$ \\
& $+15$ on $7961$ & $5.49\%$ & $9.39\%$ & $2.03\times$ & $3/20$ & $0.010$ / $0.50$ \\
& $+15$ on $1471$ & $5.49\%$ & $\mathbf{12.10\%}$ & $\mathbf{2.62\times}$ & $\mathbf{8/20}$ & $\mathbf{0.002}$ / $\mathbf{0.016}$ \\
\midrule
\multirow{2}{*}{Flan-T5-large} & none & $5.32\%$ & $13.79\%$ & $2.98\times$ & $10/20$ & --- \\
& $-15$ on $1471$ & $5.17\%$ & $10.09\%$ & $2.18\times$ & $6/20$ & $0.067$ / $0.39$ \\
\bottomrule
\end{tabular}
\end{table}

At equal accuracy, the gain grows with the frequency of the category the
added answers point to: none for random categories, a significant gain for
the second most frequent category (where $5$ of $20$ runs settle on that
category instead), and the largest gain for the dominant one. Removing the
dominant-category answers from Flan-T5-large lowers its multiplier but not
significantly; its other answers still lead $6$ of $20$ runs to the
dominant category. Individual seeds are sensitive to any change in the
routed data, which can tip a run toward or away from a category, so we
compare rates over seeds rather than reading individual pairs.

\section{Confidence-Gated Meta-Algorithm: Full Description}
\label{app:meta-algorithm}

This appendix expands the brief description of
Section~\ref{sec:handoff} into the full round-by-round walkthrough,
states the deterministic confidence-gated routing rule, and provides
the pseudocode used in our experiments. The Router is the adaptive-$\tau$ version of the handoff hybrid
$\Hyb_\tau$ of Section~\ref{sec:preliminaries}: it chooses, round by round,
whether the expert or LinUCB acts, and feeds every observation to LinUCB.
Proposition~\ref{thm:handoff} applies to it with the acquisition term summed over
expert-routed rounds.

Rather than encoding prior knowledge through a single improved
initialization, we assume access to \emph{multiple frozen policy
estimators} $\{(\widehat\theta_i,V^{(i)})\}_{i=1}^m$. Each pair should
be thought of as a constant policy obtained offline:
$\widehat\theta_i$ is a fixed linear predictor and $V^{(i)}$ is the
corresponding design matrix that captures the geometry of the data on
which that predictor was trained. Our goal is not to compete with the
single best expert in hindsight, an objective closer in spirit to
online model-selection methods such as CORRAL
\cite{agarwal2017corralling,chen2024llmdecision}, which provide
robustness guarantees of the form ``perform nearly as well as the best
base algorithm.'' Here we want to collect the data-value term $V_\tau$ of
Proposition~\ref{prop:decomp} when it is available while paying as little
acquisition regret as possible, and to fall back to LinUCB otherwise.

\paragraph{Setting.}
At each round $t=1,\dots,T$, the learner observes a finite action set
$\mathcal A_t$ with feature vectors
$\{x_{t,a}\}_{a\in\mathcal A_t}\subseteq\mathbb R^d$, selects an action
$A_t\in\mathcal A_t$, and observes reward
$r_t = \langle\theta^\star,x_{t,A_t}\rangle + \eta_t$,
with $\eta_t$ conditionally mean-zero sub-Gaussian.
We are given $m$ frozen estimators
$\{(\widehat\theta_i,V^{(i)})\}_{i=1}^m$ and an online LinUCB learner
with state $(V_t,b_t,\widehat\theta_t)$, where
$\widehat\theta_t = V_t^{-1}b_t$. The frozen estimators are never
updated; only LinUCB is updated online.

\paragraph{Expert proposals and geometry-aware confidence.}
Each frozen expert $i\in[m]$ proposes the greedy action
$a_t^{(i)}\in\arg\max_{a\in\mathcal A_t}\langle\widehat\theta_i,x_{t,a}\rangle$.
To measure how trustworthy expert $i$ is on the current round, we
evaluate the uncertainty of its proposed action under its own design
geometry,
\[
u_t^{(i)} \;:=\; \bigl\|x_{t,a_t^{(i)}}\bigr\|_{(V^{(i)})^{-1}},
\qquad
\|x\|_{M^{-1}} \;:=\; \sqrt{x^\top M^{-1}x}.
\]
The interpretation is simple. If the proposed action lies in directions
that are well covered by the offline data underlying expert $i$, then
$u_t^{(i)}$ is small and the expert is more trustworthy; if the current
context points in directions poorly represented in $V^{(i)}$, then
$u_t^{(i)}$ is large and the expert should be trusted less. We treat
LinUCB itself as an additional candidate, indexed by $0$, with proposal
$a_t^{(0)}\in\arg\max_a\langle\widehat\theta_t,x_{t,a}\rangle+\alpha_t\|x_{t,a}\|_{V_t^{-1}}$
and uncertainty score $u_t^{(0)} := \|x_{t,a_t^{(0)}}\|_{V_t^{-1}}$.

\paragraph{Deterministic confidence-gated selection.}
The natural rule is to prefer the frozen expert whose proposed action
has the smallest uncertainty score, unless the online learner is
already at least as confident. Concretely, with
$i_t^\star\in\arg\min_{i\in[m]}u_t^{(i)}$,
\[
I_t \;=\;
\begin{cases}
i_t^\star, & \text{if }u_t^{(i_t^\star)}<u_t^{(0)},\\[2pt]
0, & \text{otherwise},
\end{cases}
\qquad A_t \;:=\; a_t^{(I_t)}.
\]
The master is entirely deterministic. Early in learning, LinUCB
typically has a large confidence radius, so a frozen expert may be
preferred; as data accumulate, $V_t$ grows and the LinUCB uncertainty
shrinks, and eventually the adaptive learner becomes the most confident
option and takes over automatically.

\paragraph{Optional permanent handover to LinUCB.}
To formalize the idea that LinUCB should eventually dominate, we include
a one-way switching criterion. With upper-confidence scores
$U_t^{(i)} = \langle\widehat\theta_i,x_{t,a_t^{(i)}}\rangle + \alpha_i\|x_{t,a_t^{(i)}}\|_{(V^{(i)})^{-1}}$
for each frozen expert and lower-confidence score
$L_t^{(0)} = \langle\widehat\theta_t,x_{t,a_t^{(0)}}\rangle - \alpha_t\|x_{t,a_t^{(0)}}\|_{V_t^{-1}}$
for LinUCB, if $L_t^{(0)} \ge \max_{i\in[m]} U_t^{(i)}$ for $K$
consecutive rounds we permanently hand control to LinUCB.

\begin{algorithm}[t]
\caption{Confidence-Gated Master over Frozen Policies and LinUCB}
\label{alg:confidence-gated}
\begin{algorithmic}[1]
\Require Frozen estimators $\{(\widehat{\theta}_i,V^{(i)})\}_{i=1}^m$, regularization $\lambda>0$, confidence parameters $\{\alpha_t\}_{t=1}^T$
\State Initialize LinUCB state:
\[
V_1=\lambda I_d,\qquad b_1=0,\qquad \widehat{\theta}_1=0
\]
\State Set $\textsc{handover}\gets \mathrm{False}$
\For{$t=1,\dots,T$}
    \State Observe action set $\mathcal{A}_t$
    \State Compute LinUCB proposal
    \[
    a_t^{(0)} \in \arg\max_{a\in\mathcal{A}_t}
    \left\{
    \langle \widehat{\theta}_t,x_{t,a}\rangle
    + \alpha_t \|x_{t,a}\|_{V_t^{-1}}
    \right\}
    \]
    \State Set $u_t^{(0)} \gets \|x_{t,a_t^{(0)}}\|_{V_t^{-1}}$
    \For{$i=1,\dots,m$}
        \State Compute frozen proposal
        \[
        a_t^{(i)} \in \arg\max_{a\in\mathcal{A}_t}
        \langle \widehat{\theta}_i,x_{t,a}\rangle
        \]
        \State Set $u_t^{(i)} \gets \|x_{t,a_t^{(i)}}\|_{(V^{(i)})^{-1}}$
    \EndFor
    \If{$\textsc{handover}=\mathrm{True}$}
        \State $A_t \gets a_t^{(0)}$
    \Else
        \State Choose $i_t^\star \in \arg\min_{i\in[m]} u_t^{(i)}$
        \If{$u_t^{(i_t^\star)} < u_t^{(0)}$}
            \State $A_t \gets a_t^{(i_t^\star)}$
        \Else
            \State $A_t \gets a_t^{(0)}$
        \EndIf
    \EndIf
    \State Observe reward $r_t$
    \State Update LinUCB:
    \[
    V_{t+1}\gets V_t + x_{t,A_t}x_{t,A_t}^\top,
    \qquad
    b_{t+1}\gets b_t + r_t x_{t,A_t},
    \qquad
    \widehat{\theta}_{t+1}\gets V_{t+1}^{-1}b_{t+1}
    \]
    \State Compute $U_t^{(i)}$ for all $i\in[m]$ and $L_t^{(0)}$
    \If{$L_t^{(0)} \ge \max_{i\in[m]} U_t^{(i)}$ for $K$ consecutive rounds}
        \State $\textsc{handover}\gets \mathrm{True}$
    \EndIf
\EndFor
\end{algorithmic}
\end{algorithm}

\paragraph{Distinction from randomized model selection.}
The deterministic, sample-by-sample geometric routing distinguishes
Algorithm~\ref{alg:confidence-gated} from CORRAL-style randomized
masters in two ways. First, it can exploit context-level information
about which estimator is reliable on \emph{this particular} action
set; second, because it is deterministic, the choice of which rounds to route
is a function of the current widths rather than of a sampled index. Both
rules feed the played action and the observed reward into $(V_t,b_t)$, so
both produce the online warm start $(M_\tau,c_\tau)$ whose value is
$V_\tau$ in Proposition~\ref{prop:decomp}; the decomposition itself does not
distinguish them, and any difference between them shows up in which rounds
are routed, i.e., in $A_\tau$ and in the composition of $M_\tau$.


\section{Pre-Runtime Estimates: Definitions and Consistency}
\label{app:diagnostics}

This appendix gives the formal definitions behind Section~\ref{sec:diagnostics}
and the elementary consistency statements that justify estimating them on a
held-out sample. Throughout, $\mathcal D=\{(\mathcal A_s,\{x_{s,a}\}_{a\in\mathcal A_s},\{r_{s,a}\}_{a\in\mathcal A_s})\}_{s=1}^{n}$
is a labeled held-out sample of $n$ contexts drawn from the same distribution
as the online stream, with a reward label for every candidate action (for
binary Precision@1 this is the ground-truth label), and $a_s^{\Exp}$ is the
expert's proposal on context $s$. Nothing in the online trajectory is used.

\subsection{Definitions}
\label{app:diag-defs}

\paragraph{Gap estimates.} Let $\widehat\theta_{\mathrm{or}}$ be the ridge
estimate fitted on all $(x_{s,a},r_{s,a})$ pairs of $\mathcal D$ with
regularization $\lambda$, and
$\widehat\Delta_{s,a}:=\max_{a'}\langle\widehat\theta_{\mathrm{or}},x_{s,a'}\rangle-\langle\widehat\theta_{\mathrm{or}},x_{s,a}\rangle$.
For binary reward with a single correct action one may instead use the label
directly, $\widehat\Delta_{s,a}=\ind\{a\ne a_s^\star\}$.

\paragraph{Expert Gram matrix.}
$\widehat M_{\Exp}:=\sum_{s=1}^nx_{s,a_s^{\Exp}}x_{s,a_s^{\Exp}}^\top$. For a target
number $\tau$ of routed rounds, $(\tau/n)\widehat M_{\Exp}$ is an unbiased
estimate of $\E[M_\tau]$ when the expert is followed for $\tau$ rounds on
contexts from the same distribution.

\paragraph{Width-limited regret under a design.} For a positive definite
$V$ and radius $\beta$,
\[
\widehat\Phi(V):=\frac1n\sum_{s=1}^n\max_{a\in\mathcal A_s}\min\bigl\{\widehat\Delta_{s,a},\,2\beta\|x_{s,a}\|_{V^{-1}}\bigr\}.
\]
We evaluate it at $V=\lambda I$ (cold start) and at
$V=\lambda I+(\tau/n)\widehat M_{\Exp}$ (after $\tau$ expert rounds). The
difference $\widehat\Phi(\lambda I)-\widehat\Phi(\lambda I+(\tau/n)\widehat M_{\Exp})$
is the estimated reduction, at a single future round, in the regret LinUCB
can still incur, attributable to the expert's data. It is zero whenever the
expert's proposals lie in directions orthogonal to every competitor whose
width exceeds its gap (Lemma~\ref{lem:loewner}), which is the failure mode the
principal-angle proxy detects coarsely.

\paragraph{Acquisition regret.}
$\widehat C_{\Exp}(\tau):=\tau\cdot\frac1n\sum_s\widehat\Delta_{s,a_s^{\Exp}}$.

\paragraph{Tail check.} $\widehat{\mathrm{acc}}_{\Exp}:=\frac1n\sum_s\ind\{a_s^{\Exp}=a_s^\star\}$
and $\widehat{\mathrm{acc}}_{\mathrm{or}}:=\frac1n\sum_s\ind\{\arg\max_a\langle\widehat\theta_{\mathrm{or}},x_{s,a}\rangle=a_s^\star\}$.
The learner's online accuracy is at most about $\widehat{\mathrm{acc}}_{\mathrm{or}}$
(up to the usual generalization gap), so $B_\tau>0$ requires
$\widehat{\mathrm{acc}}_{\Exp}<\widehat{\mathrm{acc}}_{\mathrm{or}}$ by a margin.

\paragraph{Net-value rule.} With $n_{\mathrm{fut}}=T-\tau$ future rounds,
\[
\widehat N(\tau):=n_{\mathrm{fut}}\bigl[\widehat\Phi(\lambda I)-\widehat\Phi(\lambda I+(\tau/n)\widehat M_{\Exp})\bigr]-\widehat C_{\Exp}(\tau).
\]
A handoff is predicted to multiply when $\widehat N(\tau)$ is a large
positive fraction of $n_{\mathrm{fut}}\widehat\Phi(\lambda I)$ for some
moderate $\tau$ and the tail check passes. We use ``large'' to mean at least
one third; this threshold is provisional. Two caveats bound what the rule
can certify. First, $\widehat\Phi(\lambda I)$ is the regret cold LinUCB could
incur at round one, not what it incurs on average over its own trajectory,
so a positive $\widehat N$ is not sufficient. Second, $\widehat N$ counts only
the saving on \emph{future} rounds and therefore misses the concurrent
channel of Proposition~\ref{prop:decomp}: on the instance of
Theorem~\ref{thm:price} the saving is carried by $A_\tau$ (cold pays for its
exploration inside the routed window at the higher price), the warm widths on
the bad arms stay near the gap, $\widehat\Phi$ barely moves, and
$\widehat N(\tau_0)<0$ even though the handoff provably multiplies. So a
negative $\widehat N$ argues against a tail-channel ($V_\tau$) multiplier but
does not exclude a concurrent-channel ($A_\tau$) one; on the instance of
Proposition~\ref{prop:cat-tail-saving}, where the same price ratio is
collected as $V_\tau$, $\widehat\Phi$ falls below the gap on every bad
direction and $\widehat N(\tau_0)$ is positive once the tail is long enough.
The estimate that covers both channels at once is a \emph{dry run}: replay the
held-out contexts (resampled to a horizon $T'$) with rewards drawn from the
labels or the ridge fit, simulate cold LinUCB and LinUCB warm-started from
$\lambda I+(\tau/n)\widehat M_{\Exp}$, and compare
$\widehat R_{T'}(\Cold)$ against
$\widehat C_{\Exp}(\tau)+\widehat R_{[\tau+1,T']}(\mathrm{warm})$. The dry run
uses no LLM calls beyond the held-out proposals already collected and no
online rounds; its cost is one bandit simulation per candidate $\tau$.

\paragraph{What these estimators do on our data, and why.}
Neither predicts the size of the observed multiplier, and
Section~\ref{sec:two-channels} says why. On all twelve configurations
$\widehat\Phi(\lambda I)$ saturates at $1.0000$: with $13{,}330$ actions in
$768$ dimensions the width term $2\beta\|x\|_{(\lambda I)^{-1}}$ exceeds the
unit gap for every held-out context, so the clipping in \eqref{eq:phi} is
active everywhere, $\widehat\Phi$ cannot move, and $\widehat N(\tau)=-\widehat
C_{\Exp}(\tau)<0$ for every expert and every $\tau$ we tried
($\tau\in\{1000,2000,5000,10000\}$). The dry run, run at the full horizon
$T'=50{,}000$ over ten seeds with all $13{,}330$ actions scored, does reproduce
the observed ordering at $\tau=1600$ (small $1.021\pm0.029$, base
$1.057\pm0.087$, large $1.084\pm0.082$; large over small is significant,
paired $t=2.91$) and its cold-start accuracy matches the online runs
($8.46\%$ against $8.39\%$), but it underestimates the magnitude. The dry
run's multiplier is a ratio of regrets, so it is to be compared with the
observed regret ratios: its excesses over one, $2.2\%$, $5.7\%$ and $8.4\%$,
stand against observed excesses of $7.0\%$, $18.7\%$ and $26.6\%$, a shortfall
by a factor of about $3.2$ that is nearly the same for all three experts. It
ranks the experts correctly and misses their magnitude by a common factor.

Both failures have the same cause. The estimators are built from $\widehat
M_{\Exp}$, that is from coverage, and by
Corollary~\ref{cor:coverage-complete} coverage cannot distinguish experts on
this data---the price ratios, measured post hoc, are $1.73$, $1.59$ and
$1.58$, ordered opposite to the multipliers. What distinguishes them is the transfer, and by
Proposition~\ref{prop:transfer} the transfer is the ridge projection of the
residuals \emph{at the actions the expert chooses}, which is not a function of
$\widehat M_{\Exp}$ and is not available before the expert is run. Predicting
the multiplier of a given expert before running it therefore remains open, and
Proposition~\ref{prop:transfer-regret} bounds what the second channel can
deliver without identifying which expert will realize it.

\subsection{Consistency}
\label{app:consistency}

\begin{lemma}[Concentration of the held-out quantities]
\label{lem:diag-concentration}
Suppose contexts are i.i.d.\ and $\widehat\Delta_{s,a}=\Delta_{s,a}$ (labels are
exact). Then for any fixed $V$, $\widehat\Phi(V)$ is an average of $n$ i.i.d.\
variables in $[0,\Delta_{\max}]$ with mean $\E[\Phi_t(V)]$, so
$|\widehat\Phi(V)-\E\Phi_t(V)|\le\Delta_{\max}\sqrt{\log(2/\delta)/(2n)}$ with
probability $1-\delta$ (Hoeffding); the same holds for
$\widehat C_{\Exp}(\tau)/\tau$ and for the two accuracies. For the warm design,
$\|\tfrac1n\widehat M_{\Exp}-\E[xx^\top]\|_{\mathrm{op}}\le CL^2\sqrt{\log(2d/\delta)/n}$
with probability $1-\delta$ (matrix Bernstein), and since
$V\mapsto\|x\|_{V^{-1}}$ is Lipschitz on $\{V\succeq\lambda I\}$ with constant
$\|x\|_2/(2\lambda^{3/2})$ in operator norm, $\widehat\Phi(\lambda I+(\tau/n)\widehat M_{\Exp})$
is within $O\bigl(\beta L\tau\lambda^{-3/2}L^2\sqrt{\log(2d/\delta)/n}\bigr)$ of
its population counterpart.
\end{lemma}

\begin{proof}
The scalar statements are Hoeffding's inequality applied to bounded i.i.d.\
summands; $\Phi_t(V)\in[0,\Delta_{\max}]$ because it is a maximum of minima
each bounded by a gap. The matrix statement is the matrix Bernstein inequality
for i.i.d.\ PSD summands of operator norm at most $L^2$. For the Lipschitz
claim, $\|x\|^2_{V^{-1}}-\|x\|^2_{V'^{-1}}=x^\top V^{-1}(V'-V)V'^{-1}x$, bounded by
$\|x\|_2^2\|V'-V\|_{\mathrm{op}}/\lambda^2$, and $|\sqrt u-\sqrt v|\le|u-v|/(2\sqrt{\min\{u,v\}})$
with $\min\{u,v\}\ge\|x\|_2^2/(\lambda+L^2\tau)$; the stated constant follows
after simplification. Plugging into the definition of $\widehat\Phi$ and
using that $\min\{\cdot,\cdot\}$ and $\max$ are $1$-Lipschitz completes the
proof.
\end{proof}

When labels are not exact, $\widehat\Delta$ carries the error of the ridge
fit; this affects the first argument of the minimum only and is the usual
price of a gap-dependent quantity.

\subsection{Relation to the Four Proxies}
\label{app:proxies}

The four statistics of Table~\ref{tab:diagnostics} are cheaper than
$\widehat\Phi$ and each bounds one ingredient of it:
\begin{itemize}[itemsep=1pt,topsep=1pt,leftmargin=*]
\item \emph{Effective rank} $\widehat r=\mathrm{tr}(\widehat\Sigma)/\|\widehat\Sigma\|_{\mathrm{op}}$
of the context covariance. By the elliptical potential bound, cold LinUCB can
have at most $O(\widehat r\log T)$ rounds at which some played width exceeds
a fixed level, so a small $\widehat r$ means the cold-start widths, and with
them $\Phi_t(V^{\Cold}_t)$, fall quickly on their own.
\item \emph{Coherence} $\widehat{\bar\alpha}$, the median over actions of the
minimum pairwise cosine. By Lemma~\ref{lem:transfer} the cosine is the
transfer coefficient at a cold start up to the factor $\|x\|^2/(\lambda+\|x\|^2)$;
positive coherence means cold observations shrink the widths of other
arms from the first round.
\item \emph{Gap tail}, the fraction of held-out contexts whose estimated
runner-up gap is below $\eta$. Small gaps cap the first argument of the
minimum in $\Phi_t$; a heavy tail means many rounds are cheap whatever the
widths are, which by itself argues \emph{against} a large multiplier.
\item \emph{Subspace mismatch}, the largest principal angle between the
top-$\widehat r$ eigenspaces of $\widehat M_{\Exp}$ and of the context
covariance. An angle near $90^\circ$ means the expert's data covers
directions that carry little of the online context mass, so by
Lemma~\ref{lem:loewner} it cannot reduce the widths that matter. The
statistic is a coarse indicator that $\widehat\Phi$ supersedes; it is not a
deployment rule.
\end{itemize}

\section*{NeurIPS Paper Checklist}




\begin{enumerate}

\item {\bf Claims}
    \item[] Question: Do the main claims made in the abstract and introduction accurately reflect the paper's contributions and scope?
    \item[] Answer: \answerYes{} % Replace by \answerYes{}, \answerNo{}, or \answerNA{}.
    \item[] Justification: The abstract and introduction state the problem setting (LLM warm-starts in linear contextual bandits), the exact decomposition of the hybrid's advantage (Proposition~\ref{prop:decomp}), the price bound on the online multiplier (Theorem~\ref{thm:general-price}), the completeness of that bound under linear rewards (Proposition~\ref{prop:response-sum} and Corollary~\ref{cor:coverage-complete}), the transfer term that misspecification introduces (Propositions~\ref{prop:transfer} and~\ref{prop:transfer-regret}), and the experiments that measure the decomposition and isolate the transfer. Claims are stated with the scope the proofs support; in particular we state explicitly that Proposition~\ref{prop:transfer-regret} bounds what the second channel can deliver without predicting which expert will realize it, and Appendix~\ref{app:diagnostics} reports that the dry run ranks the experts correctly but no pre-runtime estimator predicts the size of the multiplier.
    \item[] Guidelines:
    \begin{itemize}
        \item The answer \answerNA{} means that the abstract and introduction do not include the claims made in the paper.
        \item The abstract and/or introduction should clearly state the claims made, including the contributions made in the paper and important assumptions and limitations. A \answerNo{} or \answerNA{} answer to this question will not be perceived well by the reviewers. 
        \item The claims made should match theoretical and experimental results, and reflect how much the results can be expected to generalize to other settings. 
        \item It is fine to include aspirational goals as motivation as long as it is clear that these goals are not attained by the paper. 
    \end{itemize}

\item {\bf Limitations}
    \item[] Question: Does the paper discuss the limitations of the work performed by the authors?
    \item[] Answer: \answerYes{} % Replace by \answerYes{}, \answerNo{}, or \answerNA{}.
    \item[] Justification: Limitations are discussed at the end of Section~\ref{sec:conclusion}: linear reward, and the fact that on our main dataset the channel the linear theory leaves out (the transfer) is the larger one and has no pre-runtime predictor; the post-hoc nature of the price ratio on real data; the constructed nature of the instance of Theorem~\ref{thm:price} and the unoptimized constants of the catalog, the one-sided nature of the pre-runtime estimates and their provisional thresholds, the dependence on a labeled held-out sample and on the expert's proposals, the restriction to LinUCB-style optimistic selection in Lemma~\ref{lem:width-limited}, the use of only two real datasets, and the lack of multi-expert stress tests.
    \item[] Guidelines:
    \begin{itemize}
        \item The answer \answerNA{} means that the paper has no limitation while the answer \answerNo{} means that the paper has limitations, but those are not discussed in the paper. 
        \item The authors are encouraged to create a separate ``Limitations'' section in their paper.
        \item The paper should point out any strong assumptions and how robust the results are to violations of these assumptions (e.g., independence assumptions, noiseless settings, model well-specification, asymptotic approximations only holding locally). The authors should reflect on how these assumptions might be violated in practice and what the implications would be.
        \item The authors should reflect on the scope of the claims made, e.g., if the approach was only tested on a few datasets or with a few runs. In general, empirical results often depend on implicit assumptions, which should be articulated.
        \item The authors should reflect on the factors that influence the performance of the approach. For example, a facial recognition algorithm may perform poorly when image resolution is low or images are taken in low lighting. Or a speech-to-text system might not be used reliably to provide closed captions for online lectures because it fails to handle technical jargon.
        \item The authors should discuss the computational efficiency of the proposed algorithms and how they scale with dataset size.
        \item If applicable, the authors should discuss possible limitations of their approach to address problems of privacy and fairness.
        \item While the authors might fear that complete honesty about limitations might be used by reviewers as grounds for rejection, a worse outcome might be that reviewers discover limitations that aren't acknowledged in the paper. The authors should use their best judgment and recognize that individual actions in favor of transparency play an important role in developing norms that preserve the integrity of the community. Reviewers will be specifically instructed to not penalize honesty concerning limitations.
    \end{itemize}

\item {\bf Theory assumptions and proofs}
    \item[] Question: For each theoretical result, does the paper provide the full set of assumptions and a complete (and correct) proof?
    \item[] Answer: \answerYes{} % Replace by \answerYes{}, \answerNo{}, or \answerNA{}.
    \item[] Justification: Section~\ref{sec:preliminaries} states the linear contextual-bandit setup and the three policies compared; Section~\ref{sec:catalog} states the decomposition (Proposition~\ref{prop:decomp}) and the catalog. Section~\ref{sec:width} states the confidence event under which Lemma~\ref{lem:width-limited} holds, and Lemmas~\ref{lem:loewner} and~\ref{lem:transfer} are stated as exact linear-algebra identities. Section~\ref{sec:handoff} states Proposition~\ref{thm:handoff} with the exact determinant ratio. All proofs are in Appendix~\ref{app:proofs}; the consistency statement for the held-out estimates is Lemma~\ref{lem:diag-concentration} in Appendix~\ref{app:diagnostics}.
    \item[] Guidelines:
    \begin{itemize}
        \item The answer \answerNA{} means that the paper does not include theoretical results. 
        \item All the theorems, formulas, and proofs in the paper should be numbered and cross-referenced.
        \item All assumptions should be clearly stated or referenced in the statement of any theorems.
        \item The proofs can either appear in the main paper or the supplemental material, but if they appear in the supplemental material, the authors are encouraged to provide a short proof sketch to provide intuition. 
        \item Inversely, any informal proof provided in the core of the paper should be complemented by formal proofs provided in appendix or supplemental material.
        \item Theorems and Lemmas that the proof relies upon should be properly referenced. 
    \end{itemize}

    \item {\bf Experimental result reproducibility}
    \item[] Question: Does the paper fully disclose all the information needed to reproduce the main experimental results of the paper to the extent that it affects the main claims and/or conclusions of the paper (regardless of whether the code and data are provided or not)?
    \item[] Answer: \answerYes{} % Replace by \answerYes{}, \answerNo{}, or \answerNA{}.
    \item[] Justification: Section~\ref{sec:exp-setup} specifies the datasets (OneShotWikiLinks-311 and AmazonCat-13K), baselines (Router/Algorithm~\ref{alg:confidence-gated}, LLM-only, Cold LinUCB), embedding pipelines (BERT-768d for Flan-T5, BGE-1024d for Qwen2.5), hyperparameters ($\lambda=0.1$, $\alpha_t=\max(0.1, 0.9985^t)$), the confidence-parameter grids and the protocol that selects $c$ on seeds $11$--$20$ and reports seeds $1$--$10$, and the round-level update protocol for $(V_t,b_t)$. Sections~\ref{sec:wiki} and~\ref{sec:amazon} report results for every configuration, including the WikiLinks confidence-sweep table (Table~\ref{tab:wiki-csweep}) and AmazonCat learning-dynamics figure (Figure~\ref{fig:amazon-dynamics}). The diagnostics computed pre-runtime are formally defined in Appendix~\ref{app:diagnostics}. All datasets and pretrained models used are publicly available.
    \item[] Guidelines:
    \begin{itemize}
        \item The answer \answerNA{} means that the paper does not include experiments.
        \item If the paper includes experiments, a \answerNo{} answer to this question will not be perceived well by the reviewers: Making the paper reproducible is important, regardless of whether the code and data are provided or not.
        \item If the contribution is a dataset and\slash or model, the authors should describe the steps taken to make their results reproducible or verifiable. 
        \item Depending on the contribution, reproducibility can be accomplished in various ways. For example, if the contribution is a novel architecture, describing the architecture fully might suffice, or if the contribution is a specific model and empirical evaluation, it may be necessary to either make it possible for others to replicate the model with the same dataset, or provide access to the model. In general. releasing code and data is often one good way to accomplish this, but reproducibility can also be provided via detailed instructions for how to replicate the results, access to a hosted model (e.g., in the case of a large language model), releasing of a model checkpoint, or other means that are appropriate to the research performed.
        \item While NeurIPS does not require releasing code, the conference does require all submissions to provide some reasonable avenue for reproducibility, which may depend on the nature of the contribution. For example
        \begin{enumerate}
            \item If the contribution is primarily a new algorithm, the paper should make it clear how to reproduce that algorithm.
            \item If the contribution is primarily a new model architecture, the paper should describe the architecture clearly and fully.
            \item If the contribution is a new model (e.g., a large language model), then there should either be a way to access this model for reproducing the results or a way to reproduce the model (e.g., with an open-source dataset or instructions for how to construct the dataset).
            \item We recognize that reproducibility may be tricky in some cases, in which case authors are welcome to describe the particular way they provide for reproducibility. In the case of closed-source models, it may be that access to the model is limited in some way (e.g., to registered users), but it should be possible for other researchers to have some path to reproducing or verifying the results.
        \end{enumerate}
    \end{itemize}


\item {\bf Open access to data and code}
    \item[] Question: Does the paper provide open access to the data and code, with sufficient instructions to faithfully reproduce the main experimental results, as described in supplemental material?
    \item[] Answer: \answerNo{} % Replace by \answerYes{}, \answerNo{}, or \answerNA{}.
    \item[] Justification: To preserve double-blind review, this submission does not include a code-release URL. The datasets used (OneShotWikiLinks-311 and AmazonCat-13K) are public benchmarks cited in Section~\ref{sec:exp-setup}, and the pretrained models (Flan-T5, Qwen2.5, BERT, BGE) are publicly available with citations in the body. All experimental settings, hyperparameters, prompts/embedding pipelines, and seed protocols are reported in Sections~\ref{sec:exp-setup}--\ref{sec:amazon} in sufficient detail to reproduce the results; code will be released after the review period.
    \item[] Guidelines:
    \begin{itemize}
        \item The answer \answerNA{} means that paper does not include experiments requiring code.
        \item Please see the NeurIPS code and data submission guidelines (\url{https://neurips.cc/public/guides/CodeSubmissionPolicy}) for more details.
        \item While we encourage the release of code and data, we understand that this might not be possible, so \answerNo{} is an acceptable answer. Papers cannot be rejected simply for not including code, unless this is central to the contribution (e.g., for a new open-source benchmark).
        \item The instructions should contain the exact command and environment needed to run to reproduce the results. See the NeurIPS code and data submission guidelines (\url{https://neurips.cc/public/guides/CodeSubmissionPolicy}) for more details.
        \item The authors should provide instructions on data access and preparation, including how to access the raw data, preprocessed data, intermediate data, and generated data, etc.
        \item The authors should provide scripts to reproduce all experimental results for the new proposed method and baselines. If only a subset of experiments are reproducible, they should state which ones are omitted from the script and why.
        \item At submission time, to preserve anonymity, the authors should release anonymized versions (if applicable).
        \item Providing as much information as possible in supplemental material (appended to the paper) is recommended, but including URLs to data and code is permitted.
    \end{itemize}


\item {\bf Experimental setting/details}
    \item[] Question: Does the paper specify all the training and test details (e.g., data splits, hyperparameters, how they were chosen, type of optimizer) necessary to understand the results?
    \item[] Answer: \answerYes{} % Replace by \answerYes{}, \answerNo{}, or \answerNA{}.
    \item[] Justification: Section~\ref{sec:exp-setup} reports all training/test details: regularization $\lambda=0.1$, confidence schedule $\alpha_t=\max(0.1, 0.9985^t)$, the confidence-parameter grids with the disjoint tuning (seeds $11$--$20$) and evaluation (seeds $1$--$10$) blocks, embedding pipelines for both feature spaces, and the round-level update rule for the LinUCB sufficient statistics. The held-out sample size for diagnostic computation (5{,}000 contexts) and the Router's routing rule are also specified. Per-dataset configurations are detailed in Sections~\ref{sec:wiki} and~\ref{sec:amazon}, including the WikiLinks $c$-sweep results (Table~\ref{tab:wiki-csweep}). The provisional threshold of the net-value rule is stated in Appendix~\ref{app:diagnostics}.
    \item[] Guidelines:
    \begin{itemize}
        \item The answer \answerNA{} means that the paper does not include experiments.
        \item The experimental setting should be presented in the core of the paper to a level of detail that is necessary to appreciate the results and make sense of them.
        \item The full details can be provided either with the code, in appendix, or as supplemental material.
    \end{itemize}

\item {\bf Experiment statistical significance}
    \item[] Question: Does the paper report error bars suitably and correctly defined or other appropriate information about the statistical significance of the experiments?
    \item[] Answer: \answerYes{} % Replace by \answerYes{}, \answerNo{}, or \answerNA{}.
    \item[] Justification: All reported results (Tables~\ref{tab:wiki-results}, \ref{tab:amazon-results}, \ref{tab:amazon-decomp}, \ref{tab:two-channel}, \ref{tab:baselines}) are mean $\pm$ standard deviation over seeds $1$--$10$, and every confidence parameter is chosen on the disjoint block of seeds $11$--$20$ (Tables~\ref{tab:bge-csweep} and~\ref{tab:wiki-csweep}), so no reported number comes from a seed used for tuning. Multipliers are reported both as accuracy ratios and as regret ratios. Comparisons between experts in Table~\ref{tab:two-channel} are paired $t$-tests over the ten shared seeds, with the threshold $t^\star=2.26$ stated. Seed-to-seed variability is far larger on AmazonCat than on WikiLinks: cold LinUCB alone ranges over $5.1$--$11.1\%$ across seeds, since with $13{,}330$ actions and $50{,}000$ rounds the order in which arms are first explored matters a great deal; we therefore report the ratio of seed-averaged accuracies rather than the average of per-seed ratios, and state both where they differ. The confidence sweeps span four orders of magnitude in $c$, so the negative results are not an artifact of the gate's hyperparameter. The controlled answer-injection tests of Appendix~\ref{app:dominant} use $20$ paired seeds per condition and report paired $t$-tests on accuracy and exact McNemar tests on the settled category.
    \item[] Guidelines:
    \begin{itemize}
        \item The answer \answerNA{} means that the paper does not include experiments.
        \item The authors should answer \answerYes{} if the results are accompanied by error bars, confidence intervals, or statistical significance tests, at least for the experiments that support the main claims of the paper.
        \item The factors of variability that the error bars are capturing should be clearly stated (for example, train/test split, initialization, random drawing of some parameter, or overall run with given experimental conditions).
        \item The method for calculating the error bars should be explained (closed form formula, call to a library function, bootstrap, etc.)
        \item The assumptions made should be given (e.g., Normally distributed errors).
        \item It should be clear whether the error bar is the standard deviation or the standard error of the mean.
        \item It is OK to report 1-sigma error bars, but one should state it. The authors should preferably report a 2-sigma error bar than state that they have a 96\% CI, if the hypothesis of Normality of errors is not verified.
        \item For asymmetric distributions, the authors should be careful not to show in tables or figures symmetric error bars that would yield results that are out of range (e.g., negative error rates).
        \item If error bars are reported in tables or plots, the authors should explain in the text how they were calculated and reference the corresponding figures or tables in the text.
    \end{itemize}

\item {\bf Experiments compute resources}
    \item[] Question: For each experiment, does the paper provide sufficient information on the computer resources (type of compute workers, memory, time of execution) needed to reproduce the experiments?
    \item[] Answer: \answerYes{} % Replace by \answerYes{}, \answerNo{}, or \answerNA{}.
    \item[] Justification: Each LLM (Flan-T5 at 60M--780M parameters \cite{chung2024scaling}, Qwen2.5 at 0.5B--7B \cite{hui2024qwen2}) is run once over each dataset with greedy decoding and its proposals are cached, so no LLM is called during the online experiments; their cost is dominated by scoring every action with LinUCB at every round. All online runs used one machine with a single NVIDIA RTX~4000 Ada GPU (20\,GB), a 24-core CPU and 46\,GB of RAM, with 10--14 runs sharing the GPU. At that sharing level an AmazonCat-13K run ($T=50{,}000$, $13{,}330$ actions) takes about one hour and a OneShotWikiLinks-311 run ($T=622{,}000$, $311$ actions) about seven to eight hours; memory, not GPU throughput, limits how many WikiLinks runs fit at once. The runner also has a CPU-only path, verified to give results identical to the GPU path. The full study, including the confidence sweeps and ablations, is on the order of $10^3$ runs; runs from earlier versions of this work that are not reported were of the same kind.
    \item[] Guidelines:
    \begin{itemize}
        \item The answer \answerNA{} means that the paper does not include experiments.
        \item The paper should indicate the type of compute workers CPU or GPU, internal cluster, or cloud provider, including relevant memory and storage.
        \item The paper should provide the amount of compute required for each of the individual experimental runs as well as estimate the total compute. 
        \item The paper should disclose whether the full research project required more compute than the experiments reported in the paper (e.g., preliminary or failed experiments that didn't make it into the paper). 
    \end{itemize}
    
\item {\bf Code of ethics}
    \item[] Question: Does the research conducted in the paper conform, in every respect, with the NeurIPS Code of Ethics \url{https://neurips.cc/public/EthicsGuidelines}?
    \item[] Answer: \answerYes{} % Replace by \answerYes{}, \answerNo{}, or \answerNA{}.
    \item[] Justification: To the best of our knowledge, the research conforms with the NeurIPS Code of Ethics. The paper develops theory and algorithms for contextual bandits using public benchmark datasets and public pretrained models; it does not involve human subjects, private personal data collection, sensitive demographic data, or real-world deployment.
    \item[] Guidelines:
    \begin{itemize}
        \item The answer \answerNA{} means that the authors have not reviewed the NeurIPS Code of Ethics.
        \item If the authors answer \answerNo, they should explain the special circumstances that require a deviation from the Code of Ethics.
        \item The authors should make sure to preserve anonymity (e.g., if there is a special consideration due to laws or regulations in their jurisdiction).
    \end{itemize}


\item {\bf Broader impacts}
    \item[] Question: Does the paper discuss both potential positive societal impacts and negative societal impacts of the work performed?
    \item[] Answer: \answerYes{} % Replace by \answerYes{}, \answerNo{}, or \answerNA{}.
    \item[] Justification: A brief broader-impact paragraph appears at the end of Section~\ref{sec:conclusion}. Positive impacts include reducing unnecessary LLM inference in sequential decision systems by identifying when cold-start LinUCB suffices, which lowers compute cost and energy use. Potential negative impacts include the risk that LLM-guided contextual-bandit systems may inherit biases from pretrained models or logged data if deployed in recommendation, ranking, or other consequential decision-making pipelines without auditing. The work is primarily theoretical and diagnostic, so concrete deployment harms depend on follow-on applications not studied here.
    \item[] Guidelines:
    \begin{itemize}
        \item The answer \answerNA{} means that there is no societal impact of the work performed.
        \item If the authors answer \answerNA{} or \answerNo, they should explain why their work has no societal impact or why the paper does not address societal impact.
        \item Examples of negative societal impacts include potential malicious or unintended uses (e.g., disinformation, generating fake profiles, surveillance), fairness considerations (e.g., deployment of technologies that could make decisions that unfairly impact specific groups), privacy considerations, and security considerations.
        \item The conference expects that many papers will be foundational research and not tied to particular applications, let alone deployments. However, if there is a direct path to any negative applications, the authors should point it out. For example, it is legitimate to point out that an improvement in the quality of generative models could be used to generate Deepfakes for disinformation. On the other hand, it is not needed to point out that a generic algorithm for optimizing neural networks could enable people to train models that generate Deepfakes faster.
        \item The authors should consider possible harms that could arise when the technology is being used as intended and functioning correctly, harms that could arise when the technology is being used as intended but gives incorrect results, and harms following from (intentional or unintentional) misuse of the technology.
        \item If there are negative societal impacts, the authors could also discuss possible mitigation strategies (e.g., gated release of models, providing defenses in addition to attacks, mechanisms for monitoring misuse, mechanisms to monitor how a system learns from feedback over time, improving the efficiency and accessibility of ML).
    \end{itemize}
    
\item {\bf Safeguards}
    \item[] Question: Does the paper describe safeguards that have been put in place for responsible release of data or models that have a high risk for misuse (e.g., pre-trained language models, image generators, or scraped datasets)?
    \item[] Answer: \answerNA{} % Replace by \answerYes{}, \answerNo{}, or \answerNA{}.
    \item[] Justification: The paper does not release a new pretrained model, image generator, scraped dataset, or other artifact that poses a high risk of misuse. The experiments use existing public datasets and existing pretrained language/embedding models, and the contribution is a theoretical and algorithmic framework with diagnostic experiments.
    \item[] Guidelines:
    \begin{itemize}
        \item The answer \answerNA{} means that the paper poses no such risks.
        \item Released models that have a high risk for misuse or dual-use should be released with necessary safeguards to allow for controlled use of the model, for example by requiring that users adhere to usage guidelines or restrictions to access the model or implementing safety filters. 
        \item Datasets that have been scraped from the Internet could pose safety risks. The authors should describe how they avoided releasing unsafe images.
        \item We recognize that providing effective safeguards is challenging, and many papers do not require this, but we encourage authors to take this into account and make a best faith effort.
    \end{itemize}

\item {\bf Licenses for existing assets}
    \item[] Question: Are the creators or original owners of assets (e.g., code, data, models), used in the paper, properly credited and are the license and terms of use explicitly mentioned and properly respected?
    \item[] Answer: \answerYes{} % Replace by \answerYes{}, \answerNo{}, or \answerNA{}.
    \item[] Justification: All datasets and pretrained models used are publicly released and credited. Datasets: OneShotWikiLinks-311 (a 311-class subset of the OneShotWikiLinks named-entity recognition benchmark, \url{https://www.kaggle.com/generall/oneshotwikilinks}) and AmazonCat-13K (the extreme-classification benchmark \cite{bhatia2016extreme}, \url{http://manikvarma.org/downloads/XC/XMLRepository.html}), described in Section~\ref{sec:exp-setup}. Pretrained models: Flan-T5 \cite{chung2024scaling}, Qwen2.5-Instruct \cite{hui2024qwen2}, Sentence-BERT-style embeddings \cite{reimers2019sentence}, and BAAI/bge-large-en-v1.5 \cite{xiao2024c}, all cited in Sections~\ref{sec:wiki} and~\ref{sec:amazon}. The paper uses these assets for research purposes consistent with their respective licenses.
    \item[] Guidelines:
    \begin{itemize}
        \item The answer \answerNA{} means that the paper does not use existing assets.
        \item The authors should cite the original paper that produced the code package or dataset.
        \item The authors should state which version of the asset is used and, if possible, include a URL.
        \item The name of the license (e.g., CC-BY 4.0) should be included for each asset.
        \item For scraped data from a particular source (e.g., website), the copyright and terms of service of that source should be provided.
        \item If assets are released, the license, copyright information, and terms of use in the package should be provided. For popular datasets, \url{paperswithcode.com/datasets} has curated licenses for some datasets. Their licensing guide can help determine the license of a dataset.
        \item For existing datasets that are re-packaged, both the original license and the license of the derived asset (if it has changed) should be provided.
        \item If this information is not available online, the authors are encouraged to reach out to the asset's creators.
    \end{itemize}

\item {\bf New assets}
    \item[] Question: Are new assets introduced in the paper well documented and is the documentation provided alongside the assets?
    \item[] Answer: \answerNA{} % Replace by \answerYes{}, \answerNo{}, or \answerNA{}.
    \item[] Justification: The paper does not introduce a new dataset or pretrained model. If the authors release new code for the confidence-gated router, diagnostic computation, and experiments, it should be documented separately with installation instructions, reproduction commands, and an explicit license.
    \item[] Guidelines:
    \begin{itemize}
        \item The answer \answerNA{} means that the paper does not release new assets.
        \item Researchers should communicate the details of the dataset\slash code\slash model as part of their submissions via structured templates. This includes details about training, license, limitations, etc. 
        \item The paper should discuss whether and how consent was obtained from people whose asset is used.
        \item At submission time, remember to anonymize your assets (if applicable). You can either create an anonymized URL or include an anonymized zip file.
    \end{itemize}

\item {\bf Crowdsourcing and research with human subjects}
    \item[] Question: For crowdsourcing experiments and research with human subjects, does the paper include the full text of instructions given to participants and screenshots, if applicable, as well as details about compensation (if any)? 
    \item[] Answer: \answerNA{} % Replace by \answerYes{}, \answerNo{}, or \answerNA{}.
    \item[] Justification: The paper does not involve crowdsourcing, user studies, or research with human subjects. All evaluations are offline contextual-bandit experiments using existing benchmark datasets.
    \item[] Guidelines:
    \begin{itemize}
        \item The answer \answerNA{} means that the paper does not involve crowdsourcing nor research with human subjects.
        \item Including this information in the supplemental material is fine, but if the main contribution of the paper involves human subjects, then as much detail as possible should be included in the main paper. 
        \item According to the NeurIPS Code of Ethics, workers involved in data collection, curation, or other labor should be paid at least the minimum wage in the country of the data collector. 
    \end{itemize}

\item {\bf Institutional review board (IRB) approvals or equivalent for research with human subjects}
    \item[] Question: Does the paper describe potential risks incurred by study participants, whether such risks were disclosed to the subjects, and whether Institutional Review Board (IRB) approvals (or an equivalent approval/review based on the requirements of your country or institution) were obtained?
    \item[] Answer: \answerNA{} % Replace by \answerYes{}, \answerNo{}, or \answerNA{}.
    \item[] Justification: The paper does not involve crowdsourcing, user studies, or research with human subjects, so IRB approval or equivalent review is not applicable.
    \item[] Guidelines:
    \begin{itemize}
        \item The answer \answerNA{} means that the paper does not involve crowdsourcing nor research with human subjects.
        \item Depending on the country in which research is conducted, IRB approval (or equivalent) may be required for any human subjects research. If you obtained IRB approval, you should clearly state this in the paper. 
        \item We recognize that the procedures for this may vary significantly between institutions and locations, and we expect authors to adhere to the NeurIPS Code of Ethics and the guidelines for their institution. 
        \item For initial submissions, do not include any information that would break anonymity (if applicable), such as the institution conducting the review.
    \end{itemize}

\item {\bf Declaration of LLM usage}
    \item[] Question: Does the paper describe the usage of LLMs if it is an important, original, or non-standard component of the core methods in this research? Note that if the LLM is used only for writing, editing, or formatting purposes and does \emph{not} impact the core methodology, scientific rigor, or originality of the research, declaration is not required.
    %this research? 
    \item[] Answer: \answerYes{} % Replace by \answerYes{}, \answerNo{}, or \answerNA{}.
    \item[] Justification: LLMs are central to the research question and experimental design. The paper studies frozen LLM-induced policies as warm-starting and routing components in contextual bandits, and the experimental sections specify the LLM model families used as frozen policies. Any LLM use for writing or editing, if applicable, is separate from the core scientific use of LLMs as decision-making components.
    \item[] Guidelines:
    \begin{itemize}
        \item The answer \answerNA{} means that the core method development in this research does not involve LLMs as any important, original, or non-standard components.
        \item Please refer to our LLM policy in the NeurIPS handbook for what should or should not be described.
    \end{itemize}

\end{enumerate}

\end{document}